\documentclass[12pt,a4paper]{article}
\usepackage[utf8]{inputenc}
\usepackage[T2A]{fontenc}
\usepackage[russian]{babel}
\usepackage{amsmath,amssymb}
\usepackage{graphicx}
\usepackage{array}
\usepackage{microtype}
\usepackage{geometry}
\usepackage{hyperref}
\usepackage{xurl}
\geometry{margin=2.3cm}
\setlength{\emergencystretch}{2em}

\title{Трактат об устройстве мира\\
Том I: постановка исследовательской программы и карта структурных режимов}
\author{}
\date{}

\begin{document}

\maketitle

\begin{abstract}
Настоящий текст представляет первый том Трактата об устройстве мира. Его задача --- не объявить завершённую физическую теорию, а зафиксировать исходную исследовательскую программу, ввести рабочий язык описания и собрать в одной схеме геометрические, спектральные, топологические, фазовые и корреляционные мотивы, которые в дальнейшем предполагается связать с единым структурным источником. В этом томе формулируются исходные онтологические тезисы, вводятся инварианты и статусные различения, предлагается полуформальный вариационный каркас и строятся первые иллюстративные карты режимов. Тем самым том задаёт рамку дальнейшей работы: от философской гипотезы --- к уточняемому формализму и к будущим проверяемым моделям.
\end{abstract}

\section{Введение}

Этот трактат начинается не с готового доказательства и не с законченной системы аксиом. Он начинается с ощущения формы. Такое ощущение возникает редко: не тогда, когда найдено одно красивое совпадение, а тогда, когда слишком многие независимые линии начинают указывать в одну и ту же сторону. После многих циклов изучения различных теорий, математических выкладок и концепций, которые мы последовательно рассматривали, постепенно стало проступать нечто общее. Не один частный результат, не одна удачная формула и не одна изолированная интерпретация, а единая интуиция о том, что в основании мира лежит более глубокая организующая сущность.

Эта интуиция не была дана заранее. Она родилась не из желания навязать материалу единство, а, напротив, из сопротивления материала. Снова и снова приходилось проходить через разные языки описания: через геометрию и спектр, через топологию и фазу, через локальные коэффициенты и глобальные инварианты, через корреляционные картины и операторные редукции. Каждый из этих языков по отдельности был недостаточен. Но именно их длительное взаимное наложение стало рождать предчувствие, что за ними стоит не набор независимых механизмов, а один более глубокий порядок.

Если выразить это предчувствие самым кратким образом, оно состоит в следующем: мир в своём основании является не совокупностью отдельных вещей, а структурой согласованности. Частицы, поля, геометрии, заряды, массы и фазы могут оказаться не первичными сущностями, а разными режимами чтения одной и той же глубинной ткани. Тогда геометрия перестаёт быть просто ареной, спектр перестаёт быть только вычислительным инструментом, а топология перестаёт быть декоративным усложнением формализма. Все они становятся различимыми сечениями одного источника.

В такой картине пространство не является пустым контейнером, в котором затем размещаются объекты. Оно само участвует в порождении допустимых состояний. Спектр оказывается не вторичной характеристикой уже данных сущностей, а более фундаментальным описанием того, какие режимы вообще разрешены. Топологические и голономные структуры перестают быть техническими дополнениями и начинают читаться как следы глобальной связности мира. А то, что в обычном языке называется частицей, может быть понято как устойчивый локализованный режим внутри более общей структурной ткани.

Из этого постепенно вырастает и более сильная мысль. Возможно, основание реальности не является ни сугубо геометрическим, ни сугубо квантово-полевым, ни сугубо информационным в обычном смысле. Возможно, его лучше всего понимать как организованную систему допустимых состояний, переходов и инвариантов, то есть как глубинную ткань согласованности, у которой есть несколько равноправных способов проявления. Один язык будет видеть её как геометрию. Другой --- как спектральный каркас. Третий --- как корреляционный источник. Четвёртый --- как набор эффективных физических наблюдаемых. Но объект, порождающий эти чтения, может быть один.

Именно поэтому этот трактат не начинается с утверждения, что какая-либо одна из ранее рассматривавшихся теорий уже является окончательной. Наоборот, его отправная установка состоит в том, что ни одна из них по отдельности, вероятно, не исчерпывает целое. Каждая схватывает свой срез. Каждая удерживает свой тип строгости. Каждая освещает свой слой одной и той же глубины. То, что в одной теории появляется как геометрический опорный каркас, в другой может читаться как спектральный инвариант. То, что в одной выступает как голономная фаза, в другой может быть выражено как глобальный корреляционный след. То, что в одной кажется внутренней секторной структурой, в другой может оказаться эффективной тенью более общего закона организации.

Отсюда возникает рабочая философская гипотеза этого трактата: первична не субстанция, а структура; первичен не объект, а закон его допустимого становления; первична не отдельная форма материи, а единый источник, из которого геометрия, спектр, фаза, заряд, корреляция и взаимодействие проявляются как разные редукции. Если эта гипотеза верна хотя бы частично, то задача фундаментальной физики должна быть поставлена заново. Нужно искать не просто новую частицу, не просто новое поле и не просто ещё одно уравнение, а минимальную сущность, из которой одновременно и естественно следуют все эти режимы чтения.

Настоящий трактат и задуман как первый шаг в этом направлении. Его цель --- не закрыть вопрос окончательно, а зафиксировать форму того, что стало возможно увидеть после длительного исследования. Здесь будет предпринята попытка перейти от рассеянных результатов к единому образу мира; от частных проверок --- к общему источнику; от сравнения теорий --- к поиску той сущности, которая, быть может, и делает возможным само их частичное совпадение.

\section{Исходная установка}

В дальнейшем изложении мы будем исходить из того, что накопленные ранее геометрические, спектральные, топологические и операторные результаты не следует понимать как конкурирующие фрагменты. Гораздо плодотворнее читать их как следы одной ещё не до конца сформулированной структуры. Именно эту структуру и предстоит попытаться выделить, описать и затем подвергнуть новому самостоятельному анализу.

\section{Дальнейшие разделы}

Этот текст пока открыт и будет расширяться. Предварительно в трактат предполагается включить следующие разделы:
\begin{itemize}
    \item онтологические тезисы о минимальной глубинной структуре мира;
    \item сопоставление геометрического, спектрального и корреляционного языков;
    \item формулировку рабочей гипотезы о единой сущности;
    \item новый исследовательский протокол от постановки до критериев проверки;
    \item обсуждение совместимости с существующей физикой и границ применимости гипотезы.
\end{itemize}

\section{Онтологические тезисы}

Чтобы перейти от общей интуиции к более строгой форме разговора, зафиксируем несколько исходных тезисов. Эти тезисы пока не являются окончательными аксиомами. Их правильнее понимать как рабочую онтологию, выведенную из многократного сопоставления уже накопленных результатов.

\subsection{Тезис о первичности структуры}

Первичным является не вещество и не набор уже готовых объектов, а структура допустимых состояний, переходов и ограничений. Мир не сначала состоит из ``вещей'', а затем подчиняется законам. Скорее, сами вещи возникают как устойчивые режимы внутри некоторого более глубокого закона согласованности.

Из этого следует важный сдвиг. Частица в такой картине есть не первичный кирпич мира, а стабильная форма организации. Поле есть не субстанция, разлитая по пространству, а способ локального и нелокального связывания допустимых состояний. Геометрия же выступает не фоном, а выражением тех ограничений, внутри которых подобные режимы вообще могут существовать.

\subsection{Тезис о первичности допустимости над локальностью}

Обычная физическая интуиция часто исходит из локальной картины: что где находится, как движется и как взаимодействует вблизи. Однако накопленный материал подсказывает, что глубже локальности лежит более фундаментальный слой --- слой допустимости. Не всякая локально мыслимая конфигурация допустима глобально. Не всякая локальная степень свободы совместима с топологией, спектром, спин-структурой или голономии.

Поэтому глубинная физика должна описывать прежде всего пространство разрешённых режимов, а уже затем --- их локальные приближения. В таком подходе спектр, граничные условия, топологические классы и глобальные фазы оказываются не побочными техническими деталями, а прямыми носителями фундаментальной информации.

\subsection{Тезис о многослойности реальности}

Реальность, по-видимому, имеет несколько уровней чтения, и ни один из них не исчерпывает её полностью. По крайней мере в рамках уже изученного материала можно различить следующие слои:
\begin{itemize}
    \item \textbf{геометрический слой,} где фиксируются объёмы, кривизны, компактные пространства и их инварианты;
    \item \textbf{спектральный слой,} где мир читается через допустимые моды, операторы и асимптотические коэффициенты;
    \item \textbf{тополого-фазовый слой,} где существенны фундаментальная группа, голономии, спин-структуры и дискретные ветви;
    \item \textbf{корреляционный слой,} где одни и те же структуры проявляются как глобальные отношения, а не только как локальные наблюдаемые величины;
    \item \textbf{эффективный физический слой,} где все предыдущие уровни редуцируются к тем величинам, которые наблюдаются как массы, заряды, константы связи и отклики.
\end{itemize}

Рабочая мысль трактата состоит в том, что эти слои не независимы. Они являются различными редукциями одной и той же подлежащей сущности.

\subsection{Тезис о происхождении частиц и взаимодействий}

Если предыдущие тезисы верны, то частицы и взаимодействия должны пониматься заново. Частица есть устойчивая мода глубинной ткани. Заряд есть ориентированная фазовая ветвь допустимого переноса. Масса есть мера структурной инерции соответствующего режима. Взаимодействие есть перестройка режима согласованности между различными секторами одной и той же основы.

Такой взгляд позволяет по-новому читать отношения между геометрией и физикой. Геометрические инварианты тогда не просто сопровождают физику, а кодируют часть самого механизма возникновения наблюдаемых. Спектральные коэффициенты перестают быть внешними числами и становятся отпечатками внутреннего строения допустимых мод. Глобальные фазы перестают быть математическим украшением и начинают выражать то, как реальность сохраняет согласованность при обходе нестягиваемых циклов и переходах между секторами.

\subsection{Тезис о едином источнике}

Из всех предыдущих рассуждений вытекает центральный для этого трактата тезис: существует единый источник, из которого геометрический, спектральный, топологический, корреляционный и эффективный физический уровни возникают как взаимосвязанные чтения. Этот источник пока ещё не имеет окончательной формализации. Тем не менее уже сейчас можно указать на его предполагаемые свойства.

Он должен:
\begin{itemize}
    \item порождать допустимый спектр состояний;
    \item нести глобальную топологическую и фазовую структуру;
    \item допускать локальную редукцию к эффективным наблюдаемым;
    \item обладать внутренней секторностью, позволяющей отличать опорный каркас, фазовый сектор и сектор внутренней нагрузки;
    \item быть читаемым несколькими различными теоретическими языками без потери общей идентичности.
\end{itemize}

Именно этот источник и будет далее рассматриваться как кандидат на минимальную глубинную сущность мира.

\section{Предварительная схема единой сущности}

Пока ещё рано давать окончательное определение, однако в первом приближении искомую сущность можно описать как \textit{структурно-спектральную ткань согласованности}. Это выражение пока не является термином строгой теории; оно вводится как временное имя для объекта, который должен объяснить совместное происхождение:
\begin{itemize}
    \item геометрических инвариантов;
    \item спектральных коэффициентов и операторных следов;
    \item голономии- и фазово-структур;
    \item внутренних секторных разбиений;
    \item устойчивых эффективных наблюдаемые величины.
\end{itemize}

Если пользоваться более образным языком, эта сущность напоминает не вещество, а ткань резонансной организации. Но в отличие от чистой метафоры здесь важно следующее: речь идёт о ткани, которая должна быть способна к строгому чтению. Именно поэтому трактат не останавливается на философской интуиции, а должен перейти к новому исследовательскому протоколу.

\section{Переход к исследовательской программе}

С этого места философский язык должен постепенно уступать место языку построения. Если предшествующее предчувствие не является иллюзией, то оно должно быть переведено в систему исследовательских шагов. Нужно определить, какие свойства единой сущности являются необходимыми, какие --- производными, какие --- наблюдаемыми, а какие --- зависят от языка описания. Иначе говоря, следующий шаг трактата состоит в том, чтобы превратить онтологическую интуицию в проверяемую программу.

\section{Сопоставление трёх исследовательских языков}

Чтобы перейти от общей онтологии к рабочей гипотезе, необходимо сначала прояснить, какими именно языками уже описывалась одна и та же глубинная реальность. В накопленном материале особенно отчётливо выделяются три способа чтения: геометрический, спектральный и корреляционный. По отдельности каждый из них неполон. Вместе они начинают образовывать более цельную картину.

\subsection{Геометрический язык}

Геометрический язык описывает мир через форму допустимого пространства, его объёмы, кривизны, факторизации, фундаментальные группы и минимальные циклы. Его сила заключается в том, что он сразу даёт жёсткие ограничения. Он показывает, что не всякая конфигурация пространства совместима с наличием спин-структуры, топологической нетривиальности, устойчивых фазовых ветвей и глобальных инвариантов.

На этом языке мир выглядит как иерархия вложенных ограничений. Геометрия здесь не пассивна. Она отбирает, какие классы состояний вообще возможны, какие циклы являются физически существенными, какие фазы допускаются и какие глобальные сектора оказываются доступными. В этом смысле геометрический язык есть язык формы допустимости.

Однако у него есть и граница. Сама по себе геометрия ещё не объясняет, каким образом её структура превращается в набор наблюдаемых величин, коэффициентов и режимов отклика. Она указывает каркас, но не всегда даёт достаточно прозрачный механизм чтения этого каркаса в физические числа.

\subsection{Спектральный язык}

Спектральный язык делает следующий шаг. Он читает геометрию не непосредственно через образ пространства, а через её допустимые моды, операторы, следы и асимптотические коэффициенты. Здесь пространство становится читаемым как спектральная структура. В этом языке особенно важны тепло-ядерные коэффициенты, спектральные действия, нормировки, операторы дираковского типа и устойчивость соответствующих комбинаций при деформациях постановки.

Сила спектрального языка состоит в его дисциплине. Он не просто утверждает, что геометрия важна, а показывает, как именно геометрический каркас отпечатывается в числах. Через него становится видно, что объёмный сектор, кривизновый сектор и некоторые глобальные contributions могут проявляться как устойчивые компоненты одного и того же операторного чтения.

Но и здесь остаётся граница. Спектральный язык превосходно описывает допустимые моды и их инварианты, однако не всегда сам по себе отвечает на вопрос, почему именно эти моды должны быть прочитаны как носители более общего корреляционного содержания. Он силён в форме и расчёте, но ещё не обязательно завершён в интерпретации основания.

\subsection{Корреляционный язык}

Корреляционный язык появляется там, где речь идёт уже не только о спектре как таковом, а о том, как различные сектора мира соотносятся как части одной организованной целостности. Его интересует не только то, какие моды допустимы, но и то, как из них складывается связная картина глобального источника. В этом языке особенно важны не столько локальные коэффициенты, сколько сама идея общего корреляционного объекта, из которого различные функциональные чтения могут быть получены как тени.

Этот язык особенно ценен тем, что он удерживает целое. Он позволяет не застрять в отдельных коэффициентах и не свести реальность к набору вычислительных рецептов. Именно он подсказывает, что разные теоретические формализации могут быть не взаимоисключающими описаниями, а различными проекциями одной и той же глубины.

Его граница, однако, симметрична геометрическому и спектральному языкам: корреляционный подход легче указывает смысл, чем строгую форму. Он естественным образом схватывает общность, но нуждается в опоре на геометрию и спектр, чтобы не раствориться в слишком широкой метафизике.

\subsection{Почему нужен их синтез}

Теперь можно сформулировать главную причину, по которой в трактате необходимо совместное рассмотрение всех трёх языков. Геометрия без спектра рискует остаться слишком молчаливой в отношении конкретных чисел и наблюдаемых. Спектр без корреляционного горизонта рискует остаться слишком локальным или слишком инструментальным. Корреляционный язык без геометрии и спектра рискует потерять дисциплину и превратиться в общую метафору.

Но когда все три уровня читаются совместно, возникает гораздо более сильная картина:
\begin{itemize}
    \item геометрия задаёт пространство допустимого;
    \item спектр даёт строгую форму этого допустимого;
    \item корреляционный язык раскрывает, почему разные спектральные и геометрические проявления принадлежат одному источнику.
\end{itemize}

Именно этот синтез и подводит нас к следующему шагу: к формулировке рабочей гипотезы о том, что существует единая сущность, допускающая одновременно геометрическое, спектральное и корреляционное чтение.

\section{Рабочая гипотеза о единой сущности}

На основании предшествующих рассуждений можно сформулировать первую рабочую гипотезу трактата.

\begin{quote}
Существует единая структурно-спектральная сущность, обладающая геометрическим, спектральным и корреляционным аспектами, из которой наблюдаемые физические режимы возникают как различные редукции и проекции.
\end{quote}

Эта гипотеза пока ещё не фиксирует окончательного математического объекта. Она задаёт только минимальное требование к будущей теории: искомая сущность должна быть такой, чтобы одни и те же её свойства могли быть прочитаны тремя взаимодополняющими способами.

\subsection{Минимальные требования к гипотезе}

Если гипотеза имеет смысл, то искомая сущность должна удовлетворять нескольким условиям.

Во-первых, она должна допускать \textit{геометрическую редукцию}, то есть иметь такой режим чтения, в котором становятся осмысленными объёмы, кривизны, циклы, компактные секторы и топологические классы.

Во-вторых, она должна допускать \textit{спектральную редукцию}, то есть порождать операторы, спектры, асимптотические коэффициенты и устойчивые следы, пригодные для строгого анализа.

В-третьих, она должна допускать \textit{корреляционную редукцию}, в которой различные локальные и глобальные проявления можно рассматривать как части одной более глубокой организованности.

В-четвёртых, она должна объяснять, почему некоторые числа и структуры оказываются устойчивыми при изменении режима чтения. Иначе говоря, она должна содержать инвариантное ядро, которое не исчезает при переходе от одного языка описания к другому.

\subsection{Первое рабочее следствие}

Если такая сущность существует, то фундаментальная физика должна быть переосмыслена как теория согласованных редукций. Тогда задача состоит не только в поиске правильного уравнения, но и в построении правильного перехода между языками. Необходимо будет научиться показывать, как геометрическая структура переходит в спектральную, как спектральная структура выражает корреляционную, и как все они совместно редуцируются к наблюдаемым физическим величинам.

Такое понимание сразу меняет акцент: вместо спора о том, какая теория окончательна, возникает вопрос, какой объект лежит глубже тех теорий, которые уже обнаружили его частные тени.

\section{Новый исследовательский протокол}

Если трактат не должен остаться лишь философским текстом, необходимо немедленно зафиксировать форму дальнейшей работы. Ниже предлагается первый вариант исследовательского протокола, который переводит общую интуицию в проверяемую программу.

\subsection{Цель протокола}

Цель состоит не в том, чтобы произвольно объединить уже имеющиеся теории, а в том, чтобы выявить минимальную сущность, которая:
\begin{itemize}
    \item объясняет устойчивость геометрических, спектральных и фазовых структур;
    \item допускает несколько согласованных режимов чтения;
    \item воспроизводит уже обнаруженные инварианты и не разрушает прежние проверки;
    \item открывает путь к новым нетривиальным следствиям.
\end{itemize}

Иначе говоря, протокол должен вести не к декларативному синтезу, а к новому объекту исследования.

\subsection{Первая фаза: инвариантное ядро}

На первом этапе необходимо выделить то, что не исчезает при смене языка описания. Под инвариантным ядром здесь понимается совокупность структур, которые воспроизводятся в геометрическом, спектральном и корреляционном чтении как одна и та же глубинная закономерность.

Предварительно к этому ядру относятся:
\begin{itemize}
    \item устойчивые геометрические масштабы;
    \item устойчивые спектральные коэффициенты;
    \item глобальные фазовые и голономные секторы;
    \item различение опорный каркас, фазовый сектор и сектор внутренней нагрузки;
    \item воспроизводимость структуры при переходе между несколькими режимами редукции.
\end{itemize}

Задача первой фазы --- определить, какие из этих элементов являются действительно первичными, а какие лишь производными формами записи.

\subsection{Вторая фаза: минимальная формализация сущности}

После выделения инвариантного ядра необходимо предложить минимальную математическую форму искомой сущности. На этом этапе нельзя требовать полной завершённости. Достаточно, чтобы формализация удовлетворяла трём требованиям:
\begin{enumerate}
    \item она допускает геометрическое чтение;
    \item она допускает спектральное чтение;
    \item она допускает корреляционное чтение.
\end{enumerate}

При этом предпочтение должно отдаваться не самой красивой, а самой экономной форме: той, где число исходных структур минимально, но уже достаточно для воспроизведения наблюдаемых следствий.

\subsection{Третья фаза: вывод редукций}

Далее необходимо не просто постулировать наличие нескольких языков, а явно показать переходы между ними. Это означает построение трёх основных отображений:
\begin{itemize}
    \item от глубинной сущности к геометрическому каркасу;
    \item от глубинной сущности к спектральному операторному чтению;
    \item от глубинной сущности к корреляционному описанию глобальной организованности.
\end{itemize}

Только после этого будет возможно строго обсуждать, какие аспекты различных теорий являются подлинными редукциями общего источника, а какие --- лишь удобными приближениями.

\subsection{Четвёртая фаза: проверка на прежнем материале}

Новая гипотеза не имеет права игнорировать уже накопленный корпус проверок. Поэтому следующая фаза протокола состоит в ретроспективной верификации. Нужно проверить, способен ли новый объект естественным образом объяснить уже полученные результаты, в частности:
\begin{itemize}
    \item устойчивость геометрических инвариантов;
    \item устойчивость спектральных коэффициентов при деформациях постановки;
    \item отделение локального сектора от глобального голономии-сектора;
    \item устойчивость картины обогащённого оператора;
    \item появление структуры операторной атрибуции.
\end{itemize}

Если новая сущность не может воспроизвести эти узлы, она должна быть отвергнута или существенно переработана.

\subsection{Пятая фаза: новые предсказательные различия}

Наиболее важной фазой будет та, на которой новая программа начнёт давать различия. Недостаточно повторить уже сделанные результаты. Необходимо получить хотя бы один класс новых следствий, которые не были заранее заложены в исходных построениях. Это могут быть:
\begin{itemize}
    \item новые связи между секторами;
    \item новые критерии устойчивости;
    \item новые ограничения на допустимые компактные пространства;
    \item новые операторных соответствия;
    \item новые численные или структурные тесты.
\end{itemize}

Именно на этом этапе станет ясно, является ли трактат новым началом или лишь новой формой комментария к уже известному.

\section{Критерии подтверждения и критерии отказа}

Всякая новая программа рискует стать неуязвимой риторикой, если заранее не определить, что именно будет считаться её успехом, а что --- её неудачей. Поэтому необходимо сразу обозначить критерии подтверждения и критерии отказа.

\subsection{Что будет считаться подтверждением}

Подтверждением гипотезы о единой сущности будет считаться ситуация, в которой удаётся одновременно показать следующее:
\begin{itemize}
    \item один и тот же минимальный объект действительно допускает три режима чтения;
    \item этот объект естественно воспроизводит уже обнаруженные устойчивые структуры;
    \item переход между языками не требует специально введённых вставок для каждого отдельного результата;
    \item новая схема даёт хотя бы один нетривиальный новый вывод, отсутствовавший в исходных построениях;
    \item совокупная картина становится проще, а не сложнее, чем сумма разрозненных объяснений.
\end{itemize}

Именно последний пункт особенно важен. Настоящая глубинная сущность должна сокращать избыточность описания, а не множить её.

\subsection{Что будет считаться признаком неудачи}

Неудачей программы следует считать любую из следующих ситуаций:
\begin{itemize}
    \item новая сущность не даёт общего инвариантного ядра, а лишь переименовывает старые результаты;
    \item переход между геометрическим, спектральным и корреляционным языками требует независимой подгонки на каждом этапе;
    \item новая формализация не удерживает уже подтверждённые структурные узлы;
    \item теория оказывается слишком общей и допускает любые исходы без возможности различения;
    \item язык трактата остаётся красивым, но перестаёт производить строгие и локально проверяемые следствия.
\end{itemize}

Принципиально важно сохранить право на отказ. Без этого новая программа утратит научный характер и превратится лишь в мировоззренческий проект.

\section{Граница между философией и теорией}

Настоящий текст сознательно начинается с философского уровня, потому что без него невозможно увидеть целое. Но он не имеет права оставаться на этом уровне слишком долго. Философия здесь нужна не как замена теории, а как среда, в которой впервые становится видима форма будущего объекта. После этого она обязана уступить место более строгому построению.

Именно поэтому дальнейшие разделы трактата должны быть всё более дисциплинированными. Сначала необходимо зафиксировать карту соответствий между уже существующими теориями и режимами чтения. Затем --- предложить минимальную математическую форму искомой сущности. Затем --- построить новый анализ, независимый от прежних текстов по форме, но опирающийся на уже накопленные проверки по существу. Только такой путь даст право говорить не о красивом синтезе, а о новом исследовательском шаге.

\section{Карта соответствий между RPFT, UGSM и TOE}

Следующий необходимый шаг состоит в том, чтобы явно указать, как уже накопленные исследовательские корпуса соотносятся друг с другом. Без такой карты трактат рискует остаться слишком абстрактным. С картой же становится видно, что разные теории могут быть поняты не как конкуренты, а как разные режимы доступа к одному и тому же содержанию.

\subsection{RPFT как широкий каркас}

В этой картине RPFT естественно читается как наиболее широкий уровень. Она задаёт общую мета-рамку: выбор фундаментальной геометризации, аксиоматику минимального пространства, топологические требования, роль компактные секторы, логику допустимости и онтологическую установку в целом. Именно на этом уровне впервые появляется мысль, что физика может быть выведена не из набора внешне данных параметров, а из ограничений, налагаемых самой структурой мира.

Поэтому RPFT в трактате следует понимать не как одну частную модель среди прочих, а как язык общей постановки. Она отвечает прежде всего на вопросы вида: почему внутреннее пространство должно быть нетривиальным; почему важны спин-структуры; почему глобальная топология не может быть сведена к локальной геометрии; почему фундаментальные безразмерные числа могут иметь геометрическое происхождение.

Иными словами, RPFT есть язык \textit{архитектурной постановки мира}.

\subsection{UGSM как строгая редукция}

UGSM в этой схеме занимает иное место. Она не дублирует весь широкий каркас RPFT, а выделяет из него строгий спектрально-геометрический сектор. Именно здесь общая постановка подвергается математической дисциплине: появляются операторы, тепло-ядерные коэффициенты, спектральные действия, вычислимые нормировки, строгие и полустрогие derivations. Если RPFT задаёт пространство смыслов, то UGSM задаёт пространство проверяемых структур.

Поэтому UGSM следует понимать как \textit{строгую редукцию RPFT}, а не как отдельный и самодостаточный метафизический универсум. Её сила в том, что она делает широкий каркас читаемым в числах, коэффициентах, инвариантах и операторных конструкциях. Именно UGSM даёт возможность говорить о контролируемой устойчивости и проверяемости.

Иначе говоря, UGSM есть язык \textit{строгого спектрального чтения мира}.

\subsection{TOE как корреляционно-архитектурное чтение}

TOE, в свою очередь, естественно занимает третью позицию. Она не совпадает ни с широкой геометрической аксиоматикой RPFT, ни со строгой дисциплиной UGSM. Её сила в другом: она удерживает целое как архитектуру взаимосвязей и позволяет рассматривать различные сектора как проявления одного корреляционного источника. Там, где UGSM сильна в операторных чтении, TOE сильна в схватывании глубинной связности.

На языке трактата TOE разумнее всего понимать как \textit{корреляционно-архитектурное чтение} той же подлежащей сущности. Она важна потому, что показывает следующее: один и тот же фундаментальный источник может быть выражен не только через оператор и спектр, но и через более глобальный язык отношения, организации и тени.

Иными словами, TOE есть язык \textit{целостного корреляционного чтения мира}.

\subsection{Объединяющая схема}

Теперь можно свести эти различия в одну схему:
\begin{itemize}
    \item \textbf{RPFT} отвечает за общую архитектурную и геометрическую постановку;
    \item \textbf{UGSM} отвечает за строгую спектральную редукцию и вычислимые следствия;
    \item \textbf{TOE} отвечает за корреляционное и целостное чтение общего источника.
\end{itemize}

При таком понимании между ними исчезает значительная часть кажущихся противоречий. Они возникают в основном тогда, когда одну из этих форм пытаются принять за полное описание целого. Но если считать их разными сечениями одной сущности, то различия начинают работать как дополнения.

\subsection{Почему эта карта важна}

Эта карта важна не только как удобная схема. Она даёт методологическое правило на будущее. Всякий раз, когда между корпусами возникает напряжение, нужно спрашивать не только о том, кто прав, но и о том, не принадлежат ли сравниваемые утверждения к разным режимам чтения. Возможно, одно утверждение относится к уровню общей архитектуры, другое --- к уровню строгой спектральной редукции, а третье --- к уровню корреляционного смысла. В таком случае они не обязаны быть тождественными по формулировке, чтобы быть совместимыми по основанию.

Именно этот принцип и делает возможным следующий шаг: поиск не ``лучшой'' из имеющихся теорий, а того общего объекта, из которого все три режима чтения выводятся как естественные редукции.

\section{Первый набросок минимальной формы}

Философская и методологическая часть трактата уже позволяют сделать следующий рискованный, но необходимый шаг: попытаться в самом общем виде наметить, какой тип математического объекта мог бы играть роль искомой глубинной сущности.

Разумеется, речь пока не идёт о завершённой формализации. На текущем этапе можно лишь выделить требования к форме такого объекта.

\subsection{Требование 1: объект должен нести структуру состояний}

Прежде всего, искомый объект не может быть просто геометрическим многообразием в узком смысле. Он должен кодировать пространство допустимых состояний или мод, а значит включать в себя не только форму пространства, но и закон, по которому различаются и организуются допустимые режимы. В этом смысле объект должен быть по меньшей мере \textit{структурой состояний}, а не только формой протяжённости.

\subsection{Требование 2: объект должен порождать операторное чтение}

Далее, он должен естественно вести к оператору или семейству операторов, через которые становятся доступными спектр, след, тепло-ядерные и асимптотические коэффициенты. Без этого он не сможет породить строгий сектор и останется лишь философской картиной.

Следовательно, искомая сущность должна быть операторно-читаемой. Иначе говоря, её математическая форма обязана иметь такой режим, где она проявляется как источник спектральной структуры.

\subsection{Требование 3: объект должен нести глобальную фазовую и топологическую память}

Наконец, объект должен удерживать не только локальные моды, но и глобальные различия между ветвями. Это означает, что в его структуре должны быть различимы:
\begin{itemize}
    \item топологические классы;
    \item фазовые сектора;
    \item голономные структуры;
    \item внутренние ветви редукции.
\end{itemize}

Без этой памяти объект не сможет объяснить, почему глобальные фазы и внутренние секторного расщепленияs оказываются устойчивыми и физически содержательными.

\subsection{Рабочее предположение о форме}

В первом приближении можно предположить, что искомая сущность должна быть описываема как \textit{операторно-организованная структура на допустимом пространстве состояний с геометрической и топологической редукцией}. Это определение пока намеренно широко. Оно оставляет открытым вопрос, является ли объект единичным оператором, семейством операторов, категорией редукций, спектральной мерой или более сложной смешанной структурой. Но оно уже задаёт направление поиска.

Смысл этого предположения таков: то, что в геометрическом языке выглядит как компактное многообразие с инвариантами, в спектральном языке должно появляться как след некоторой операторной организации, а в корреляционном языке --- как проекция единой согласованной ткани состояний.

\section{Формальная постановка задачи}

После философского, онтологического и методологического развёртывания трактат должен наконец перейти к первой строгой формулировке. На этом этапе мы ещё не фиксируем окончательную теорию. Мы вводим минимальный объект и минимальный набор требований, которым он обязан удовлетворять, чтобы вообще быть кандидатом на роль глубинной сущности мира.

\subsection{Базовый объект}

Обозначим через
\[
\mathfrak{S}
\]
искомую глубинную сущность. На текущем этапе \(\mathfrak{S}\) понимается как минимальная структурно-спектральная организация допустимых состояний, обладающая тремя согласованными режимами чтения:
\begin{equation}
\mathcal{R}_g(\mathfrak{S}),\qquad \mathcal{R}_s(\mathfrak{S}),\qquad \mathcal{R}_c(\mathfrak{S}),
\end{equation}
где:
\begin{itemize}
    \item \(\mathcal{R}_g\) есть геометрическая редукция;
    \item \(\mathcal{R}_s\) есть спектральная редукция;
    \item \(\mathcal{R}_c\) есть корреляционная редукция.
\end{itemize}

Здесь слово ``редукция'' не означает потери сущности в отрицательном смысле. Напротив, оно означает, что единый объект допускает несколько согласованных способов проявления, каждый из которых подчеркивает свой аспект, не исчерпывая целого.

\subsection{Минимальный набор постулатов}

В качестве первой рабочей системы примем следующие постулаты.

\subsubsection*{Постулат I: постулат допустимости}

Существует пространство допустимых состояний
\[
\mathcal{H}_{adm},
\]
определяемое не произвольным перечислением конфигураций, а внутренним законом согласованности, задаваемым самой структурой \(\mathfrak{S}\).

Смысл этого постулата состоит в том, что мир не строится из всех логически мыслимых конфигураций. Он строится только из тех, которые совместимы с глубинной структурой допустимости.

\subsubsection*{Постулат II: постулат операторной порождаемости}

Сущность \(\mathfrak{S}\) индуцирует по меньшей мере один эффективный оператор
\[
\mathcal{D}_{\mathfrak{S}}
\]
или семейство операторов
\[
\{\mathcal{D}_{\mathfrak{S}}^{(a)}\}_{a\in A},
\]
действующих на \(\mathcal{H}_{adm}\), так что спектральная редукция определяется через спектр, следы и асимптотики этих операторов.

Этот постулат фиксирует принципиальный переход от философии к вычислимости. Без него нельзя получить ни спектральные коэффициенты, ни операторных структуру, ни связь с прежними проверками.

\subsubsection*{Постулат III: постулат геометрической реализуемости}

Существует такой режим чтения \(\mathcal{R}_g(\mathfrak{S})\), в котором часть структуры \(\mathfrak{S}\) выражается как геометрический каркас \((M,K,\Gamma,\Sigma,\ldots)\), где становятся осмысленными объёмы, кривизны, топологические классы, спин-структуры, циклы и компактные секторы.

Здесь важно, что геометрия не постулируется как единственная реальность. Она возникает как один из привилегированных режимов выражения глубинной структуры.

\subsubsection*{Постулат IV: постулат корреляционной целостности}

Существует такой режим чтения \(\mathcal{R}_c(\mathfrak{S})\), в котором различные локальные и глобальные проявления не рассматриваются как независимые объекты, а организуются как части единого корреляционного целого. Иначе говоря, структура \(\mathfrak{S}\) должна содержать механизм, делающий разные сектора совместимыми как проявления одного источника.

Этот постулат необходим для того, чтобы не свести всю программу к набору не связанных между собой операторных совпадений.

\subsubsection*{Постулат V: постулат инвариантного ядра}

Существует непустой набор величин, отношений или структур
\[
\mathcal{I}(\mathfrak{S})
\]
таких, что они сохраняют идентичность при переходе между \(\mathcal{R}_g\), \(\mathcal{R}_s\) и \(\mathcal{R}_c\). Именно \(\mathcal{I}(\mathfrak{S})\) и составляет инвариантное ядро глубинной сущности.

Без этого постулата переходы между режимами были бы лишь метафорическими аналогиями. С ним возникает возможность говорить о подлинной единой сущности.

\section{Первая система обозначений}

Для дальнейшей работы полезно зафиксировать минимальный словарь обозначений, который позволит не смешивать уровни анализа.

\begin{itemize}
    \item \(\mathfrak{S}\) --- глубинная структурно-спектральная сущность.
    \item \(\mathcal{H}_{adm}\) --- пространство допустимых состояний.
    \item \(\mathcal{D}_{\mathfrak{S}}\) --- операторный генератор спектрального чтения.
    \item \(\mathcal{R}_g\) --- геометрическая редукция.
    \item \(\mathcal{R}_s\) --- спектральная редукция.
    \item \(\mathcal{R}_c\) --- корреляционная редукция.
    \item \(\mathcal{I}(\mathfrak{S})\) --- инвариантное ядро, сохраняющееся при переходе между режимами чтения.
    \item \(\mathcal{B}\) --- сектор опорного каркаса, то есть устойчивый геометрически-спектральный каркас.
    \item \(\Phi\) --- фазово-голономный сектор.
    \item \(\mathcal{L}_{int}\) --- сектор внутренней нагрузки, отвечающий за внутренние субведущие структуры и секторные деформации.
\end{itemize}

В дальнейшем предполагается, что
\begin{equation}
\mathcal{I}(\mathfrak{S}) \supset \{\mathcal{B},\Phi,\mathcal{L}_{int}\},
\end{equation}
хотя точное определение этих компонент ещё должно быть выведено, а не просто постулировано.

\section{Первая строгая задача трактата}

После введения базового объекта и системы постулатов можно сформулировать первую строгую задачу.

\begin{quote}
Найти такую минимальную реализацию \(\mathfrak{S}\), для которой геометрическая, спектральная и корреляционная редукции воспроизводят одно и то же инвариантное ядро, а соответствующие наблюдаемые величины и секторные структуры оказываются естественными проявлениями этого ядра.
\end{quote}

Эта задача имеет несколько уровней сложности. На слабом уровне достаточно показать внутреннюю согласованность схемы. На среднем уровне --- воспроизвести уже известные устойчивые результаты. На сильном уровне --- получить новые нетривиальные следствия. Вся дальнейшая программа трактата будет строиться вокруг постепенного перехода от слабого уровня к сильному.

\section{Почему именно такая постановка минимальна}

Может возникнуть вопрос: не является ли предложенный формализм слишком широким? Ответ состоит в том, что на данной стадии он намеренно минимален, а не максимален. Мы не вводим здесь лишних полей, лишних симметрий, лишних измерений и лишних уравнений. Вместо этого мы фиксируем лишь те требования, без которых нельзя одновременно удержать:
\begin{itemize}
    \item геометрическое происхождение ограничений;
    \item спектральную вычислимость;
    \item корреляционную целостность;
    \item существование инвариантного ядра.
\end{itemize}

Если дальнейший анализ покажет, что для реализации \(\mathfrak{S}\) требуются дополнительные структуры, они должны будут вводиться не заранее, а только по необходимости. Такой принцип экономии и будет защищать трактат от избыточной спекуляции.

\section{Первая пробная реализация сущности \texorpdfstring{$\mathfrak{S}$}{S}}

После формальной постановки следующий шаг состоит в том, чтобы предложить хотя бы первую пробную реализацию искомой сущности. Она ещё не обязана быть окончательной. Но она должна быть достаточно содержательной, чтобы:
\begin{itemize}
    \item связывать уже накопленные геометрические и спектральные результаты;
    \item давать осмысленную интерпретацию фазово- и секторных структур;
    \item позволять начать новый анализ не с пустого места, а с минимального рабочего объекта.
\end{itemize}

\subsection{Пробное определение}

В первом приближении будем рассматривать \(\mathfrak{S}\) как тройку вида
\begin{equation}
\mathfrak{S} = (\mathcal{H}_{adm},\,\mathcal{D}_{\mathfrak{S}},\,\mathfrak{C}),
\end{equation}
где:
\begin{itemize}
    \item \(\mathcal{H}_{adm}\) есть пространство допустимых состояний;
    \item \(\mathcal{D}_{\mathfrak{S}}\) есть операторный генератор спектрального чтения;
    \item \(\mathfrak{C}\) есть правило корреляционной сборки, задающее, каким образом различные сектора интерпретируются как части единой структуры.
\end{itemize}

Это пробное определение важно тем, что оно сразу соединяет три уровня, которые ранее рассматривались раздельно. Пространство состояний отвечает за допустимость. Оператор отвечает за строгий спектральный режим. Корреляционная сборка отвечает за целостность и за то, почему разные операторных и геометрических проявления принадлежат одному объекту.

\subsection{Геометрическая реализация пробной сущности}

В геометрическом режиме данная тройка должна допускать редукцию к структуре вида
\begin{equation}
\mathcal{R}_g(\mathfrak{S}) = (M_4, K, \Gamma, \Sigma),
\end{equation}
где:
\begin{itemize}
    \item \(M_4\) обозначает эффективный внешний макроскопический сектор;
    \item \(K\) обозначает компактное внутреннее пространство;
    \item \(\Gamma\) кодирует топологическую или факторизационную структуру;
    \item \(\Sigma\) кодирует спиновых и фазово-чувствительных выборов редукции.
\end{itemize}

На этой стадии мы ещё не утверждаем, что такая форма является единственной. Однако она позволяет сразу встроить в схему уже накопленный опыт: компактные секторы, топологические классы, спин-структуры, minimal cycles, фактор-пространства и различия, порождаемые голономией.

Именно в этом режиме появляются объёмные и кривизновые инварианты, а также глобальные геометрические следы, которые позже должны быть прочитаны спектрально.

\subsection{Спектральная реализация пробной сущности}

В спектральном режиме та же сущность должна допускать описание через операторный след. В первом приближении это означает наличие объекта вида
\begin{equation}
K_{\mathfrak{S}}(t)=\mathrm{Tr}\,e^{-t\mathcal{D}_{\mathfrak{S}}^2},
\end{equation}
или, более общо,
\begin{equation}
S_f(\mathfrak{S};\Lambda)=\mathrm{Tr}\,f\!\left(\frac{\mathcal{D}_{\mathfrak{S}}^2}{\Lambda^2}\right),
\end{equation}
где разные выборы функции \(f\) соответствуют различным функциональным окнам чтения одной и той же структурной основы.

В таком режиме инвариантное ядро должно проявляться как устойчивость определённых комбинаций коэффициентов, асимптотик и их масштабирования. Иначе говоря, если \(\mathfrak{S}\) действительно существует как единая сущность, то её спектральное чтение не должно быть хаотичным: оно обязано воспроизводить повторяющиеся структурные мотивы при смене режимов, окон и деформаций.

\subsection{Корреляционная реализация пробной сущности}

Корреляционная часть пробной реализации связана с объектом \(\mathfrak{C}\). В минимальном виде его можно понимать как правило, которое сопоставляет различным редукциям не независимые интерпретации, а совместные проявления. Формально это можно записать схематически как отображение
\begin{equation}
\mathfrak{C}: \bigl(\mathcal{R}_g(\mathfrak{S}),\mathcal{R}_s(\mathfrak{S})\bigr) \longrightarrow \mathcal{R}_c(\mathfrak{S}),
\end{equation}
то есть корреляционный слой строится не отдельно, а как согласование геометрического и спектрального чтений.

Это особенно важно. Корреляционный язык не должен висеть в воздухе как независимая философская надстройка. Напротив, он должен рождаться именно как целостная интерпретация того, что геометрия и спектр уже предъявили по отдельности. Тогда он становится не украшением, а синтетическим уровнем.

\section{Секторная декомпозиция пробной сущности}

На следующем шаге полезно не просто говорить о едином объекте вообще, а сразу ввести его внутреннее разбиение на три предварительных сектора. Это разбиение пока не является окончательным, но оно естественно продолжает уже выявленные различия.

\subsection{Сектор опорного каркаса}

Обозначим через \(\mathcal{B}\) сектор опорного каркаса --- ту часть структуры \(\mathfrak{S}\), которая отвечает за устойчивый геометрически-спектральный каркас. Именно здесь должны находиться те инварианты, которые мало чувствительны к вторичным деформациям и выражают саму опорную архитектуру мира.

Именно сектор опорного каркаса должен объяснять, почему некоторые объёмные и кривизновые комбинации оказываются не случайными числами, а повторяющимися отпечатками одной и той же глубинной структуры.

\subsection{Фазовый сектор}

Обозначим через \(\Phi\) фазовый сектор --- ту часть структуры \(\mathfrak{S}\), которая отвечает за голономии, дискретные ветви, фазовые различия, сдвиги, чувствительные к спин-структуре и глобальные циклические следы. Этот сектор особенно важен потому, что именно он связывает локально неразличимые режимы с глобально различными классами допустимости.

Если сектор опорного каркаса отвечает за геометрически устойчивый каркас, то фазовый сектор отвечает за то, как этот каркас допускает несколько глобально различных, но структурно связанных ветвей существования.

\subsection{Сектор внутренней нагрузки}

Обозначим через \(\mathcal{L}_{int}\) сектор внутренней нагрузки --- ту часть структуры \(\mathfrak{S}\), которая кодирует внутренние субведущие нагрузки, мультиплетную структуру, более богатые внутренние каналы и их влияние на субведущие наблюдаемые величины. Этот сектор не обязан влиять на опорный каркас столь же сильно, как фазовый сектор влияет на глобальные ветви, но он должен отвечать за различия между внутренне обогащёнными режимами одной и той же общей структуры.

Именно этот сектор позволяет переходить от минимального скелета к более реалистичной внутренней организации без разрушения общего каркаса.

\subsection{Схема секторального взаимодействия}

В первом приближении можно считать, что
\begin{equation}
\mathfrak{S} \sim \mathcal{B} \oplus \Phi \oplus \mathcal{L}_{int},
\end{equation}
где символ \(\oplus\) пока следует понимать не как строгое прямое алгебраическое разложение, а как указание на различимые, но согласованные компоненты одной общей структуры.

Тогда дальнейшая задача заключается в том, чтобы определить:
\begin{itemize}
    \item какие наблюдаемые величины преимущественно читают \(\mathcal{B}\);
    \item какие --- \(\Phi\);
    \item какие --- \(\mathcal{L}_{int}\);
    \item и какие величины являются смешанными следами нескольких секторов сразу.
\end{itemize}

\section{Первое сближение с уже известными результатами}

Теперь можно сделать осторожный, но принципиальный шаг: соотнести пробную сущность \(\mathfrak{S}\) с уже накопленными структурными результатами, не утверждая пока окончательного тождества.

\subsection{Геометрические инварианты как след сектор опорного каркаса}

Если пробная схема верна, то устойчивые геометрические и квазигеометрические комбинации следует читать как проявления сектор опорного каркаса. Это означает, что устойчивые объёмные и кривизновые структуры являются не внешними коэффициентами, а следами опорной архитектуры \(\mathfrak{S}\).

\subsection{Holonomy и фазовые ветви как след фазовый сектор}

Глобальные фазовые сдвиги, ветви, чувствительные к спин-структуре, и различия, порождаемые голономией, естественно читаются как проявления фазового сектора. Их смысл тогда состоит не в том, что к теории добавлен вспомогательный параметр, а в том, что сама глубинная сущность допускает несколько глобально различных режимов обхода и чтения.

\subsection{Внутренние более богатые внутренние структуры как след сектор внутренней нагрузки}

Там, где возникают более богатые операторные картины, внутренние multiplets и субведущие нагрузки, естественно видеть проявление сектор внутренней нагрузки. Этот сектор не отменяет опорный каркас и не подменяет фазовый сектор, а усложняет внутреннюю организацию, сохраняя при этом более глубокую идентичность объекта.

\subsection{Промежуточный вывод}

Тем самым пробная реализация \(\mathfrak{S}\) уже начинает выполнять свою роль. Она позволяет не просто перечислить ранее найденные результаты, а распределить их по внутренней структуре одного объекта. Даже если эта схема ещё далека от окончательной строгости, она обладает важным достоинством: она делает видимым, как различные линии прежнего исследования могут быть поняты как следы одной глубинной организации.

\section{Первая рабочая программа вывода наблюдаемые величины}

После введения пробной сущности необходимо ответить на следующий естественный вопрос: каким образом из \(\mathfrak{S}\) должны извлекаться наблюдаемые структуры? Без ответа на этот вопрос трактат не сможет перейти от метафизики и общей схемы к содержательной физике.

Здесь под наблюдаемые величины понимаются не только экспериментальные числа в узком смысле. Речь идёт шире: об устойчивых коэффициентах, спектральных комбинациях, фазовых ветвях, масштабированиях и секторных различиях, которые можно интерпретировать как читаемые следы глубинной сущности.

\subsection{Общий принцип чтения}

В качестве базового принципа примем следующее:
\begin{quote}
наблюдаемая величина возникает как стабильный след одной или нескольких редукций сущности \(\mathfrak{S}\), а её смысл определяется тем, какой сектор инвариантного ядра в ней доминирует.
\end{quote}

Это означает, что ни одна важная величина не должна рассматриваться как ``просто число''. Каждое существенное наблюдаемое должно отвечать на три вопроса:
\begin{enumerate}
    \item из какого режима чтения оно извлекается;
    \item какой сектор \(\mathfrak{S}\) в нём доминирует;
    \item является ли оно чистым следом одного сектора или смешанным следом нескольких.
\end{enumerate}

\subsection{Наблюдаемые величины опорного каркаса}

Назовём опорный каркас-наблюдаемые величины такие наблюдаемые, которые прежде всего читают \(\mathcal{B}\), то есть устойчивый каркас сущности. Их отличительный признак состоит в слабой чувствительности к вторичным внутренним деформациям и в относительной устойчивости при смене окна чтения.

В первом приближении к ним должны относиться:
\begin{itemize}
    \item объёмные инварианты;
    \item кривизновые комбинации, устойчивые в малой асимптотике;
    \item ведущие спектральные коэффициенты;
    \item устойчивые геометрические масштабирования.
\end{itemize}

Если новая программа верна, именно такие наблюдаемые величины должны проявлять наибольшую универсальность при переходе между различными моделями и процедурами чтения.

\subsection{Фазовые наблюдаемые величины}

Назовём фазово-наблюдаемые величины такие наблюдаемые, которые преимущественно чувствительны к \(\Phi\). Их отличительный признак --- слабая зависимость от чисто локальной геометрической деформации при сильной зависимости от глобальной ветви, класса голономии, спин-структуры или циклического обхода.

В эту группу в первом приближении должны входить:
\begin{itemize}
    \item фазовые сдвиги;
    \item ветви Вильсона и голономии;
    \item глобальные различия, чувствительные к спин-структуре;
    \item наблюдаемые величины, связанные с циклическим или нестягиваемым сектором.
\end{itemize}

Эти величины особенно важны для трактата, поскольку они выражают то, что не может быть восстановлено чисто локальным чтением. Именно через них глубинная сущность демонстрирует свою глобальную память.

\subsection{Наблюдаемые величины сектора внутренней нагрузки}

Назовём внутренней нагрузки наблюдаемые величины такие наблюдаемые, которые преимущественно читают \(\mathcal{L}_{int}\). Это наблюдаемые величины, чувствительные к внутренним каналам, мультиплетную структуру, субведущие деформациям и эффекты более богатого операторного содержания.

Их отличительный признак состоит в том, что они не обязаны радикально менять опорный каркас, но могут существенно перестраивать подведущую структуру, внутренние отношения между каналами и величины, зависящие от скрытой внутренней организации.

В первом приближении к таким наблюдаемые величины относятся:
\begin{itemize}
    \item субведущие spectral coefficients;
    \item внутренние комбинации, чувствительные к внутренней нагрузке;
    \item параметры секторного расщепления;
    \item наблюдаемые величины, меняющиеся при обогащении внутреннего операторного содержания.
\end{itemize}

\subsection{Смешанные наблюдаемые величины}

Не все наблюдаемые обязаны быть чистыми следами одного сектора. Наоборот, наиболее физически интересные величины могут оказаться смешанными. Тогда задача теории состоит не в грубом отнесении каждой величины к одному сектору, а в выделении доминирующего вклада и описании смешения.

Схематически это можно записать так:
\begin{equation}
\mathcal{O} = \mathcal{O}_{\mathcal{B}} + \mathcal{O}_{\Phi} + \mathcal{O}_{\mathcal{L}_{int}} + \mathcal{O}_{\mathrm{mix}},
\end{equation}
где \(\mathcal{O}_{\mathrm{mix}}\) обозначает несводимые перекрёстные члены. Такая запись пока лишь эвристична, но она задаёт правильное направление: наблюдаемое должно быть декомпозируемо по секторам хотя бы в первом приближении.

\section{Правило секторной атрибуции}

Чтобы новая программа могла работать системно, необходимо зафиксировать первое правило атрибуции.

\begin{quote}
Если наблюдаемая величина остаётся устойчивой при внутренних деформациях и зависит главным образом от геометрического каркаса, она атрибутируется сектору опорного каркаса. Если она радикально зависит от глобальной фазы, ветви или класса голономии, она атрибутируется фазовому сектору. Если она преимущественно чувствительна к внутренней мультиплетной или субведущей нагрузке, она атрибутируется сектору внутренней нагрузки.
\end{quote}

Это правило не является окончательной теоремой, но оно уже даёт инструмент для нового анализа. С его помощью можно повторно прочитывать уже накопленные результаты и проверять, действительно ли они естественно раскладываются по предложенной структуре.

\subsection{Сильная и слабая формы атрибуции}

Полезно различать два режима секторной атрибуции.

\textbf{Слабая атрибуция} утверждает лишь, что существует доминирующий сектор, объясняющий основную часть поведения наблюдаемой величины.

\textbf{Сильная атрибуция} требует большего: она утверждает, что величина может быть не только качественно, но и количественно декомпозирована по секторам с малым остатком смешения.

На данной стадии трактат работает в режиме слабой атрибуции. Переход к сильной атрибуции должен стать задачей следующего цикла исследований.

\section{Первый тест новой схемы}

Теперь можно зафиксировать самый простой тест, которому должна удовлетворять предложенная программа.

\begin{quote}
Если пробная сущность \(\mathfrak{S}\) и её секторное разбиение осмысленны, то уже известные устойчивые результаты должны допускать более простую и более ясную классификацию, чем до введения этой схемы.
\end{quote}

Это очень важный критерий. На данном этапе новая теория ещё не обязана выводить все числа заново. Но она уже обязана упорядочивать известное лучше, чем это делала предыдущая фрагментированная картина. Если она этого не делает, значит введение \(\mathfrak{S}\), \(\mathcal{B}\), \(\Phi\) и \(\mathcal{L}_{int}\) является только красивым переименованием.

\subsection{Что именно нужно проверить сначала}

Первый цикл такого теста должен включать по меньшей мере пять вопросов:
\begin{enumerate}
    \item какие наблюдаемые устойчиво читают сектор опорного каркаса;
    \item какие наблюдаемые устойчиво читают фазовый сектор;
    \item какие наблюдаемые устойчиво читают сектор внутренней нагрузки;
    \item какие величины являются смешанными;
    \item удаётся ли благодаря этой карте устранить хотя бы часть прежней концептуальной неоднородности.
\end{enumerate}

Лишь после этого будет возможно переходить к более амбициозной задаче --- к новому самостоятельному выводу следствий из самой сущности \(\mathfrak{S}\), а не только к повторному чтению уже известных результатов.

\section{Переход к следующему этапу}

Тем самым трактат впервые достигает состояния, в котором уже можно начинать собственный исследовательский цикл. У нас есть:
\begin{itemize}
    \item философская интуиция;
    \item онтологическая схема;
    \item карта теоретических корпусов;
    \item минимальный формальный объект;
    \item секторная декомпозиция;
    \item правило чтения наблюдаемых величин.
\end{itemize}

Следующий этап должен состоять в том, чтобы перестать говорить о сущности в общем и начать строить её первый предметный анализ. Это будет означать переход от трактата как введения к трактату как новой исследовательской работе в строгом смысле.

\section{Первый предметный анализ сущности \texorpdfstring{$\mathfrak{S}$}{S}}

Теперь необходимо сделать первый шаг от общей схемы к рабочему анализу. Этот шаг ещё не даёт полной теории, но он вводит минимальную дисциплину рассуждения. Если ранее сущность \(\mathfrak{S}\) фигурировала прежде всего как объединяющий принцип, то теперь она должна быть рассмотрена как объект, для которого можно сформулировать по меньшей мере: допустимые деформации, устойчивые характеристики, режимы редукции и критерии тождества.

Именно на этом уровне начинается различие между сильной метафизической гипотезой и исследовательской программой в собственном смысле. Исследовательская программа обязана не только говорить, \textit{что} за различными описаниями мира стоит единый источник, но и указывать, \textit{как} отличать существенное в этом источнике от случайного, какие преобразования считать несущественными и какие следы признавать инвариантными.

\subsection{Минимальное требование к предметности}

Будем считать, что сущность \(\mathfrak{S}\) введена предметно лишь в том случае, если для неё можно указать:
\begin{enumerate}
    \item класс допустимых состояний или конфигураций;
    \item класс допустимых деформаций, не разрушающих её идентичность;
    \item семейство инвариантов, сохраняющихся при этих деформациях;
    \item набор режимов чтения, в которых эти инварианты становятся наблюдаемыми или квазинаблюдаемыми величинами.
\end{enumerate}

Это требование следует понимать как минимальный порог строгости. Пока оно не выполнено, разговор о \(\mathfrak{S}\) остаётся полезной эвристикой, но ещё не становится собственным языком теории.

\subsection{Тождество сущности и допустимые деформации}

Первый вопрос, который необходимо поставить, звучит так: когда две реализации \(\mathfrak{S}_1\) и \(\mathfrak{S}_2\) следует считать одной и той же сущностью, а когда --- уже различными сущностями? Без ответа на этот вопрос невозможно ни говорить об инвариантах, ни строить карту наблюдаемых величин.

Предварительно примем следующее рабочее правило:
\begin{quote}
две реализации принадлежат одной и той же сущности \(\mathfrak{S}\), если они связаны цепью допустимых деформаций, сохраняющих опорный каркас, секторную различимость и класс глобальной согласованности.
\end{quote}

Это правило сразу разделяет деформации на три типа.

\textbf{Слабыми деформациями} назовём такие изменения, которые не нарушают ни опорный каркас \(\mathcal{B}\), ни глобальный фазовый класс \(\Phi\), ни саму возможность устойчивого чтения внутренней нагрузки \(\mathcal{L}_{int}\). Они должны рассматриваться как внутренние вариации одной и той же сущности.

\textbf{Промежуточными деформациями} назовём такие изменения, которые сохраняют общую идентичность \(\mathfrak{S}\), но меняют относительный вес секторных вкладов и тем самым перестраивают карту наблюдаемых величин. Именно этот тип деформаций, по-видимому, будет ответственен за переходы между различными эффективными физическими режимами.

\textbf{Сильными деформациями} назовём такие изменения, при которых разрушается либо опорная инвариантная структура, либо фазовый класс, либо принципиальная возможность прежней секторной декомпозиции. В этом случае речь, вероятно, должна идти уже не о новой конфигурации той же сущности, а о переходе к иной сущности или иной области теории.

\subsection{Инвариантное ядро}

Чтобы сделать предыдущую схему содержательной, введём понятие \textit{инвариантного ядра} сущности. Обозначим его через \(\mathcal{I}(\mathfrak{S})\). Под ним будем понимать минимальный набор характеристик, который:
\begin{itemize}
    \item сохраняется при всех слабых деформациях;
    \item контролирует допустимую секторную декомпозицию;
    \item ограничивает спектр возможных наблюдаемых редукций.
\end{itemize}

На этой стадии нет необходимости задавать \(\mathcal{I}(\mathfrak{S})\) окончательной формулой. Достаточно зафиксировать его функциональную роль. Оно должно быть именно тем, что делает разные реализации узнаваемыми как чтения одного источника. Иначе говоря, если \(\mathfrak{S}\) есть глубинная сущность, то \(\mathcal{I}(\mathfrak{S})\) есть её минимально сохраняемая идентичность.

В первом приближении можно считать, что инвариантное ядро имеет по меньшей мере три компоненты:
\begin{equation}
\mathcal{I}(\mathfrak{S}) = \mathcal{I}_{\mathcal{B}} \oplus \mathcal{I}_{\Phi} \oplus \mathcal{I}_{\mathcal{L}_{int}},
\end{equation}
где каждая компонента отвечает не просто отдельному сектору, а тому, что в данном секторе неустранимо без разрушения общей идентичности сущности.

\subsection{Принцип редукционного чтения}

Теперь можно сформулировать первый рабочий принцип, связывающий глубинную сущность с наблюдаемым миром:
\begin{quote}
физически читаемая структура возникает не из всей сущности \(\mathfrak{S}\) сразу, а из её редукции относительно выбранного окна чтения, масштаба и допустимого класса деформаций.
\end{quote}

Это положение имеет решающее значение. Оно означает, что неудача некоторого языка полностью описать сущность не опровергает саму сущность. Она лишь показывает, что данный язык является частичной редукцией. Геометрический язык читает один класс следов, спектральный --- другой, корреляционный --- третий. Ошибка прежних фрагментированных подходов могла состоять не в том, что они были ложны, а в том, что каждый из них молчаливо абсолютизировал собственное окно чтения.

Отсюда следует важная методологическая норма: сравнивать разные теории следует не как конкурирующие описания целого, а как разные функторы чтения, действующие на одну и ту же сущность \(\mathfrak{S}\). Если два таких чтения дают согласованные инварианты, это следует считать аргументом в пользу единого источника. Если же они дают несовместимые картины, то необходимо выяснять, связано ли это с подлинной неоднозначностью сущности или лишь с нарушением области применимости одного из окон чтения.

\subsection{Первый критерий хорошей формализации}

Теперь можно ввести и первый критерий того, что формализация \(\mathfrak{S}\) выбрана удачно.

\begin{quote}
Хорошая формализация сущности \(\mathfrak{S}\) должна одновременно уменьшать число независимых исходных предположений и увеличивать число естественно объясняемых инвариантных связей между уже известными наблюдаемыми величинами.
\end{quote}

Этот критерий особенно важен, потому что он запрещает чисто декоративную общность. Если новое описание просто вводит более абстрактные обозначения, не уменьшая концептуальную раздробленность материала, то оно не усиливает теорию. Если же одна и та же структура начинает объяснять сразу спектральную устойчивость, фазовые различия и характер некоторых эффективных физических величин, тогда абстракция оправдана.

\subsection{Пробная схема первого анализа}

С практической точки зрения первый предметный анализ \(\mathfrak{S}\) должен идти по следующей схеме:
\begin{enumerate}
    \item выделить предполагаемое инвариантное ядро \(\mathcal{I}(\mathfrak{S})\);
    \item перечислить допустимые слабые деформации;
    \item определить, какие наблюдаемые сохраняются при этих деформациях;
    \item отделить их от величин, чувствительных к промежуточным деформациям;
    \item сравнить полученную карту с уже известными спектральными, геометрическими и фазовыми результатами.
\end{enumerate}

Эта схема ещё не есть полная теория вывода. Но она уже создаёт рабочий контур, внутри которого можно впервые проверять не только красоту общей гипотезы, но и её объяснительную дисциплину.

\subsection{Рабочая гипотеза следующего шага}

Если вся предыдущая конструкция верна хотя бы качественно, то следующий шаг должен состоять в следующем: выбрать один ограниченный класс моделей или один устойчивый класс уже полученных результатов и показать, что они действительно допускают реконструкцию через \(\mathcal{I}(\mathfrak{S})\), допустимые деформации и секторную редукцию.

Именно здесь трактат должен перейти от общего языка к проверяемому разбору. До сих пор сущность \(\mathfrak{S}\) вводилась как необходимая гипотеза единства. Начиная с этого места, она должна доказать своё право на существование уже не только философской мощностью, но и способностью организовывать конкретный материал.

Следовательно, ближайшая задача новой программы может быть сформулирована так: нужно построить первый локальный пример, в котором опорный каркас, фазовый сектор и сектор внутренней нагрузки не просто перечисляются, а действительно участвуют в совместном объяснении некоторого устойчивого класса наблюдаемых величин. Только после этого можно будет говорить, что теория начала собственное движение.

\section{Первый локальный пример чтения сущности \texorpdfstring{$\mathfrak{S}$}{S}}

Чтобы сделать программу содержательной, начнём не с максимально сложной физической системы, а с предельно простого, но методологически показательного класса ситуаций. Рассмотрим схему, в которой одна и та же глубинная сущность \(\mathfrak{S}\) редуцируется к эффективной картине движения мод на компактном носителе с нетривиальной глобальной фазой и ограниченной внутренней секторностью. Этот пример ещё не претендует на окончательную физическую модель. Его задача скромнее: показать, как три сектора новой схемы могут одновременно участвовать в порождении наблюдаемых различий.

\subsection{Постановка пробной модели}

Пусть фиксирован некоторый компактный носитель \(M\), на котором определён операторный режим чтения сущности \(\mathfrak{S}\). Будем считать, что в первом приближении эта сущность индуцирует:
\begin{enumerate}
    \item опорный каркас \(\mathcal{B}\), задающий класс допустимой геометрии и граничных ограничений;
    \item фазовый сектор \(\Phi\), задающий допустимую глобальную ветвь переноса;
    \item внутренний сектор \(\mathcal{L}_{int}\), задающий конечный набор внутренних каналов или мультиплетных компонент.
\end{enumerate}

В качестве наблюдаемого окна чтения выберем эффективный спектральный режим. Тогда с каждой реализацией \(\mathfrak{S}\) будет связан некоторый набор собственных частот, массоподобных масштабов или устойчивых коэффициентов отклика. Важно, что эти величины не вводятся независимо. Они должны рассматриваться как следы совместной редукции трёх секторов.

\subsection{Опорный каркас как источник ведущего масштаба}

Начнём с простейшего уровня. Если фиксировать фазовый класс и внутреннюю нагрузку, то именно опорный каркас должен задавать ведущую геометрию допустимых мод. Это означает, что грубая плотность спектра, ведущий масштаб разделения уровней и устойчивые асимптотические коэффициенты прежде всего читают \(\mathcal{B}\).

Схематически это можно выразить так:
\begin{equation}
\lambda_n \sim F_{\mathcal{B}}(n) + \delta_{\Phi}(n) + \delta_{int}(n),
\end{equation}
где главный вклад \(F_{\mathcal{B}}(n)\) определяется опорным каркасом, а члены \(\delta_{\Phi}(n)\) и \(\delta_{int}(n)\) являются поправками фазового и внутреннего происхождения. Само это разложение пока следует понимать эвристически. Его смысл в том, что геометрия каркаса задаёт не всю наблюдаемую структуру, но именно её ведущую несущую форму.

Иначе говоря, если два режима чтения различаются только малыми внутренними перестройками, то их ведущий спектральный профиль должен оставаться близким. В этом и выражается физическая роль опорного каркаса: он удерживает крупномасштабную идентичность семейства наблюдаемых.

\subsection{Фазовый сектор как источник глобального расщепления}

Теперь предположим, что геометрический каркас \(\mathcal{B}\) фиксирован, но меняется глобальная фазовая ветвь \(\Phi\). Тогда ведущая плотность состояний может остаться почти неизменной, однако часть спектра, фазовые сдвиги и правила секторного отбора должны перестраиваться. Именно это и делает фазовый сектор принципиально отличным от геометрического.

В простейшем виде эту мысль можно записать как зависимость эффективных уровней от класса голономии:
\begin{equation}
\lambda_n^{(\alpha)} = \lambda_n^{(0)} + \Delta_{\Phi}(\alpha,n),
\end{equation}
где параметр \(\alpha\) метит фазовую ветвь или класс обхода. При этом важен не конкретный аналитический вид функции \(\Delta_{\Phi}\), а сам принцип: глобальная фаза должна вызывать такие различия, которые не редуцируются к локальной деформации каркаса.

Это позволяет сразу сформулировать первый локальный диагностический признак фазового сектора: если некоторое наблюдаемое устойчиво к гладким локальным вариациям геометрии, но чувствительно к смене ветви переноса, то оно должно быть атрибутировано \(\Phi\). В уже накопленном материале именно такого типа признаки неоднократно возникали в виде голономных различий, фазовых расщеплений и зависимостей от нестягиваемых циклов.

\subsection{Внутренняя нагрузка как источник тонкой структуры}

Наконец, зафиксируем опорный каркас и фазовую ветвь и рассмотрим изменения, индуцируемые сектором \(\mathcal{L}_{int}\). Если предыдущая схема верна, то этот сектор не обязан радикально менять глобальную форму спектра, но должен перестраивать тонкую структуру: кратности, субведущие коэффициенты, относительные сдвиги каналов и характер некоторых смешанных корреляций.

Схематически это можно записать так:
\begin{equation}
\lambda_{n,a} = \lambda_n^{\mathrm{base}} + \epsilon_a(n), \qquad a \in \mathcal{A}_{int},
\end{equation}
где индекс \(a\) маркирует внутренний канал, а поправка \(\epsilon_a(n)\) отражает внутреннюю секторность сущности. Тогда то, что в эффективной физике описывается как расщепление уровней, мультиплетность или появление дополнительных селекционных правил, может быть интерпретировано не как внешний довесок к геометрии, а как прямое проявление внутренней нагрузки глубинной сущности.

Именно здесь новая программа получает шанс объяснить то, что часто выглядит как искусственная надстройка над геометрической картиной. Если внутренний сектор действительно принадлежит самой сущности \(\mathfrak{S}\), то часть эффективных ``внутренних квантовых чисел'' должна пониматься как след её скрытой организации, а не как посторонний список добавленных ярлыков.

\subsection{Смешанное наблюдаемое как первый нетривиальный тест}

Наиболее важен, однако, не случай чистых вкладов, а случай смешанного наблюдаемого. Рассмотрим величину \(\mathcal{O}\), которая зависит одновременно от ведущего масштаба каркаса, от глобальной фазовой ветви и от внутреннего секторного расщепления. Тогда в первом порядке можно ожидать схему вида:
\begin{equation}
\mathcal{O} = \mathcal{O}_0(\mathcal{B}) + \mathcal{O}_1(\Phi) + \mathcal{O}_2(\mathcal{L}_{int}) + \mathcal{O}_{12}(\mathcal{B},\Phi) + \mathcal{O}_{23}(\Phi,\mathcal{L}_{int}) + \cdots .
\end{equation}

Смысл этой записи состоит не в том, чтобы сразу получить готовую вычислительную формулу, а в том, чтобы дисциплинировать анализ. Теперь любое устойчивое наблюдаемое можно рассматривать как вопрос: какая часть его поведения контролируется каркасом, какая --- фазой, какая --- внутренней нагрузкой, и где возникают несводимые перекрёстные члены.

Если новая схема плодотворна, то именно такие смешанные величины должны стать её главным полигоном. На чистых величинах теория лишь подтверждает уже ожидаемое. На смешанных же она либо показывает свою объяснительную силу, либо распадается на неработающее переименование старых различий.

\subsection{Первое следствие для физической интерпретации}

Из пробной модели уже следует важное следствие. Масса, заряд, фазовая ветвь и внутренняя мультиплетность не должны больше мыслиться как полностью независимые исходные данные. По крайней мере в рамках рассматриваемой программы они должны интерпретироваться как разные проекции единого инвариантного ядра \(\mathcal{I}(\mathfrak{S})\) в различных окнах чтения.

Это не означает, что все физические величины немедленно сводятся друг к другу. Речь о другом: их независимость перестаёт быть методологически исходной. Вместо этого возникает задача реконструкции, то есть задача понять, какие из них совместно кодируют одну и ту же глубинную структуру, а какие действительно принадлежат разным уровням описания.

\subsection{Какой именно тест должен следовать дальше}

Теперь можно указать и следующий уже вполне конкретный шаг. Нужно выбрать один класс ранее полученных результатов --- например, семейство спектральных коэффициентов с фазовым расщеплением или класс геометрических конфигураций с устойчивой внутренней мультиплетностью --- и проверить, допускает ли он систематическое разложение по трём секторам:
\begin{equation}
\text{данные} \longrightarrow (\mathcal{B},\Phi,\mathcal{L}_{int}).
\end{equation}

Если такое разложение окажется устойчивым и не произвольным, то теория сделает первый настоящий шаг к предметности. Если же один и тот же результат можно будет объяснить одинаково хорошо бесконечным числом несовместимых секторных карт, это будет означать, что понятие сущности \(\mathfrak{S}\) пока ещё не достигло достаточной строгости.

\subsection{Промежуточный вывод}

Итак, первый локальный пример показывает следующее. Новая сущность \(\mathfrak{S}\) может быть понята не как туманная метафизическая подложка, а как принцип организации, чьи три сектора читаются в наблюдаемом мире различными, но взаимосвязанными способами. Опорный каркас удерживает ведущую форму допустимых режимов. Фазовый сектор кодирует глобальные различия, неустранимые локальной деформацией. Внутренняя нагрузка отвечает за тонкую структуру и дополнительную канализацию наблюдаемых.

Тем самым мы получаем первый образ той физики, которую ищет трактат: физики, где наблюдаемое не просто извлекается из готовой геометрии или из отдельного поля, а возникает как структурированный след глубинной сущности, читаемой через несколько совместимых редукций. Следующий шаг должен состоять в том, чтобы превратить эту пробную схему в более строгий протокол реконструкции.

\section{Протокол реконструкции наблюдаемых из \texorpdfstring{$\mathfrak{S}$}{S}}

Теперь необходимо перевести общую схему и локальный пример в форму рабочего исследовательского протокола. Под протоколом реконструкции будем понимать упорядоченную процедуру, которая позволяет не просто описывать наблюдаемые величины постфактум, а систематически ставить вопрос о том, какие компоненты глубинной сущности \(\mathfrak{S}\) ответственны за их устойчивые свойства. Иначе говоря, задача состоит в том, чтобы двигаться не от готовой интерпретации к данным, а от данных к допустимой структуре сущности.

\subsection{Исходная установка протокола}

Предположим, что дан некоторый класс наблюдаемых величин или устойчивых результатов \(\mathcal{D}\). Это могут быть спектральные коэффициенты, геометрические инварианты, фазовые расщепления, корреляционные структуры или эффективные параметры модели. Протокол реконструкции должен отвечать на следующий вопрос:
\begin{quote}
какая минимальная секторная структура сущности \(\mathfrak{S}\) необходима, чтобы объяснить устойчивую организацию данных \(\mathcal{D}\)?
\end{quote}

Уже сама такая постановка меняет направление анализа. Мы больше не спрашиваем только о том, какой формальный аппарат удобен для описания данных. Мы спрашиваем, какой тип глубинной организации обязан стоять за ними, если они действительно являются следами одной и той же согласованной сущности.

\subsection{Шаг первый: выделение устойчивого ядра данных}

Первый шаг состоит в том, чтобы отделить устойчивое содержание данных от случайных или модельно-зависимых деталей. Обозначим это устойчивое содержание через \(\mathcal{D}_{\mathrm{stab}}\). В него должны входить только те свойства, которые:
\begin{itemize}
    \item воспроизводятся в нескольких режимах вычисления или описания;
    \item не исчезают при малых допустимых деформациях модели;
    \item допускают осмысленную интерпретацию как инвариантные или квазиинвариантные следы.
\end{itemize}

Этот шаг принципиален. Пока не выделено \(\mathcal{D}_{\mathrm{stab}}\), реконструкция остаётся бессодержательной, потому что теория начинает подгоняться под шум. Следовательно, первая добродетель новой программы --- не абстрактность, а дисциплина отбора.

\subsection{Шаг второй: первичная секторная атрибуция}

После выделения устойчивого ядра необходимо провести первичную атрибуцию наблюдаемых. Для каждого элемента \(O \in \mathcal{D}_{\mathrm{stab}}\) ставятся три вопроса:
\begin{enumerate}
    \item чувствителен ли он прежде всего к опорному каркасу;
    \item чувствителен ли он прежде всего к глобальной фазе;
    \item чувствителен ли он прежде всего к внутренней нагрузке.
\end{enumerate}

Результат этого шага удобно записывать в виде пробного отображения
\begin{equation}
\mathfrak{A}_0 : \mathcal{D}_{\mathrm{stab}} \longrightarrow \{\mathcal{B},\Phi,\mathcal{L}_{int},\mathrm{mix}\}.
\end{equation}

Здесь индекс \(0\) подчёркивает, что речь идёт пока только о слабой атрибуции. Она не обязана давать количественное разложение, но обязана фиксировать доминирующий характер вклада. Если даже такая слабая карта не удаётся, теория ещё не получила права на дальнейшую конкретизацию.

\subsection{Шаг третий: реконструкция инвариантного ядра}

После первичной атрибуции нужно поставить более трудный вопрос: существует ли такой минимальный набор структурных характеристик, из которого полученная карта наблюдаемых следует естественно? Именно этот набор и должен играть роль кандидатного инвариантного ядра \(\mathcal{I}(\mathfrak{S})\).

На этом этапе реконструкция движется в обратном направлении по сравнению с обычным моделированием. Мы не задаём заранее полный объект и затем считаем его следствия. Мы ищем минимальную структуру, необходимую для объяснения уже обнаруженного набора устойчивых следов. Формально это можно выразить как задачу найти такое \(\mathcal{I}^*\), что
\begin{equation}
\mathcal{I}^* \Rightarrow \mathfrak{A}_0(\mathcal{D}_{\mathrm{stab}})
\end{equation}
с минимальным числом независимых предположений.

Именно здесь проявляется главный критерий экономии: хорошая реконструкция не размножает сущности сверх необходимости, а наоборот, стягивает разные наблюдаемые к одному и тому же инвариантному источнику.

\subsection{Шаг четвёртый: проверка деформационной устойчивости}

Однако одной только структурной экономии недостаточно. Кандидатное ядро должно пройти проверку на устойчивость при допустимых деформациях. Пусть \(\delta \in \Delta_{adm}\) --- допустимая деформация в смысле ранее введённой схемы. Тогда реконструкция считается содержательной только в том случае, если при действии \(\delta\) сохраняется либо сама карта атрибуции, либо по меньшей мере её доминирующий секторный профиль.

Схематически это требование можно записать так:
\begin{equation}
\mathfrak{A}_0(O) \approx \mathfrak{A}_0(\delta O), \qquad O \in \mathcal{D}_{\mathrm{stab}}, \; \delta \in \Delta_{adm}.
\end{equation}

Если же малая деформация радикально перестраивает всю карту секторной атрибуции, то перед нами либо неверно выделенное устойчивое ядро данных, либо слишком грубая гипотеза о сущности. В обоих случаях протокол обязан вернуть исследование на предыдущий шаг.

\subsection{Шаг пятый: различение чистых и смешанных величин}

Следующий шаг состоит в различении чистых и смешанных наблюдаемых. Это необходимо, потому что именно смешанные величины представляют главную трудность для всякой глубинной программы. Если для чистых величин ещё можно довольствоваться качественной атрибуцией, то смешанные требуют хотя бы приближённого выделения доминирующих и перекрёстных вкладов.

Пусть \(O_{mix}\) --- наблюдаемое, для которого карта \(\mathfrak{A}_0\) возвращает значение \(\mathrm{mix}\). Тогда нужно искать разложение вида
\begin{equation}
O_{mix} = O_{\mathcal{B}} + O_{\Phi} + O_{\mathcal{L}_{int}} + O_{\times},
\end{equation}
где \(O_{\times}\) содержит несводимые перекрёстные вклады. Даже если такое разложение не является строгим в аналитическом смысле, оно должно быть осмысленным в структурном отношении: теория обязана уметь указать, какой тип смешения имеет место и почему он не устраним более грубой редукцией.

\subsection{Шаг шестой: сравнение конкурирующих окон чтения}

После этого необходимо проверить, согласуются ли разные теоретические окна чтения в отношении реконструированного ядра. Пусть один и тот же класс данных допускает геометрическое, спектральное и корреляционное описание. Тогда хорошая теория не должна строить три несвязанных объяснения. Она должна показывать, что эти описания являются различными редукциями одной и той же реконструированной структуры.

Тем самым вводится условие межъязыковой согласованности:
\begin{equation}
R_{geom}(\mathcal{I}^*) \sim R_{spec}(\mathcal{I}^*) \sim R_{corr}(\mathcal{I}^*),
\end{equation}
где знак \(\sim\) обозначает не буквальное совпадение формул, а согласованность инвариантного содержания. Если это условие систематически нарушается, то либо \(\mathcal{I}^*\) выбрано неверно, либо один из языков применяется за пределами своей естественной области.

\subsection{Шаг седьмой: критерий непустоты теории}

Теперь можно сформулировать важнейший критерий того, что протокол действительно производит знание, а не просто организует уже известное. Реконструкция считается нетривиальной тогда и только тогда, когда из найденного инвариантного ядра следуют хотя бы некоторые ограничения на ещё не классифицированные наблюдаемые или на допустимые формы их смешения.

Иначе говоря, хорошая реконструкция должна работать не только назад, но и вперёд. Она обязана не просто объяснять задним числом, а ограничивать будущую интерпретацию. Если теория допускает всё, она не объясняет ничего.

\subsection{Итоговая схема протокола}

В собранном виде протокол реконструкции можно записать так:
\begin{align}
\mathcal{D}
    &\longrightarrow \mathcal{D}_{\mathrm{stab}} \\
    &\longrightarrow \mathfrak{A}_0(\mathcal{D}_{\mathrm{stab}}) \\
    &\longrightarrow \mathcal{I}^*(\mathfrak{S}) \\
    &\longrightarrow \text{проверка на } \Delta_{adm} \\
    &\longrightarrow \text{межъязыковая согласованность} \\
    &\longrightarrow \text{новые ограничения и предсказания}.
\end{align}

Эта схема ещё не есть окончательная математика новой теории. Но она уже задаёт исследовательский ритм, без которого теория не сможет перейти от вдохновляющей онтологии к строгой работе. С этого места трактат впервые получает форму процедуры, а не только форму мировоззрения.

\subsection{Переход к следующей задаче}

Следующий этап теперь ясен. Нужно взять один конкретный класс результатов и прогнать его через весь описанный протокол без пропуска шагов. Только такой полный прогон покажет, действительно ли \(\mathfrak{S}\) является рабочим понятием, а не просто общим именем для набора красивых интуиций.

Именно поэтому дальнейшая часть трактата должна перейти к примеру полного цикла: от данных --- к устойчивому ядру, от ядра --- к секторной атрибуции, от атрибуции --- к реконструкции инвариантной структуры, и затем --- к проверяемым следствиям. Лишь в этом случае новая программа начнёт становиться теорией в строгом смысле слова.


\section{Первый полный цикл реконструкции на пробном классе данных}

Теперь необходимо выполнить то, к чему подводили предыдущие разделы: не просто описать протокол в общем виде, а прогнать через него один конкретный пробный класс данных. Для этой цели выберем не максимально богатую систему, а достаточно простый, но содержательный класс спектрально-фазовых данных, в котором одновременно присутствуют: устойчивый ведущий масштаб, глобальное фазовое расщепление и тонкая внутренняя структура. Именно такой класс особенно удобен для первого полного цикла, потому что он сразу затрагивает все три сектора сущности \(\mathfrak{S}\).

\subsection{Выбор пробного класса данных}

Пусть дан набор наблюдаемых величин
\begin{equation}
\mathcal{D}^{(1)} = \{\rho(\Lambda),\, \Delta_{\alpha},\, m_a,\, C_{a\alpha}\},
\end{equation}
где \(\rho(\Lambda)\) обозначает усреднённую спектральную плотность до масштаба \(\Lambda\), \(\Delta_{\alpha}\) --- фазовое расщепление между глобальными ветвями, \(m_a\) --- эффективные массоподобные уровни внутренних каналов, а \(C_{a\alpha}\) --- смешанные коэффициенты отклика, зависящие одновременно от внутреннего индекса \(a\) и фазовой ветви \(\alpha\).

Этот набор выбран не случайно. В нём уже на уровне обозначений видно разделение ролей: \(\rho(\Lambda)\) естественно тяготеет к геометрическому каркасу, \(\Delta_{\alpha}\) --- к фазовому сектору, \(m_a\) --- к внутренней нагрузке, а коэффициенты \(C_{a\alpha}\) сразу претендуют на статус смешанных наблюдаемых.

\subsection{Шаг 1: выделение устойчивого ядра данных}

Первый вопрос состоит в том, какие элементы набора \(\mathcal{D}^{(1)}\) следует считать устойчивыми. Предположим, что сравнение различных режимов вычисления показывает следующее:
\begin{itemize}
    \item грубая асимптотика \(\rho(\Lambda)\) устойчива при малых локальных деформациях;
    \item величины \(\Delta_{\alpha}\) сохраняют качественный профиль при перестройке локальных параметров, но чувствительны к смене глобальной ветви;
    \item семейство \(m_a\) устойчиво как набор относительных расщеплений, хотя отдельные значения могут слегка дрейфовать;
    \item коэффициенты \(C_{a\alpha}\) частично нестабильны по значению, но сохраняют устойчивую схему зависимостей по индексам.
\end{itemize}

Тогда естественно определить устойчивое ядро как
\begin{equation}
\mathcal{D}^{(1)}_{\mathrm{stab}} = \{\rho_{as},\, [\Delta_{\alpha}],\, [m_a],\, [C_{a\alpha}]_{pat}\},
\end{equation}
где квадратные скобки обозначают не точные числа, а устойчивые классы поведения: профиль расщепления, относительную иерархию и паттерн смешанной зависимости.

Этот шаг крайне важен методологически. Он показывает, что новая программа стремится реконструировать не все численные детали сразу, а инвариантную организацию данных. Именно она и должна быть первым предметом теории.

\subsection{Шаг 2: первичная секторная атрибуция}

Теперь применим слабую атрибуцию \(\mathfrak{A}_0\) к элементам \(\mathcal{D}^{(1)}_{\mathrm{stab}}\). Получаем предварительную карту:
\begin{align}
\mathfrak{A}_0(\rho_{as}) &= \mathcal{B}, \\
\mathfrak{A}_0([\Delta_{\alpha}]) &= \Phi, \\
\mathfrak{A}_0([m_a]) &= \mathcal{L}_{int}, \\
\mathfrak{A}_0([C_{a\alpha}]_{pat}) &= \mathrm{mix}.
\end{align}

Смысл этой карты ясен. Ведущая спектральная плотность читается как след опорного каркаса. Фазовые расщепления прямо указывают на \(\Phi\). Внутренняя иерархия каналов наиболее естественно относится к \(\mathcal{L}_{int}\). Наконец, коэффициенты \(C_{a\alpha}\), зависящие и от внутреннего индекса, и от фазовой ветви, уже на первом взгляде не допускают чистой атрибуции.

Именно здесь теория делает первый нетривиальный шаг: она не довольствуется общим тезисом о многослойности, а строит рабочую карту того, какие следы какой слой преимущественно несёт.

\subsection{Шаг 3: реконструкция кандидатного инвариантного ядра}

После атрибуции нужно спросить, какая минимальная структура способна породить именно такую карту. В первом приближении кандидатное ядро \(\mathcal{I}^{*(1)}\) должно содержать три типа данных:
\begin{equation}
\mathcal{I}^{*(1)} = (G_{core},\, \Theta_{hol},\, \Sigma_{int}),
\end{equation}
где \(G_{core}\) кодирует геометрию опорного каркаса, \(\Theta_{hol}\) --- глобальный класс фазового переноса, а \(\Sigma_{int}\) --- внутреннюю канализацию и допустимые мультиплетные различия.

Тогда предполагается следующая схема порождения наблюдаемых:
\begin{align}
\rho_{as} &\leftarrow G_{core}, \\
[\Delta_{\alpha}] &\leftarrow \Theta_{hol}, \\
[m_a] &\leftarrow \Sigma_{int}, \\
[C_{a\alpha}]_{pat} &\leftarrow (\Theta_{hol},\Sigma_{int}).
\end{align}

Важно, что последний пункт сразу показывает необходимость смешанного уровня. Если некоторый наблюдаемый нельзя привязать к одному сектору без потери устойчивой структуры, теория обязана признать это прямо, а не насильно редуцировать его к одному источнику.

\subsection{Шаг 4: проверка деформационной устойчивости}

Теперь проверим, сохраняется ли реконструкция при допустимых слабых деформациях. Пусть \(\delta_g\) обозначает малую геометрическую деформацию, не меняющую глобальный фазовый класс, а \(\delta_{int}\) --- малую внутреннюю перестройку, не разрушающую опорный каркас.

Тогда от хорошей реконструкции требуется следующее:
\begin{align}
\delta_g &: \rho_{as}\ \text{может слегка изменяться, но остаётся в том же секторе } \mathcal{B}, \\
\delta_g &: [\Delta_{\alpha}]\ \text{не должны терять фазовую атрибуцию}, \\
\delta_{int} &: [m_a]\ \text{могут перераспределяться, но сохраняют принадлежность } \mathcal{L}_{int}, \\
\delta_{int} &: [C_{a\alpha}]_{pat}\ \text{сохраняют смешанный характер.}
\end{align}

Если именно так и происходит, то реконструкция проходит первый тест устойчивости. Это означает, что секторная карта не является случайной классификацией под один-единственный режим вычисления.

\subsection{Шаг 5: анализ смешанного наблюдаемого}

Рассмотрим теперь наиболее важный элемент --- коэффициенты \(C_{a\alpha}\). По самому своему виду они зависят от двух индексов, принадлежащих разным секторам. В рамках новой программы естественно предположить разложение:
\begin{equation}
C_{a\alpha} = C^{(0)} + U_a + V_{\alpha} + W_{a\alpha},
\end{equation}
где \(U_a\) кодирует внутренний вклад, \(V_{\alpha}\) --- фазовый, а \(W_{a\alpha}\) --- несводимое смешение.

Именно член \(W_{a\alpha}\) представляет наибольший интерес. Если он систематически ненулевой и устойчив как паттерн, это означает, что фазовый сектор и внутренняя нагрузка в действительности не независимы, а связаны в более глубоком источнике. Тем самым смешанное наблюдаемое становится прямым аргументом в пользу сущности \(\mathfrak{S}\), а не только в пользу сосуществования нескольких независимых описаний.

\subsection{Шаг 6: межъязыковая проверка}

Теперь необходимо проверить, не зависит ли вся реконструкция слишком сильно от спектрального языка. Для этого нужно спросить, может ли тот же класс данных быть прочитан геометрически и корреляционно.

В геометрическом окне \(\rho_{as}\) должна соответствовать эффективному каркасному масштабу, \([\Delta_{\alpha}]\) --- различию глобальных обходов, а \([m_a]\) --- внутренней секторной нагрузке, влияющей на подведущую структуру. В корреляционном окне тот же набор должен давать различие долгоживущих режимов связи, фазовую память и внутреннюю канализацию отклика. Если обе эти редукции сохраняют ту же инвариантную карту, то реконструкция проходит межъязыковой тест.

Именно такой тест отличает серьёзную исследовательскую программу от удачного частного формализма. Теория должна уметь удерживать идентичность сущности при переходе между языками, а не только внутри одного технического аппарата.

\subsection{Шаг 7: первое ограничивающее следствие}

Теперь можно спросить, даёт ли проведённая реконструкция хоть какое-то новое ограничение. Даёт. Если карта выше верна, то не всякое изменение фазового класса может происходить независимо от смешанных коэффициентов \(C_{a\alpha}\). Иначе говоря, при фиксированном опорном каркасе изменения \(\Theta_{hol}\) должны перестраивать не только \(\Delta_{\alpha}\), но и часть паттерна \(W_{a\alpha}\).

Это означает, что теория уже на первом цикле запрещает некоторый класс произвольных моделей, в которых фазовое расщепление и смешанная внутренняя структура меняются независимо. В этом и состоит первый зачаток предсказательной силы: не в точном числе, а в нетривиальном ограничении допустимого пространства моделей.

\subsection{Промежуточный итог цикла}

Таким образом, первый полный цикл реконструкции даёт следующую цепочку:
\begin{enumerate}
    \item данные редуцируются к устойчивому ядру;
    \item устойчивое ядро получает секторную карту;
    \item карта стягивается к минимальному кандидатному ядру \(\mathcal{I}^{*(1)}\);
    \item это ядро проходит проверку на допустимых деформациях;
    \item смешанные наблюдаемые обнаруживают внутреннюю связанность секторов;
    \item теория производит первое ограничение на допустимые интерпретации.
\end{enumerate}

Этот результат ещё не является доказательством истинности всей программы. Но он показывает, что введённые понятия --- \(\mathfrak{S}\), \(\mathcal{B}\), \(\Phi\), \(\mathcal{L}_{int}\), \(\mathcal{I}(\mathfrak{S})\) --- действительно могут работать совместно в одном исследовательском цикле, а не только украшать общий философский каркас.

\subsection{Что требуется дальше}

Следующий шаг теперь становится гораздо более определённым. Нужно перейти от пробного класса \(\mathcal{D}^{(1)}\) к более жёсткому классу данных, где хотя бы часть структур допускает уже не только качественную, но и полуколичественную реконструкцию. Это будет означать переход от первого демонстрационного цикла к началу собственно теоретической разработки.

Именно на этом рубеже трактат впервые получает шанс стать не просто программой о возможном устройстве мира, а формирующейся теорией о том, как глубинная сущность мира должна читаться через устойчивые структуры, инварианты и ограничения.


\section{Пределы применимости и критерии опровержения программы}

Всякая теория, претендующая на фундаментальность, должна уметь не только расширять поле интерпретации, но и ограничивать его. Если новая программа допускает любое наблюдаемое как совместимое с \(\mathfrak{S}\), то она не обладает объяснительной силой. Поэтому теперь необходимо зафиксировать не только условия её продуктивности, но и условия возможного провала. Иначе трактат рискует превратиться в язык, который умеет сопровождать любой результат, но не умеет ничему противостоять.

\subsection{Почему вопрос об опровержении принципиален}

До сих пор трактат двигался от интуиции единства к рабочим процедурам реконструкции. Однако именно на этом этапе возникает опасность, характерная для всякой глубокой объединяющей гипотезы: чем более общо её формулировать, тем труднее отличить силу от расплывчатости. Поэтому критерий опровержения здесь играет не вторичную, а конститутивную роль. Он определяет, является ли \(\mathfrak{S}\) реальным понятием теории или только метафорой для уже наблюдаемой согласованности.

Иначе говоря, вопрос следует поставить жёстко:
\begin{quote}
при каких типах результатов программа с сущностью \(\mathfrak{S}\), секторной декомпозицией и инвариантным ядром должна быть признана недостаточной или неверной?
\end{quote}

Только после такого вопроса теория получает право называться строгой исследовательской программой.

\subsection{Первый предел: отсутствие устойчивого ядра данных}

Наиболее фундаментальное условие провала касается самого начала реконструкции. Если для широкого класса существенных результатов не удаётся выделить устойчивое ядро \(\mathcal{D}_{\mathrm{stab}}\), то программа теряет почву. Это означало бы, что рассматриваемые наблюдаемые не образуют инвариантной организации, а распадаются на набор слабо связанных чисел и эффектов, для которых поиск единого источника оказывается насильственным.

Следовательно, первый критерий отрицания можно сформулировать так:
\begin{quote}
если существенный класс данных не допускает устойчивого инвариантного содержания при смене допустимых режимов описания, то реконструкция через \(\mathfrak{S}\) не запускается.
\end{quote}

Это не уничтожает автоматически саму идею глубинной структуры, но показывает, что именно в данной форме она неприменима к соответствующему материалу.

\subsection{Второй предел: неустойчивость секторной атрибуции}

Предположим, что устойчивое ядро данных выделить удалось. Тогда следующий возможный провал связан с картой атрибуции. Если малые допустимые деформации или переход между естественными окнами чтения радикально меняют секторную принадлежность одних и тех же наблюдаемых, то декомпозиция на \(\mathcal{B}\), \(\Phi\) и \(\mathcal{L}_{int}\) оказывается произвольной.

В таком случае не удаётся различить, что в наблюдаемом действительно принадлежит каркасу, что --- фазе, а что --- внутренней нагрузке. Тогда секторная структура перестаёт быть реальным аналитическим инструментом и вырождается в язык свободной переклассификации.

Поэтому второй критерий опровержения таков:
\begin{quote}
если для естественного класса наблюдаемых не существует устойчивой слабой атрибуции, сохраняющейся при допустимых деформациях и смене окна чтения, то секторная схема теории должна быть пересмотрена или отвергнута.
\end{quote}

\subsection{Третий предел: бесконечная недоопределённость инвариантного ядра}

Даже устойчивая карта атрибуции ещё не гарантирует успеха программы. Может оказаться, что одному и тому же набору наблюдаемых соответствуют бесконечно многие несовместимые инвариантные ядра, одинаково хорошо объясняющие данные. Тогда теория не извлекает структуру, а лишь переименовывает недоопределённость.

В этом случае \(\mathcal{I}(\mathfrak{S})\) перестаёт быть содержательным понятием. Оно уже не минимально и не реконструируемо; оно лишь служит местом, куда можно поместить любой недостающий механизм. Именно такой сценарий особенно опасен для фундаментальных программ, потому что внешне он похож на гибкость, а по сути означает отсутствие реальной дисциплины.

Отсюда следует третий критерий:
\begin{quote}
если реконструкция допускает бесконечное семейство структурно несовместимых ядер без естественного принципа минимизации, то теория в данной области не обладает объяснительной определённостью.
\end{quote}

\subsection{Четвёртый предел: распад межъязыковой согласованности}

До сих пор одной из сильнейших надежд программы было то, что разные языки --- геометрический, спектральный, корреляционный --- являются не конкурирующими теориями, а разными редукциями одной сущности. Но именно здесь лежит ещё один возможный провал. Если реконструированное ядро в одном языке систематически не переносится в другой без потери инвариантного содержания, то гипотеза о едином источнике ослабевает.

Здесь особенно важно различать технический перевод и инвариантное соответствие. Теория не обязана давать одинаковые формулы во всех языках. Но она обязана сохранять один и тот же структурный смысл. Если же геометрическое чтение требует одного ядра, а спектральное --- принципиально другого, причём никакая естественная редукция между ними не строится, то тезис об их общем происхождении оказывается под серьёзным сомнением.

Следовательно, четвёртый критерий таков:
\begin{quote}
если для устойчивого класса данных разные естественные окна чтения требуют несводимых и структурно несовместимых инвариантных ядер, то гипотеза о единой сущности \(\mathfrak{S}\) должна считаться по меньшей мере неполной.
\end{quote}

\subsection{Пятый предел: отсутствие новых ограничений}

Даже если все предыдущие шаги формально проходят, теория может всё ещё оставаться пустой. Это случается тогда, когда реконструкция красиво организует уже известный материал, но не даёт ни одного нового ограничения на допустимые интерпретации, ни одного исключаемого класса моделей, ни одного нетривиального соотношения между разными наблюдаемыми.

В таком случае программа может обладать высокой эвристической ценностью, но ещё не имеет статуса зрелой теории. Она объясняет лишь задним числом. Для фундаментального проекта этого недостаточно.

Поэтому пятый критерий необходимо сформулировать максимально прямо:
\begin{quote}
если из реконструкции \(\mathcal{I}(\mathfrak{S})\) не следуют никакие ограничения на ещё не проанализированные наблюдаемые, формы смешения или допустимые модели, то теория остаётся эвристикой, а не становится строгой программой объяснения.
\end{quote}

\subsection{Сильная и слабая формы опровержения}

Полезно различать два режима неуспеха программы.

\textbf{Слабое опровержение} означает, что теория не работает на некотором конкретном классе данных или в некотором конкретном окне чтения. Тогда речь может идти не о ложности всей программы, а о недостаточности текущей формализации, о слишком грубой секторной схеме или о неверно выбранной области применимости.

\textbf{Сильное опровержение} означало бы большее: что в широком спектре принципиально различных классов данных не обнаруживается ни устойчивого ядра, ни устойчивой атрибуции, ни межъязыковой согласованности, ни нетривиальных ограничений. В таком случае сама идея описывать глубинную структуру мира через \(\mathfrak{S}\) в её текущем виде должна быть оставлена.

Это различие важно, потому что защищает программу одновременно от двух крайностей: от догматической неуязвимости и от преждевременного отказа при первом локальном сбое.

\subsection{Граница между пределом и областью применимости}

Не всякий отрицательный результат является опровержением. Иногда он лишь показывает, что анализируемый объект не принадлежит той области, где программа обязана работать. Поэтому необходимо строго различать два случая:
\begin{enumerate}
    \item теория терпит неудачу там, где сама обещала реконструируемую структуру;
    \item теория сталкивается с объектом, который с самого начала лежит вне её предполагаемого режима применимости.
\end{enumerate}

Это различие особенно существенно для ранней стадии трактата. Пока формализация не доведена до окончательной точности, слишком широкий замах может создавать ложную видимость опровержения. Поэтому зрелая версия программы должна будет явно указывать, какие классы систем считаются допустимой ареной для сущности \(\mathfrak{S}\), а какие пока нет.

\subsection{Первый список проверяемых провалов}

На текущей стадии можно уже перечислить конкретные типы результатов, которые были бы особенно опасны для программы:
\begin{itemize}
    \item устойчивые данные, не допускающие никакой секторной карты;
    \item фазовые эффекты, полностью несвязанные с какими-либо смешанными или каркасными следами;
    \item внутренние мультиплетные структуры, которые нельзя согласовать ни с одной минимальной инвариантной организацией;
    \item систематическое расхождение между геометрическим и спектральным чтением одних и тех же устойчивых величин;
    \item успешные альтернативные описания, не требующие никакого аналога сущности \(\mathfrak{S}\), но дающие при этом более сильные ограничения.
\end{itemize}

Последний пункт особенно важен. Теория должна выдерживать конкуренцию не только с хаосом данных, но и с более экономными объяснениями. Если другой подход объясняет то же самое проще и с большей предсказательной силой, то новая программа обязана либо превзойти его, либо уступить ему место.

\subsection{Промежуточный вывод}

Тем самым трактат впервые фиксирует собственную уязвимость. Это важный шаг не меньше, чем введение новых понятий. До этого момента программа могла расширяться почти без сопротивления. Теперь она получает внутреннюю меру строгости: устойчивое ядро, атрибутивная стабильность, минимальность реконструкции, межъязыковая согласованность и наличие нетривиальных ограничений.

Если эти условия выполняются, теория укрепляется. Если они систематически нарушаются, она должна быть изменена или отвергнута. Именно так глубинная гипотеза начинает превращаться в дисциплинированную форму знания.

\subsection{Переход к следующему этапу}

После фиксации пределов и условий опровержения следующий шаг должен быть двояким. С одной стороны, нужно усилить положительную часть программы и перейти к более жёстким полуколичественным реконструкциям. С другой стороны, необходимо сознательно искать трудные для теории случаи --- не для разрушения текста, а для проверки его реальной несущей силы.

Только в таком двойном движении --- к большей конструктивности и к большей уязвимости --- трактат сможет претендовать на статус новой исследовательской программы, а не только на статус мощной философской интуиции.


\section{Полуколичественная схема реконструкции}

До сих пор трактат двигался в основном на уровне качественной структуры: выделялись инвариантные ядра, вводились секторные карты, строились протоколы реконструкции и формулировались критерии опровержения. Однако программа не сможет стать зрелой, если она не начнёт переходить к более жёстким отношениям между величинами. Этот переход ещё не означает окончательной математизации, но он требует следующего шага: ввести полуколичественную схему, где разные секторы сущности \(\mathfrak{S}\) проявляются уже не просто как смысловые роли, а как вклады, масштабы и относительные веса в наблюдаемой структуре.

\subsection{Зачем нужен полуколичественный уровень}

Между чистой философской онтологией и полной вычислительной теорией существует промежуточный слой. Именно на нём решается, обладает ли программа реальной несущей силой. Если теория не умеет хотя бы схематически упорядочить относительную величину вкладов, характер смешения и масштаб чувствительности наблюдаемых, то она остаётся качественной герменевтикой. Если же такие отношения удаётся построить, появляется возможность двигаться к более строгим формулам без потери исходного единства.

Поэтому полуколичественная схема нужна не как украшение и не как псевдоматематика. Её задача проста: ввести язык, в котором наблюдаемое уже можно сравнивать по степени каркасной, фазовой и внутренней обусловленности.

\subsection{Базовая идея весовой декомпозиции}

Пусть \(O\) --- наблюдаемая величина, полученная из некоторого окна чтения сущности \(\mathfrak{S}\). Тогда в первом полуколичественном приближении будем представлять её в виде
\begin{equation}
O \approx w_{\mathcal{B}}(O)\,O_{\mathcal{B}} + w_{\Phi}(O)\,O_{\Phi} + w_{int}(O)\,O_{int} + O_{\times},
\end{equation}
где коэффициенты \(w_{\mathcal{B}}(O)\), \(w_{\Phi}(O)\), \(w_{int}(O)\) задают относительный вес трёх секторных вкладов, а член \(O_{\times}\) отвечает за несводимое смешение.

Здесь особенно важно подчеркнуть: веса не обязаны быть численно точными на данной стадии. Они выступают как организующие параметры сравнения. Две наблюдаемые величины считаются близкими по типу, если их весовой профиль близок; два режима чтения считаются структурно различными, если хотя бы один из весов или смешанный остаток меняется качественно.

\subsection{Нормировка весов}

Чтобы весовая схема не превратилась в произвольный набор обозначений, необходимо ввести нормировку. В первом приближении удобно требовать
\begin{equation}
0 \leq w_{\mathcal{B}}, w_{\Phi}, w_{int} \leq 1,
\qquad
w_{\mathcal{B}} + w_{\Phi} + w_{int} + w_{\times} = 1,
\end{equation}
где \(w_{\times}\) есть эффективная мера смешанного остатка. В таком виде веса читаются как доли структурного влияния, а не как буквальные вероятности.

Такая нормировка нужна по двум причинам. Во-первых, она заставляет сравнивать разные наблюдаемые в одном масштабе. Во-вторых, она сразу делает видимым различие между почти чистыми и по-настоящему смешанными величинами. Если, например, \(w_{\Phi}\) велик, а \(w_{\times}\) мал, то перед нами почти чистое фазовое наблюдаемое. Если же заметен и фазовый, и внутренний вклад, и при этом остаток смешения не исчезает, то теория обязана рассматривать такое наблюдаемое как существенный тест глубинной структуры.

\subsection{Каркасные, фазовые и внутренние масштабы}

Теперь введём три характерных масштаба чувствительности:
\begin{equation}
\Lambda_{\mathcal{B}}, \qquad \Lambda_{\Phi}, \qquad \Lambda_{int}.
\end{equation}

Под \(\Lambda_{\mathcal{B}}\) будем понимать масштаб, на котором наблюдаемое начинает уверенно чувствовать опорный каркас. Под \(\Lambda_{\Phi}\) --- масштаб проявления глобальной фазовой ветви. Под \(\Lambda_{int}\) --- масштаб, на котором становится различима внутренняя канализация или мультиплетная структура.

Тогда уже на полуколичественном уровне появляется естественная классификация. Если
\begin{equation}
\Lambda_{\mathcal{B}} \ll \Lambda_{\Phi}, \Lambda_{int},
\end{equation}
то ведущая структура наблюдаемого, скорее всего, каркасна. Если глобальная фазовая чувствительность проявляется раньше внутренней, то можно ожидать доминирование \(\Phi\) над \(\mathcal{L}_{int}\) в соответствующем окне чтения. И наоборот, раннее проявление внутреннего сектора делает более вероятной мультиплетную или канализованную природу эффекта.

\subsection{Полуколичественное чтение спектрального класса}

Применим эту схему к уже введённому пробному набору \(\mathcal{D}^{(1)}\). Для ведущей спектральной плотности \(\rho_{as}\) естественно ожидать профиль вида
\begin{equation}
w_{\mathcal{B}}(\rho_{as}) \approx 1,
\qquad
w_{\Phi}(\rho_{as}) \ll 1,
\qquad
w_{int}(\rho_{as}) \ll 1.
\end{equation}

Для фазового расщепления \([\Delta_{\alpha}]\), напротив, разумно ожидать
\begin{equation}
w_{\Phi}([\Delta_{\alpha}]) \approx 1,
\qquad
w_{\mathcal{B}}([\Delta_{\alpha}]),\, w_{int}([\Delta_{\alpha}]) \ll 1,
\end{equation}
если только речь не идёт о сильно смешанном режиме. Для внутренней иерархии \([m_a]\) естественен профиль с доминированием \(w_{int}\). Наконец, для коэффициентов \([C_{a\alpha}]_{pat}\) следует ожидать принципиально иной режим:
\begin{equation}
0 < w_{\Phi}([C_{a\alpha}]_{pat}),\, w_{int}([C_{a\alpha}]_{pat}) < 1,
\qquad
w_{\times}([C_{a\alpha}]_{pat}) \not\approx 0.
\end{equation}

Тем самым качественная карта из предыдущего раздела приобретает полуколичественный вид. Теперь уже можно говорить не просто о принадлежности к сектору, а о степени доминирования и силе смешения.

\subsection{Матрица чувствительности}

Для систематизации удобно ввести матрицу чувствительности наблюдаемых:
\begin{equation}
\mathbf{S}(O) =
\begin{pmatrix}
\partial_{\mathcal{B}} O \\
\partial_{\Phi} O \\
\partial_{int} O
\end{pmatrix},
\end{equation}
где символы \(\partial_{\mathcal{B}}\), \(\partial_{\Phi}\), \(\partial_{int}\) обозначают не буквальные производные в аналитическом смысле, а нормированные меры чувствительности к малым допустимым изменениям соответствующих секторов.

Тогда веса можно мыслить как нормированные величины этих чувствительностей:
\begin{equation}
w_{\mathcal{B}}(O) \sim \frac{|\partial_{\mathcal{B}} O|}{|\partial_{\mathcal{B}} O| + |\partial_{\Phi} O| + |\partial_{int} O| + |\partial_{\times} O|},
\end{equation}
и аналогично для остальных компонент. Даже если на ранней стадии теория не умеет вычислять эти величины точно, сама форма такого определения задаёт дисциплину будущих вычислений.

\subsection{Первое полуколичественное ограничение}

Теперь можно сформулировать первое более жёсткое требование к программе. Если наблюдаемое классифицировано как почти чисто фазовое, то при малых каркасных деформациях должно выполняться соотношение вида
\begin{equation}
|\partial_{\mathcal{B}} O| \ll |\partial_{\Phi} O|.
\end{equation}

Если же наблюдаемое объявлено существенно смешанным, то теория, напротив, не имеет права занулять остаток смешения без потери объяснительной силы. Это означает, что для некоторых классов наблюдаемых должно выполняться
\begin{equation}
|\partial_{\times} O| \gtrsim \min\{ |\partial_{\Phi} O|, |\partial_{int} O| \}.
\end{equation}

Эти соотношения ещё не являются строгими законами. Но они уже задают полуколичественную геометрию допустимых объяснений. Они позволяют различать правдоподобные и неправдоподобные интерпретации внутри самой программы.

\subsection{Переход от весов к классам универсальности}

Следующий полезный шаг состоит в объединении наблюдаемых по близости их весовых профилей. Так возникают \textit{классы универсальности чтения}. Две величины могут быть совершенно разными по физическому смыслу, но принадлежать одному классу, если они имеют близкие соотношения между \(w_{\mathcal{B}}\), \(w_{\Phi}\), \(w_{int}\) и \(w_{\times}\).

Это даёт программе новое преимущество. Вместо того чтобы анализировать каждую наблюдаемую изолированно, теория может искать семейства объектов с общей структурной динамикой чувствительности. Именно такие семейства, вероятно, и окажутся естественной ареной для более строгой математизации \(\mathfrak{S}\).

\subsection{Пределы самой полуколичественной схемы}

Однако и этот уровень не следует переоценивать. Весовая декомпозиция полезна только при двух условиях. Во-первых, если действительно существует устойчивый способ сравнивать чувствительности. Во-вторых, если введение весов не скрывает под одной формой принципиально разные механизмы. Иначе схема снова рискует стать лишь более изящной версией прежней неопределённости.

Следовательно, полуколичественный уровень сам подлежит тем же критериям строгости, что и вся программа. Он должен быть устойчив, минимален, межъязыково согласован и давать новые ограничения, а не только новые обозначения.

\subsection{Промежуточный вывод}

Таким образом, полуколичественная схема реконструкции делает следующий важный шаг. Она переводит разговор о секторах из плоскости чистой атрибуции в плоскость относительных вкладов, масштабов чувствительности и структурных весов. Благодаря этому теория впервые получает возможность говорить не только \textit{что} доминирует в наблюдаемом, но и \textit{насколько} оно доминирует относительно других источников.

Именно в этом состоит её переходная сила. Она ещё не является полной математической теорией, но уже перестаёт быть исключительно качественной. Следующий шаг должен состоять в том, чтобы выбрать один класс наблюдаемых и попробовать построить для него первый устойчивый полуколичественный инвариант.

\subsection{Вычислимый пример первой весовой реконструкции}

Чтобы полуколичественная схема не оставалась только декларацией, полезно сразу рассмотреть пример, стилистически ближе к физическому языку трактата. Пусть \(M_{eff}\) обозначает некоторый \textit{эффективный отклик} наблюдаемого режима --- например, массообразный сдвиг или энергетическую цену локализованной моды --- а его значение определяется тремя типами вклада: каркасным фоном \(b\), фазовой деформацией \(\varphi\) и внутренней нагрузкой \(u\). В простейшей пробной схеме можно записать
\begin{equation}
M_{eff}(b,\varphi,u)=2b+5\varphi+3u+0.4\,\varphi u.
\end{equation}

Здесь первый член отвечает за жёсткость опорного каркаса, второй --- за чисто фазовый сдвиг, третий --- за внутреннюю структурную нагрузку, а билинейный член \(0.4\,\varphi u\) описывает их смешение. Тем самым модель уже не является просто абстрактной функцией трёх переменных: она читается как минимальная эффективная формула для отклика, чувствительного сразу к нескольким слоям глубинной организации.

Тогда секторные чувствительности в точке \((b,\varphi,u)=(1,2,1)\) равны
\begin{align}
\partial_{\mathcal{B}}M_{eff} &= \frac{\partial M_{eff}}{\partial b}=2, \\n\partial_{\Phi}M_{eff} &= \frac{\partial M_{eff}}{\partial \varphi}=5+0.4u=5.4, \\n\partial_{int}M_{eff} &= \frac{\partial M_{eff}}{\partial u}=3+0.4\varphi=3.8.
\end{align}
а смешанный вклад естественно оценить через величину фазово-внутреннего канала
\begin{equation}
\partial_{\times}M_{eff}\sim \left|0.4\,\varphi u\right|=0.8.
\end{equation}

После нормировки по сумме модулей
\begin{equation}
N=|2|+|5.4|+|3.8|+|0.8|=12,
\end{equation}
получаем первый явный весовой профиль:
\begin{equation}
w_{\mathcal{B}}=\frac{2}{12}\approx 0.17,
\qquad
w_{\Phi}=\frac{5.4}{12}=0.45,
\qquad
w_{int}=\frac{3.8}{12}\approx 0.32,
\qquad
w_{\times}=\frac{0.8}{12}\approx 0.07.
\end{equation}

Физический смысл такого профиля довольно прозрачен. Перед нами режим, в котором ведущую роль играет фазовый сектор, но существенная доля отклика всё ещё удерживается внутренней нагрузкой, тогда как каркасный фон остаётся подчинённым, хотя и не исчезающе малым. Смешанный канал невелик, однако ненулевой, а потому наблюдаемое нельзя считать чистой проекцией одного-единственного сектора.

Можно сделать и следующий минимальный шаг: сравнить этот же отклик с другим окном чтения, соответствующим ослаблению фазовой деформации. Если при тех же \(b\) и \(u\) фазовый параметр уменьшается от \(\varphi=2\) до \(\varphi=1\), то
\begin{equation}
\partial_{\Phi}M_{eff}=5.4 \mapsto 5.4,
\qquad
\partial_{int}M_{eff}=3.8 \mapsto 3.4,
\qquad
\partial_{\times}M_{eff}=0.8 \mapsto 0.4,
\end{equation}
и нормировка даёт уже иной профиль:
\begin{equation}
w_{\mathcal{B}}\approx 0.18,
\qquad
w_{\Phi}\approx 0.49,
\qquad
w_{int}\approx 0.31,
\qquad
w_{\times}\approx 0.04.
\end{equation}

Следовательно, ослабление фазовой деформации в этой пробной модели не уничтожает фазовое доминирование, но заметно подавляет именно канал фазово-внутреннего смешения. В языке эффективной физики это означало бы, что наблюдаемый отклик остаётся управляемым главным образом фазовым сектором, однако становится ``чище'' по составу и менее чувствительным к внутренней перестройке. Именно такого рода вычислимые переходы и нужны полуколичественной программе: они позволяют не просто называть сектор ведущим, а отслеживать, насколько устойчиво он сохраняет лидерство при контролируемой деформации окна чтения.


\section{Первый устойчивый полуколичественный инвариант}

Теперь можно сделать следующий, уже существенно более строгий шаг. Если полуколичественная схема действительно полезна, то она должна позволить построить не только распределение весов и чувствительностей, но и такую составную величину, которая остаётся сравнительно устойчивой при смене окна чтения и при допустимых слабых деформациях. Именно такую величину и естественно назвать первым устойчивым полуколичественным инвариантом новой программы.

\subsection{Что требуется от инварианта}

На этой стадии инвариант не обязан быть точным числом, одинаковым во всех моделях. Требование к нему скромнее, но всё же достаточно строгое. Он должен:
\begin{enumerate}
    \item строиться из уже введённых весов, чувствительностей и масштабов;
    \item сохранять свой структурный смысл при переходе между естественными окнами чтения;
    \item быть достаточно устойчивым к малым допустимым деформациям;
    \item различать по меньшей мере некоторые типы глубинной организации, которые на чисто качественном уровне выглядят слишком похожими.
\end{enumerate}

Иначе говоря, инвариант должен быть первой величиной, которая не просто описывает теорию, а позволяет ей \textit{сравнивать} разные реализации \(\mathfrak{S}\) на общем языке.

\subsection{Идея структурного контраста}

Естественная трудность состоит в том, что сами веса \(w_{\mathcal{B}}, w_{\Phi}, w_{int}, w_{\times}\) зависят от наблюдаемой величины. Поэтому искать инвариант отдельно для каждого коэффициента недостаточно. Нужна составная характеристика, которая отражает общий \textit{структурный контраст} между секторами.

В первом приближении введём для наблюдаемой величины \(O\) величину
\begin{equation}
\Xi(O) = \frac{w_{\Phi}(O) + w_{int}(O) + \eta\, w_{\times}(O)}{w_{\mathcal{B}}(O) + \varepsilon},
\end{equation}
где \(\eta\) есть коэффициент усиления смешанного вклада, а \(\varepsilon > 0\) --- малая регуляризующая величина, не дающая знаменателю обращаться в нуль.

Смысл \(\Xi(O)\) прост. Она измеряет, насколько не-каркасная часть наблюдаемого преобладает над каркасной. Если \(\Xi(O) \ll 1\), наблюдаемое почти полностью определяется опорным каркасом. Если \(\Xi(O) \gg 1\), перед нами величина, в которой фазовый, внутренний и смешанный режимы доминируют над геометрическим несущим фоном. Если же \(\Xi(O)\) имеет порядок единицы, это указывает на genuinely mixed regime, где разные источники соизмеримы.

\subsection{Почему именно такой тип величины полезен}

На первый взгляд выражение для \(\Xi(O)\) может показаться слишком грубым. Но именно грубость здесь и полезна. На ранней стадии теории важно не симулировать ложную точность, а найти устойчивую форму сравнения. Отношение типа \(\Xi(O)\) делает видимым то, что иначе остаётся распылённым по нескольким независимым коэффициентам: общий баланс между опорной структурой и всеми некаркасными механизмами вместе.

Это особенно важно для трактата, потому что одна из его центральных гипотез состоит именно в том, что мир не сводится к одной геометрической арене. Следовательно, первая полуколичественная величина и должна фиксировать степень отхода от чисто каркасного режима.

\subsection{Инвариант класса наблюдаемых}

Однако ещё более полезно перейти от отдельной наблюдаемой к целому классу \(\mathcal{C}\) наблюдаемых. Для такого класса определим усреднённую структурную контрастность:
\begin{equation}
\Xi[\mathcal{C}] = \left\langle \Xi(O) \right\rangle_{O \in \mathcal{C}}.
\end{equation}

Теперь \(\Xi[\mathcal{C}]\) можно читать как характеристику не отдельного эффекта, а целого семейства. Это уже гораздо ближе к инварианту теории. Если разные окна чтения --- спектральное, геометрическое, корреляционное --- приводят к близким значениям или хотя бы к близкому порядку \(\Xi[\mathcal{C}]\), то можно говорить, что данный класс наблюдаемых действительно несёт устойчивую структурную подпись сущности \(\mathfrak{S}\).

\subsection{Первое условие устойчивости}

Чтобы \(\Xi[\mathcal{C}]\) заслуживала называться инвариантом, необходимо сформулировать условие устойчивости. Пусть \(\delta \in \Delta_{adm}\) есть допустимая слабая деформация. Тогда требуется, чтобы для некоторого класса \(\mathcal{C}\) выполнялось соотношение вида
\begin{equation}
\big| \Xi[\mathcal{C}] - \Xi[\delta \mathcal{C}] \big| \leq \sigma_{\mathcal{C}},
\end{equation}
где \(\sigma_{\mathcal{C}}\) есть допустимая ширина структурного дрейфа для данного класса. Это соотношение не требует буквальной неизменности. Оно требует лишь, чтобы структурный контраст класса был устойчив в разумном диапазоне вариаций.

Именно это и отличает инвариант нового типа от жёсткой константы старого типа. Речь идёт не о точном числе, а об устойчивом режиме численного и структурного поведения.

\subsection{Применение к пробному классу данных}

Применим теперь эту схему к пробному классу \(\mathcal{D}^{(1)}\). Для ведущей спектральной плотности \(\rho_{as}\), как уже обсуждалось, ожидается почти каркасный профиль. Тогда
\begin{equation}
\Xi(\rho_{as}) \ll 1.
\end{equation}

Для фазового расщепления \([\Delta_{\alpha}]\), напротив, если каркасный вклад мал, получаем
\begin{equation}
\Xi([\Delta_{\alpha}]) \gg 1.
\end{equation}

Для коэффициентов смешения \([C_{a\alpha}]_{pat}\) наиболее интересен промежуточный случай:
\begin{equation}
\Xi([C_{a\alpha}]_{pat}) \sim \mathcal{O}(1)
\quad \text{или} \quad
\Xi([C_{a\alpha}]_{pat}) > 1
\end{equation}
в зависимости от того, насколько сильно каркас ещё удерживает соответствующее наблюдаемое.

Уже отсюда видно, что \(\Xi\) способна различать по меньшей мере три фундаментальных режима: каркасно-доминантный, некаркасно-доминантный и смешанный. Именно это делает её первой полезной полуколичественной характеристикой.

\subsection{Связь с матрицей чувствительности}

Инвариант \(\Xi\) можно выразить и через матрицу чувствительности. Если веса определяются нормированными чувствительностями, то в первом приближении имеем
\begin{equation}
\Xi(O) \sim \frac{|\partial_{\Phi} O| + |\partial_{int} O| + \eta |\partial_{\times} O|}{|\partial_{\mathcal{B}} O| + \varepsilon'}.
\end{equation}

Это представление особенно важно, потому что оно связывает инвариант не просто с формальной декомпозицией, а с реально наблюдаемой устойчивостью или неустойчивостью величины к разным классам деформаций. В дальнейшем именно такой язык может стать основой более строгих вычислительных определений.

\subsection{Критерий различимости классов}

Теперь можно сформулировать первое нетривиальное следствие. Если для двух классов наблюдаемых \(\mathcal{C}_1\) и \(\mathcal{C}_2\) выполняется
\begin{equation}
\big| \Xi[\mathcal{C}_1] - \Xi[\mathcal{C}_2] \big| \gg \max\{\sigma_{\mathcal{C}_1}, \sigma_{\mathcal{C}_2}\},
\end{equation}
то эти классы нельзя считать редукциями одного и того же режима чтения. Следовательно, инвариант \(\Xi\) уже позволяет различать не только отдельные наблюдаемые, но и семейства структур.

Это чрезвычайно важно для трактата. До сих пор теория в основном объединяла. Теперь она начинает и различать. А без способности различать никакая фундаментальная программа не может считаться полноценной.

\subsection{Почему инвариант не является окончательным}

Разумеется, \(\Xi\) ещё не есть окончательная фундаментальная величина. Она зависит от выбора параметра \(\eta\), от способа нормировки и от того, как именно определяются веса. Но именно поэтому она и называется первым устойчивым полуколичественным инвариантом, а не окончательным инвариантом теории. Её задача --- открыть сам класс таких величин и показать, что подобная конструкция вообще возможна внутри программы.

Если в дальнейшем удастся построить несколько независимых инвариантов такого типа и показать, что они согласованно ведут себя на разных классах данных, то теория сделает очень серьёзный шаг к зрелой формализации.

\subsection{Промежуточный вывод}

Тем самым новая программа впервые получает составную величину, которая делает сразу три вещи. Во-первых, она связывает качественные секторные атрибуции с полуколичественными весами. Во-вторых, она даёт критерий различия между каркасными, смешанными и некаркасными режимами. В-третьих, она создаёт язык для сравнения целых классов наблюдаемых на одном и том же структурном основании.

В этом смысле \(\Xi\) является не просто техническим приёмом, а первым знаком того, что теория начинает производить собственные объекты, а не только переосмыслять уже известные. Следующий шаг должен состоять в том, чтобы проверить, можно ли построить второй инвариант, чувствительный не к общему контрасту, а к внутренней асимметрии между фазовым и внутренним секторами.


\section{Второй полуколичественный инвариант: асимметрия фазового и внутреннего секторов}

Первый инвариант \(\Xi\) оказался полезен потому, что он измеряет общий структурный контраст между каркасным и некаркасным режимами. Но именно из-за этой общности он ещё не различает два принципиально разных случая: когда некаркасная часть наблюдаемого определяется главным образом фазовым сектором и когда она определяется главным образом внутренней нагрузкой. Для дальнейшего продвижения теории этого уже недостаточно. Следовательно, нужен второй инвариант, который не спрашивает, \textit{насколько} наблюдаемое удалено от чисто каркасного режима, а спрашивает, \textit{в какую сторону} внутри некаркасной части смещён его структурный баланс.

\subsection{Идея асимметрии}

Если \(\Xi\) измеряет вертикальный контраст между каркасом и всем остальным, то новый инвариант должен измерять горизонтальную асимметрию между \(\Phi\) и \(\mathcal{L}_{int}\). Именно такая величина позволит различать два класса явлений, которые при одинаково большом значении \(\Xi\) могут иметь совершенно разную внутреннюю природу.

В первом приближении для наблюдаемой величины \(O\) введём
\begin{equation}
\Upsilon(O) = \frac{w_{\Phi}(O) - w_{int}(O)}{w_{\Phi}(O) + w_{int}(O) + \zeta\, w_{\times}(O) + \varepsilon},
\end{equation}
где \(\zeta\) контролирует, насколько сильно смешанный вклад размывает чистую фазово-внутреннюю асимметрию, а \(\varepsilon > 0\) вновь служит регуляризацией.

Смысл \(\Upsilon(O)\) ясен сразу. Если \(\Upsilon(O) > 0\), то в некаркасной части наблюдаемого фазовый вклад преобладает над внутренним. Если \(\Upsilon(O) < 0\), то доминирует внутренняя нагрузка. Если же \(\Upsilon(O) \approx 0\), это означает либо почти симметричный баланс между фазой и внутренней структурой, либо столь сильное смешение, что прямое различение двух секторов уже невозможно без более тонкого анализа.

\subsection{Диапазон значений и режимы чтения}

При естественной нормировке величина \(\Upsilon(O)\) лежит приблизительно в диапазоне
\begin{equation}
-1 \lesssim \Upsilon(O) \lesssim 1.
\end{equation}

Тогда возникают три базовых режима:
\begin{itemize}
    \item \textbf{фазово-доминантный режим:} \(\Upsilon(O) \to 1\);
    \item \textbf{внутренне-доминантный режим:} \(\Upsilon(O) \to -1\);
    \item \textbf{асимметрично-смешанный режим:} \(\Upsilon(O) \approx 0\) при заметной некаркасной массе.
\end{itemize}

Это различие особенно важно, потому что два наблюдаемых могут иметь очень близкие значения \(\Xi\), то есть одинаково сильно уходить от чистого каркаса, но при этом один из них будет почти чисто фазовым, а другой --- почти чисто внутренним. Именно \(\Upsilon\) и должна схватывать эту разницу.

\subsection{Связь между \texorpdfstring{$\Xi$}{Xi} и \texorpdfstring{$\Upsilon$}{Upsilon}}

Теперь теория получает пару инвариантов. Первый измеряет силу некаркасности, второй --- её направление. В этом смысле \(\Xi\) и \(\Upsilon\) играют роли, отчасти похожие на модуль и ориентацию в двумерном пространстве некаркасных вкладов.

Если \(\Xi(O)\) мало, то значение \(\Upsilon(O)\) становится вторичным: наблюдаемое почти целиком каркасно, и различать фазовую и внутреннюю компоненты бессмысленно. Но если \(\Xi(O)\) велико, то \(\Upsilon(O)\) становится принципиальным. Тогда уже недостаточно знать, что наблюдаемое некаркасно; нужно знать, каким именно образом оно некаркасно.

Именно так теория начинает переходить от грубого различения режимов к более тонкой морфологии структурных классов.

\subsection{Инвариант класса наблюдаемых}

Как и в случае с \(\Xi\), полезно сразу перейти к уровню класса наблюдаемых. Для класса \(\mathcal{C}\) определим
\begin{equation}
\Upsilon[\mathcal{C}] = \left\langle \Upsilon(O) \right\rangle_{O \in \mathcal{C}}.
\end{equation}

Теперь \(\Upsilon[\mathcal{C}]\) читается как усреднённая фазово-внутренняя асимметрия класса. Если значение положительно и устойчиво, то данный класс следует считать преимущественно фазовым внутри своей некаркасной части. Если отрицательно --- преимущественно внутренним. Если близко к нулю, то либо класс реально сбалансирован, либо его смешанная природа слишком велика для грубого разложения.

\subsection{Применение к пробному классу данных}

Рассмотрим снова пробный набор \(\mathcal{D}^{(1)}\). Для фазового расщепления \([\Delta_{\alpha}]\) естественно ожидать
\begin{equation}
\Upsilon([\Delta_{\alpha}]) \approx 1.
\end{equation}

Для внутренней иерархии \([m_a]\), напротив,
\begin{equation}
\Upsilon([m_a]) \approx -1.
\end{equation}

Для смешанных коэффициентов \([C_{a\alpha}]_{pat}\) именно \(\Upsilon\) становится особенно содержательной величиной. Если
\begin{equation}
\Upsilon([C_{a\alpha}]_{pat}) > 0,
\end{equation}
то смешанный эффект следует понимать как внутреннюю структуру, ориентированную фазовым сектором. Если же
\begin{equation}
\Upsilon([C_{a\alpha}]_{pat}) < 0,
\end{equation}
то скорее наоборот: фазовый отклик оказывается встроен в более глубокую внутреннюю канализацию. Это уже не просто классификация, а реальная разница в интерпретации механизма.

\subsection{Чувствительность и асимметрия}

Инвариант \(\Upsilon\) можно выразить и через чувствительности:
\begin{equation}
\Upsilon(O) \sim
\frac{|\partial_{\Phi} O| - |\partial_{int} O|}{|\partial_{\Phi} O| + |\partial_{int} O| + \zeta |\partial_{\times} O| + \varepsilon'}.
\end{equation}

Такое представление особенно важно, потому что оно позволяет связывать асимметрию не с абстрактной словесной интерпретацией, а с измеримой реакцией наблюдаемого на разные типы допустимых вариаций. В будущем именно этот язык может привести к более строгому построению секторных метрик в пространстве данных.

\subsection{Первое новое различение классов}

Теперь уже можно сформулировать более тонкий критерий, чем раньше. Пусть два класса наблюдаемых \(\mathcal{C}_1\) и \(\mathcal{C}_2\) имеют близкие значения \(\Xi\), но существенно разные значения \(\Upsilon\). Тогда они могут быть одинаково некаркасными, но принадлежать разным структурным семействам. Формально это означает:
\begin{equation}
\Xi[\mathcal{C}_1] \approx \Xi[\mathcal{C}_2]
\quad \text{и одновременно} \quad
|\Upsilon[\mathcal{C}_1] - \Upsilon[\mathcal{C}_2]| \gg \tau,
\end{equation}
где \(\tau\) есть порог асимметрической различимости.

Это очень важный результат для трактата. Он показывает, что новая теория способна различать структуры второго порядка: не только насколько наблюдаемое отходит от геометрического каркаса, но и какой именно некаркасный механизм ведёт этот отход.

\subsection{Случай сильного смешения}

Особый интерес представляет режим, в котором и \(\Xi\) велико, и \(\Upsilon\) близко к нулю. Тогда наблюдаемое сильно некаркасно, но при этом не имеет выраженной фазовой или внутренней ориентации. В первом приближении это можно считать индикатором действительно сильного смешения, где фаза и внутренняя нагрузка уже нельзя отделить без более богатой структуры инвариантов.

Именно такой режим, вероятно, будет наиболее важен для будущей фундаментальной физики трактата. Если он устойчив и воспроизводим, то это значит, что существуют наблюдаемые, для которых сама разница между \(\Phi\) и \(\mathcal{L}_{int}\) является уже не первичной, а возникающей лишь на уровне вторичной редукции.

\subsection{Ограничения нового инварианта}

Как и \(\Xi\), величина \(\Upsilon\) не является окончательной. Она зависит от выбора схемы нормировки, от параметра \(\zeta\) и от того, насколько корректно выделяется смешанный вклад. Но её ценность состоит не в окончательности, а в том, что она открывает второй независимый канал различения структур. Если \(\Xi\) измеряет степень структурной некаркасности, то \(\Upsilon\) измеряет её внутреннюю ориентацию.

Тем самым теория получает уже не одну, а две взаимодополняющие координаты в пространстве своих наблюдаемых. Это означает, что понятийная геометрия трактата становится богаче и способна выдерживать более тонкие сравнения между классами данных.

\subsection{Промежуточный вывод}

Введение инварианта \(\Upsilon\) делает следующий важный шаг в развитии программы. Теперь теория умеет отвечать не только на вопрос, насколько наблюдаемое удалено от чисто каркасного режима, но и на вопрос, к какому из двух главных некаркасных секторов оно структурно склоняется. Благодаря этому смешанные явления начинают распадаться на более тонкие морфологические типы.

Следующий шаг должен состоять в том, чтобы объединить \(\Xi\) и \(\Upsilon\) в общую карту структурных режимов и показать, как на этой карте могут размещаться разные классы наблюдаемых. Именно это превратит набор отдельных инвариантов в первое геометрическое пространство состояний теории.


\section{Карта структурных режимов}

После введения инвариантов \(\Xi\) и \(\Upsilon\) теория впервые получает возможность перейти от отдельных характеристик к целостной геометрической картине. Если \(\Xi\) измеряет степень некаркасности, а \(\Upsilon\) --- внутреннюю ориентацию некаркасной части, то вместе они задают двумерное пространство режимов чтения. Именно в этом пространстве можно размещать наблюдаемые, классы наблюдаемых и даже целые семейства моделей. Тем самым трактат впервые получает не просто набор определений, а карту состояний собственной программы.

\subsection{Пространство режимов}

Будем рассматривать пару
\begin{equation}
\mathfrak{R}(O) = \big( \Xi(O), \, \Upsilon(O) \big)
\end{equation}
как координаты наблюдаемой величины \(O\) в пространстве структурных режимов. Аналогично для класса наблюдаемых \(\mathcal{C}\) положим
\begin{equation}
\mathfrak{R}[\mathcal{C}] = \big( \Xi[\mathcal{C}], \, \Upsilon[\mathcal{C}] \big).
\end{equation}

В таком виде каждая наблюдаемая или класс наблюдаемых получает точку на карте. Эта карта не является физическим пространством в обычном смысле. Это пространство редукций, то есть пространство того, \textit{как} глубинная сущность проявляется в наблюдаемом.

\subsection{Главные области карты}

Уже на первом приближении в пространстве \((\Xi,\Upsilon)\) естественно выделяются несколько областей.

\textbf{Каркасная область.} Если \(\Xi \ll 1\), то наблюдаемое лежит близко к оси каркасного доминирования. В этой области значение \(\Upsilon\) вторично, поскольку фазовый и внутренний секторы ещё слишком слабы, чтобы их различие играло определяющую роль.

\textbf{Фазово-некаркасная область.} Если \(\Xi \gg 1\) и одновременно \(\Upsilon \to 1\), то наблюдаемое существенно некаркасно и при этом явно ориентировано фазовым сектором.

\textbf{Внутренне-некаркасная область.} Если \(\Xi \gg 1\) и \(\Upsilon \to -1\), то наблюдаемое столь же далеко от каркасного режима, но его некаркасность определяется уже внутренней нагрузкой, а не глобальной фазой.

\textbf{Сильно-смешанная область.} Если \(\Xi \gg 1\) и \(\Upsilon \approx 0\), то речь идёт о режиме, где и фаза, и внутренняя структура существенны, но ни одна из них не получает устойчивого преобладания. Именно эта область наиболее интересна для гипотезы о глубинной сущности, потому что она указывает на режимы, где сами секторные различия могут быть лишь вторичными редукциями.

\textbf{Переходная область.} Если \(\Xi \sim \mathcal{O}(1)\), то наблюдаемое лежит между почти каркасным и явно некаркасным режимом. Здесь особенно важны траектории деформации, потому что именно они показывают, в какую сторону развивается структура при смене параметров чтения.

\subsection{Классификация пробного класса данных на карте}

Теперь можно разместить уже обсуждавшийся пробный набор \(\mathcal{D}^{(1)}\) на этой карте. Для ведущей спектральной плотности \(\rho_{as}\) естественно ожидать положение в каркасной области:
\begin{equation}
\mathfrak{R}(\rho_{as}) \in \mathcal{R}_{core}.
\end{equation}

Для фазового расщепления \([\Delta_{\alpha}]\) положение должно лежать в фазово-некаркасной области:
\begin{equation}
\mathfrak{R}([\Delta_{\alpha}]) \in \mathcal{R}_{\Phi}.
\end{equation}

Для внутренней иерархии \([m_a]\) --- во внутренне-некаркасной области:
\begin{equation}
\mathfrak{R}([m_a]) \in \mathcal{R}_{int}.
\end{equation}

Наконец, смешанные коэффициенты \([C_{a\alpha}]_{pat}\) должны располагаться либо в сильно-смешанной области, либо на границе между смешанной и одной из ориентированных областей:
\begin{equation}
\mathfrak{R}([C_{a\alpha}]_{pat}) \in \mathcal{R}_{mix} \cup \partial \mathcal{R}_{mix}.
\end{equation}

Тем самым карта режимов уже начинает работать как компактный синтез того, что в предыдущих разделах было распределено по нескольким логическим шагам.

\subsection{Границы областей как места структурного перехода}

Особенно важны не только сами области, но и их границы. На границе между \(\mathcal{R}_{core}\) и некаркасными областями теория должна ожидать чувствительные переходы, при которых малые изменения параметров чтения могут существенно перестроить относительную роль каркаса и некаркасной части. На границах между \(\mathcal{R}_{\Phi}\), \(\mathcal{R}_{int}\) и \(\mathcal{R}_{mix}\) должны возникать режимы, где один сектор начинает перехватывать доминирование у другого.

Эти границы особенно ценны, потому что именно на них теория может ожидать наиболее богатые и наиболее тонкие наблюдаемые эффекты. Далеко внутри области наблюдаемое уже устойчиво принадлежит некоторому типу. На границе же оно сохраняет память о нескольких типах сразу. Следовательно, именно границы являются естественными полигонами для поиска новых следствий теории.

\subsection{Траектории деформации}

Теперь возникает ещё одна важная идея. Допустимая деформация \(\delta \in \Delta_{adm}\) может быть прочитана как перемещение точки \(\mathfrak{R}(O)\) по карте структурных режимов. Тогда новая программа получает динамический смысл: она может описывать не только статическую принадлежность наблюдаемого к области, но и траекторию его структурной эволюции.

Схематически это можно записать как
\begin{equation}
\delta : \mathfrak{R}(O) \longmapsto \mathfrak{R}(\delta O).
\end{equation}

В таком виде сразу возникают новые исследовательские вопросы. Существуют ли устойчивые траектории, вдоль которых \(\Xi\) почти не меняется, а \(\Upsilon\) меняется существенно? Существуют ли семейства наблюдаемых, которые движутся к сильному смешению? Существуют ли запрещённые переходы между областями? Все эти вопросы уже не философские, а вполне структурно-теоретические.

\subsection{Метрика различия режимов}

Чтобы карта не оставалась чисто качественной диаграммой, полезно ввести простейшую метрику различия режимов. Для двух классов наблюдаемых \(\mathcal{C}_1\) и \(\mathcal{C}_2\) положим
\begin{equation}
 d_{\mathfrak{R}}(\mathcal{C}_1,\mathcal{C}_2)
 =
 \sqrt{
 \alpha\,\big(\Xi[\mathcal{C}_1]-\Xi[\mathcal{C}_2]\big)^2
 +
 \beta\,\big(\Upsilon[\mathcal{C}_1]-\Upsilon[\mathcal{C}_2]\big)^2
 },
\end{equation}
где \(\alpha\) и \(\beta\) задают относительную значимость вертикального и горизонтального различия.

Даже такая грубая метрика уже полезна. Она позволяет говорить не только о том, совпадают ли два класса по типу, но и насколько далеко они расположены друг от друга в пространстве режимов. Это, в свою очередь, создаёт основу для будущего понятия структурной близости между теориями, моделями и наблюдаемыми семействами.

\subsection{Первое геометрическое различение классов}

Теперь уже можно сформулировать важное следствие. Два класса наблюдаемых могут иметь близкие значения \(\Xi\), но различаться по \(\Upsilon\), и наоборот. Следовательно, полная структурная близость требует согласия по обеим координатам. Иными словами, принадлежность к одному и тому же типу чтения должна пониматься не одномерно, а геометрически.

Это важный момент для трактата. До сих пор инварианты вводились по отдельности. Теперь же теория впервые утверждает, что её объекты естественно организуются не на прямой, а на карте. Именно это и является первым зародышем собственного геометрического пространства состояний программы.

\subsection{Структурные фазы теории}

На основании карты \((\Xi,\Upsilon)\) можно уже говорить и о \textit{структурных фазах} теории. Под фазой здесь понимается не микроскопическая физическая фаза в обычном смысле, а устойчивый класс режимов чтения сущности \(\mathfrak{S}\). Тогда каркасная, фазовая, внутренняя и сильно-смешанная области становятся четырьмя первыми структурными фазами новой схемы.

Это особенно плодотворно потому, что позволяет объединить локальные наблюдаемые и глобальную архитектуру трактата. Теперь теория начинает описывать не только отдельные величины, но и целые области возможного поведения. А это уже шаг к полноценной морфологии глубинной сущности.

\subsection{Пределы карты}

Разумеется, и эта карта пока не окончательна. Она двумерна, тогда как сама сущность \(\mathfrak{S}\), вероятно, богаче. Кроме того, параметры \(\alpha\), \(\beta\), выбор нормировки и само определение областей ещё не зафиксированы строго. Но достоинство карты не в завершённости, а в том, что она впервые собирает ранее введённые понятия в одно пространство сравнения.

Именно это делает её принципиально важной. Теория получает не просто новые формулы, а новую форму видения материала.

\subsection{Промежуточный вывод}

Карта структурных режимов завершает первый большой цикл развития трактата. Начав с философской интуиции о единой сущности, программа прошла через онтологию, секторную декомпозицию, протокол реконструкции, пробный цикл, критерии опровержения, полуколичественные веса и два первых инварианта. Теперь все эти элементы впервые собраны в одном геометрическом пространстве.

Следующий шаг должен состоять в том, чтобы использовать эту карту не только для описания, но и для порождения новых гипотез: искать возможные траектории между фазами, выделять устойчивые фиксированные области и строить первый набросок динамики самой карты. Именно тогда трактат окончательно выйдет из режима статической программы и начнёт двигаться как зарождающаяся теория структурной динамики.


\section{Динамика карты структурных режимов}

Построение карты \((\Xi,\Upsilon)\) было необходимым, но ещё не достаточным шагом. Карта сама по себе даёт только геометрию возможных режимов. Чтобы трактат начал двигаться как теория, а не только как классификационная схема, необходимо понять, как точки на этой карте могут перемещаться. Иначе говоря, нужен первый набросок \textit{структурной динамики}: законов, по которым наблюдаемые, классы наблюдаемых и модели переходят из одной области режимов в другую при допустимых деформациях.

\subsection{От статической карты к потоку}

Пусть \(t\) обозначает не физическое время в прямом смысле, а параметр деформации, масштаба чтения или последовательности исследовательских редукций. Тогда каждой наблюдаемой величине \(O\) можно сопоставить траекторию
\begin{equation}
 t \longmapsto \mathfrak{R}_t(O) = \big( \Xi_t(O), \, \Upsilon_t(O) \big).
\end{equation}

В этой записи зафиксирована ключевая идея: теория изучает не просто положения точек, а их движение под действием допустимых преобразований. Это движение может отражать изменение масштаба, изменение окна чтения, слабую деформацию каркаса, перестройку фазовой ветви или внутреннего сектора. Разные типы деформаций должны индуцировать разные типы потоков на карте.

\subsection{Поле структурного потока}

В первом приближении можно ввести эффективное поле потока на карте режимов:
\begin{equation}
\mathbf{V}(\Xi,\Upsilon) = \left( \frac{d\Xi}{dt}, \, \frac{d\Upsilon}{dt} \right).
\end{equation}

Здесь компонента \(d\Xi/dt\) показывает, усиливается или ослабевает некаркасность режима, а компонента \(d\Upsilon/dt\) показывает, смещается ли внутренняя ориентация наблюдаемого в сторону фазового или внутреннего доминирования. Даже если на данной стадии эти производные не вычисляются явно, сама идея поля потока задаёт язык будущих уравнений структурной динамики.

Именно это и важно: теория начинает мыслить себя не только как геометрию режимов, но и как динамику между ними.

\subsection{Три элементарных типа движения}

Уже теперь можно различить по меньшей мере три элементарных типа траекторий.

\textbf{Вертикальные траектории.} Это деформации, при которых \(\Upsilon\) почти не меняется, а \(\Xi\) существенно растёт или убывает. Они соответствуют усилению или ослаблению общей некаркасности при сохранении внутренней ориентации режима.

\textbf{Горизонтальные траектории.} Это деформации, при которых \(\Xi\) остаётся примерно постоянной, а \(\Upsilon\) меняется существенно. Такие траектории особенно интересны: они означают переход между фазово-ориентированным и внутренне-ориентированным режимами при почти неизменной общей степени удаления от каркасного поведения.

\textbf{Диагональные траектории.} Это наиболее общие движения, при которых меняются обе координаты. Именно они, вероятно, и будут типичны для сложных физических систем, где каркас, фаза и внутренняя нагрузка перестраиваются совместно.

\subsection{Фиксированные точки и устойчивые области}

Теперь можно поставить вопрос о структурных аналогах фиксированных точек. Назовём режим \((\Xi_*,\Upsilon_*)\) \textit{фиксированным}, если при некотором классе допустимых деформаций выполняется
\begin{equation}
\mathbf{V}(\Xi_*,\Upsilon_*) = (0,0).
\end{equation}

Такой режим не обязательно является единственной точкой в буквальном смысле. Он может соответствовать устойчивой области или узкому слою карты, внутри которого структурный тип наблюдаемых сохраняется при широком наборе деформаций. В этом случае разумно говорить не просто о фиксированной точке, а о \textit{устойчивом режиме}.

Каркасная область, например, может содержать собственный устойчивый режим, если для больших классов наблюдаемых малые деформации не выводят их из \(\Xi \ll 1\). Аналогично сильно-смешанная область может оказаться собственной устойчивой фазой, а не только промежуточной зоной между чистыми режимами.

\subsection{Притягивающие и отталкивающие режимы}

Если на карту введён поток, то естественно различать и типы устойчивости. Режим назовём \textit{притягивающим}, если близкие траектории стремятся к нему при действии некоторого семейства деформаций. Режим назовём \textit{отталкивающим}, если малые отклонения уводят наблюдаемое прочь от него.

Это различие особенно важно для трактата, потому что позволяет по-новому прочитать вопрос о фундаментальности. Возможно, глубинные структуры мира соответствуют не любым точкам карты, а именно некоторым притягивающим областям. Тогда наблюдаемая устойчивость физических структур объясняется не случайностью, а тем, что структурная динамика сама ведёт большие классы состояний к определённым режимам чтения.

\subsection{Запрещённые переходы}

Не менее важна и обратная возможность: существование запрещённых или крайне неустойчивых переходов. Если теория верна, не всякая точка карты должна быть достижима из всякой другой посредством допустимых слабых деформаций. Например, может оказаться, что переход из глубоко каркасной области сразу в глубоко внутренне-некаркасную область невозможен без прохождения через промежуточную смешанную зону. Тогда карта приобретает не только геометрию и динамику, но и причинную архитектуру допустимых траекторий.

Схематически это можно выразить так:
\begin{equation}
\mathcal{R}_{core} \not\rightsquigarrow \mathcal{R}_{int}
\quad \text{без прохождения через} \quad
\mathcal{R}_{mix} \text{ или } \mathcal{R}_{trans}.
\end{equation}

Подобные ограничения особенно ценны, потому что они уже имеют характер структурных предсказаний. Теория начинает не только описывать то, что возможно, но и говорить, что именно не должно происходить напрямую.

\subsection{Динамика пробного класса данных}

Попробуем теперь прочитать пробный класс \(\mathcal{D}^{(1)}\) динамически. Ведущая спектральная плотность \(\rho_{as}\), по-видимому, должна лежать вблизи каркасного устойчивого режима. Фазовое расщепление \([\Delta_{\alpha}]\) --- в фазово-некаркасной области, а внутренняя иерархия \([m_a]\) --- во внутренне-некаркасной. Особый интерес вновь представляют коэффициенты \([C_{a\alpha}]_{pat}\). Именно для них наиболее естественно ожидать нетривиальные траектории: при изменении окна чтения они могут смещаться вдоль смешанной области, а при изменении фазового или внутреннего доминирования --- переходить к одной из ориентированных некаркасных областей.

Иначе говоря, пробный класс данных теперь читается не как набор фиксированных точек, а как набор потенциальных орбит. Тем самым трактат делает ещё один шаг от классификации к динамике.

\subsection{Эффективные уравнения структурной эволюции}

Хотя полных уравнений у теории пока нет, уже можно записать их схему. Для класса \(\mathcal{C}\) положим:
\begin{align}
\frac{d\Xi[\mathcal{C}]}{dt} &= F_{\Xi}\big(\Xi[\mathcal{C}],\Upsilon[\mathcal{C}],\mu\big), \\
\frac{d\Upsilon[\mathcal{C}]}{dt} &= F_{\Upsilon}\big(\Xi[\mathcal{C}],\Upsilon[\mathcal{C}],\mu\big),
\end{align}
где \(\mu\) обозначает совокупность параметров деформации, масштаба и режима чтения. Эти уравнения пока формальны. Но они чрезвычайно важны как исследовательская рамка: теория впервые допускает возможность собственного динамического закона на пространстве режимов.

\subsection{Первый критерий динамической согласованности}

Теперь можно сформулировать новое требование к программе. Хорошая теория должна не только правильно размещать наблюдаемые на карте, но и правильно описывать допустимые направления их движения. Если одна и та же система в разных окнах чтения требует принципиально несовместимых траекторий, то карта режимов ещё недостаточно точна или сами инварианты определены грубо.

Следовательно, динамическая согласованность требует, чтобы разные редукции одной и той же сущности индуцировали не произвольные, а согласованные траектории на карте хотя бы на уровне их общей ориентации и устойчивых предельных областей.

\subsection{Промежуточный вывод}

Динамика карты структурных режимов делает решающий шаг вперёд. Теория больше не ограничивается описанием того, где находятся её объекты. Она начинает описывать, как они движутся, к каким областям тяготеют, каких переходов избегают и какие режимы сохраняют свою форму под действием деформаций. Тем самым трактат впервые получает черты подлинной структурной динамики.

Следующий шаг должен состоять в том, чтобы выделить на карте первые кандидаты на фиксированные режимы и предложить критерии их устойчивости. Именно это даст теории зачаток собственного аналога фазового портрета.


\section{Фиксированные режимы и их устойчивость}

После введения динамики карты \((\Xi,\Upsilon)\) следующий шаг действительно состоит в том, чтобы выделить на ней первые кандидаты на фиксированные режимы. Пока теория не умеет различать такие режимы, её динамика остаётся слишком общей: она говорит о возможности потока, но не указывает, какие конфигурации этот поток сохраняет, к каким он тяготеет и какие делает неустойчивыми. Между тем именно фиксированные режимы и образуют костяк всякого фазового портрета. Они задают те структурные формы, вокруг которых затем организуется всё поле допустимых траекторий.

\subsection{Определение фиксированного режима}

Напомним, что структурный режим описывается точкой \((\Xi,\Upsilon)\) на карте. Будем называть режим \((\Xi_*,\Upsilon_*)\) фиксированным, если в некотором допустимом семействе деформаций выполняется
\begin{equation}
F_{\Xi}(\Xi_*,\Upsilon_*,\mu_*) = 0,
\qquad
F_{\Upsilon}(\Xi_*,\Upsilon_*,\mu_*) = 0,
\end{equation}
где \(\mu_*\) обозначает соответствующий режим чтения или набор параметров, при котором структура находится в равновесии.

Уже здесь важно сделать одно уточнение. В ранней версии теории фиксированный режим лучше понимать не как идеально точечный объект, а как устойчивый локальный класс, вокруг которого поток замедляется и теряет направленную тенденцию. Иначе говоря, речь может идти о фиксированной точке, фиксированном слое или малой фиксированной области.

\subsection{Четыре естественных кандидата}

На текущей стадии трактат естественно подводит к четырём первым кандидатам на фиксированные режимы.

\textbf{Каркасный фиксированный режим} \(P_{core}\):
\begin{equation}
\Xi_* \ll 1,
\qquad
\Upsilon_* \text{ несущественно.}
\end{equation}
Этот режим соответствует ситуациям, где опорный каркас настолько доминирует, что фазовая и внутренняя ориентации не успевают стать главными носителями наблюдаемого смысла.

\textbf{Фазовый фиксированный режим} \(P_{\Phi}\):
\begin{equation}
\Xi_* \gg 1,
\qquad
\Upsilon_* \approx 1.
\end{equation}
Здесь теория ожидает устойчивые классы наблюдаемых, где некаркасная часть почти целиком ориентирована глобальной фазовой ветвью.

\textbf{Внутренний фиксированный режим} \(P_{int}\):
\begin{equation}
\Xi_* \gg 1,
\qquad
\Upsilon_* \approx -1.
\end{equation}
Это режим, в котором доминирует внутренняя канализация, а фазовая структура играет вторичную роль.

\textbf{Смешанный фиксированный режим} \(P_{mix}\):
\begin{equation}
\Xi_* \gg 1,
\qquad
\Upsilon_* \approx 0.
\end{equation}
Он особенно важен, потому что именно здесь теория получает шанс описывать наблюдаемые, для которых различие между фазовым и внутренним секторами само становится только эффективной вторичной редукцией.

\subsection{Линейная устойчивость}

Чтобы различать просто фиксированные режимы и устойчивые фиксированные режимы, нужно рассмотреть поведение потока в их окрестности. Для этого естественно ввести линейное приближение поля потока:
\begin{equation}
\begin{pmatrix}
\delta \dot{\Xi} \\
\delta \dot{\Upsilon}
\end{pmatrix}
=
J_*
\begin{pmatrix}
\delta \Xi \\
\delta \Upsilon
\end{pmatrix},
\end{equation}
где \(J_*\) есть матрица Якоби поля \((F_{\Xi},F_{\Upsilon})\), вычисленная в фиксированном режиме.

Хотя сама матрица пока не задана явно, структура теории уже позволяет сформулировать критерий: если собственные значения \(J_*\) имеют отрицательные вещественные части, режим следует считать локально притягивающим; если хотя бы одно из них имеет положительную вещественную часть, режим теряет устойчивость по соответствующему направлению.

Именно здесь трактат впервые получает настоящий аналог фазового портрета в математическом смысле, пусть пока и только схематический.

\subsection{Физический смысл локальной устойчивости}

Локальная устойчивость важна не сама по себе, а потому, что она связывает геометрию карты с физическим смыслом. Если, скажем, каркасный режим локально устойчив, это означает, что множество наблюдаемых при малых допустимых деформациях возвращаются к каркасной доминанте. Если же устойчив смешанный режим, это гораздо интереснее: тогда теория предсказывает существование систем, для которых совместная роль фазы и внутренней структуры не является переходной аномалией, а представляет собой устойчивую форму глубинной организации.

Именно поэтому разные фиксированные режимы должны интерпретироваться не просто как математические точки, а как кандидаты на фундаментальные типы структурной реальности.

\subsection{Бассейны притяжения}

Следующий важный шаг состоит в введении бассейнов притяжения. Для устойчивого режима \(P\) обозначим через \(\mathcal{B}(P)\) множество тех начальных конфигураций карты, траектории которых при действии допустимого потока стремятся к \(P\):
\begin{equation}
\mathcal{B}(P) = \{ \mathfrak{R}_0 : \mathfrak{R}_t(\mathfrak{R}_0) \to P \text{ при } t \to +\infty \}.
\end{equation}

Смысл этой конструкции огромен. Теория начинает различать не только устойчивые режимы, но и то, какие семейства наблюдаемых к ним тяготеют. Тогда вопрос о фундаментальности переносится с уровня отдельных точек на уровень геометрии целых областей. Возможно, именно большие бассейны притяжения и объясняют, почему некоторые структурные режимы встречаются в природе чаще и устойчивее других.

\subsection{Границы бассейнов как зоны высокой чувствительности}

Не менее важны границы между бассейнами. Вблизи таких границ малое изменение начальных условий или режима чтения может вести к качественно разным предельным режимам. Следовательно, именно здесь теория должна ожидать максимальную чувствительность, структурные бифуркации и наиболее богатые переходные явления.

Это хорошо согласуется с уже высказанной ранее мыслью о важности границ областей на карте режимов. Теперь эта мысль получает динамическое уточнение: границы между областями могут быть не просто геометрическими рубежами, но и сепаратрисами, разделяющими разные бассейны притяжения.

\subsection{Пробный фазовый портрет}

На текущей стадии можно предложить первый пробный фазовый портрет теории. Каркасный режим \(P_{core}\) естественно рассматривать как кандидат на сильный аттрактор для тех наблюдаемых, которые прежде всего читают ведущий спектральный или геометрический масштаб. Фазовый режим \(P_{\Phi}\) и внутренний режим \(P_{int}\) могут играть роль более специализированных аттракторов для классов данных, чувствительных соответственно к глобальной ветви и к внутренней канализации. Смешанный режим \(P_{mix}\) может оказаться либо седловым, либо, наоборот, особенно интересным локальным аттрактором для genuinely hybrid observables.

Последняя возможность особенно значима. Если смешанный режим устойчив, то сама теория получает аргумент в пользу того, что глубинная сущность мира по своей природе не обязана разлагаться на почти независимые секторные чтения. Напротив, смешение может быть фундаментальной нормой, а не исключением.

\subsection{Бифуркации структуры}

Теперь естественно ввести и понятие структурной бифуркации. Пусть при изменении параметра \(\mu\) фиксированный режим теряет устойчивость или распадается на несколько режимов. Тогда будем говорить, что карта переживает бифуркацию:
\begin{equation}
P(\mu < \mu_c) \longrightarrow \{P_1,P_2,\dots\}(\mu > \mu_c).
\end{equation}

Это особенно важно для трактата, потому что позволяет по-новому понимать возникновение разных слоёв реальности. Возможно, геометрический, фазовый и внутренний режимы не просто сосуществуют, а порождаются как разные ветви единой структурной бифуркации глубинной сущности. Тогда множественность режимов мира была бы не исходным данным, а следствием динамики самой \(\mathfrak{S}\).

\subsection{Первый критерий устойчивости программы}

На этом этапе можно сформулировать и новое проверочное требование к теории. Если карта режимов и её динамика осмысленны, то разумно ожидать, что важнейшие классы наблюдаемых будут группироваться вблизи немногих устойчивых режимов, а не заполнять карту хаотически и без структуры. Иначе вся идея фазового портрета потеряет смысл.

Следовательно, программа делает содержательное предсказание: пространство реально значимых наблюдаемых должно обладать не произвольной, а кластеризованной геометрией на карте \((\Xi,\Upsilon)\). Именно это может стать одним из сильнейших косвенных аргументов в пользу глубинной сущности \(\mathfrak{S}\).

\subsection{Промежуточный вывод}

Раздел о фиксированных режимах и их устойчивости завершает ещё один важный переход. Теперь теория умеет говорить не только о потоках и траекториях, но и о том, какие режимы обладают собственной устойчивой идентичностью, какие области к ним тяготеют и где проходят чувствительные границы между типами поведения. Иначе говоря, карта структурных режимов впервые превращается в настоящий фазовый портрет.

Следующий шаг должен состоять в том, чтобы сделать ещё один ход от геометрии к содержанию: предложить, какие именно наблюдаемые величины и какие классы данных в первую очередь должны быть сопоставлены с каждым из устойчивых режимов. Это будет означать переход от абстрактной динамики к более предметной архитектуре будущей теории.


\section{Сопоставление классов наблюдаемых с устойчивыми режимами}

После того как карта структурных режимов получила динамику и первый фазовый портрет, становится возможным следующий необходимый шаг: наполнить эту карту предметным содержанием. Иначе фазовый портрет останется слишком абстрактным. Теория должна уметь сказать не только, что существуют каркасный, фазовый, внутренний и смешанный режимы, но и какие именно типы наблюдаемых следует в первую очередь относить к каждому из них. Только так геометрия карты превращается в рабочий инструмент интерпретации материала.

\subsection{Принцип сопоставления}

Будем считать, что класс наблюдаемых \(\mathcal{C}\) сопоставляется устойчивому режиму \(P\), если выполняются три условия:
\begin{enumerate}
    \item центр массы класса \(\mathfrak{R}[\mathcal{C}]\) лежит вблизи \(P\);
    \item разброс класса мал по сравнению с расстоянием до других устойчивых режимов;
    \item допустимые деформации класса не выводят его из бассейна притяжения \(\mathcal{B}(P)\).
\end{enumerate}

Эти три условия важны потому, что они превращают язык режимов в язык устойчивых ассоциаций, а не случайных совпадений. Если класс наблюдаемых систематически возвращается к одному и тому же типу режима под разными окнами чтения и малыми деформациями, то теория имеет право считать эту ассоциацию содержательной.

\subsection{Каркасный режим и ведущие инварианты}

Наиболее естественным содержанием каркасного режима \(P_{core}\) являются такие классы наблюдаемых, которые прежде всего читают опорную геометрию и её ведущую спектральную организацию. Сюда в первом приближении должны относиться:
\begin{itemize}
    \item ведущие спектральные плотности и асимптотические коэффициенты;
    \item крупномасштабные геометрические инварианты;
    \item устойчивые масштабные характеристики, слабо чувствительные к внутренней нагрузке и смене фазовой ветви;
    \item наблюдаемые, для которых \(\Xi \ll 1\) и карта чувствительности почти полностью каркасна.
\end{itemize}

Если теория верна, именно эти величины должны образовывать наиболее широкий и устойчивый бассейн. В некотором смысле это естественно: каркасный режим отвечает за те аспекты мира, которые читаются наиболее грубо и наиболее универсально.

\subsection{Фазовый режим и глобально-чувствительные наблюдаемые}

Фазовый режим \(P_{\Phi}\) должен притягивать классы наблюдаемых, чья устойчивость определяется не локальным каркасом, а глобальной ветвью переноса, голономией, обходом циклов или дискретной фазовой структурой. В первую очередь сюда следует отнести:
\begin{itemize}
    \item фазовые расщепления между ветвями;
    \item голономные различия и ветви Вильсона;
    \item глобальные различия, чувствительные к спин-структуре;
    \item такие коэффициенты отклика, у которых \(\Xi\) велико, а \(\Upsilon\) устойчиво положительно.
\end{itemize}

Содержательно это означает: фазовый режим собирает те наблюдаемые, которые нельзя восстановить одной лишь локальной геометрией. Именно через него сущность \(\mathfrak{S}\) проявляет свою глобальную память.

\subsection{Внутренний режим и канализованные структуры}

Внутреннему режиму \(P_{int}\) должны соответствовать классы наблюдаемых, где главную роль играет внутренняя нагрузка, мультиплетность, субведущая операторная организация и канализация внутренних степеней свободы. В первом приближении сюда относятся:
\begin{itemize}
    \item внутренние иерархии уровней и мультиплетные расщепления;
    \item субведущие спектральные структуры, устойчиво зависящие от внутреннего содержания;
    \item параметры внутренней секторной дифференциации;
    \item классы наблюдаемых, для которых \(\Xi\) велико, а \(\Upsilon\) устойчиво отрицательно.
\end{itemize}

Этот режим особенно важен для трактата, потому что он показывает: внутренняя организация мира не обязана быть вторичной надстройкой над геометрией, а может являться равноправным способом чтения глубинной сущности.

\subsection{Смешанный режим и гибридные наблюдаемые}

Наиболее богатым по содержанию остаётся смешанный режим \(P_{mix}\). Именно сюда должны стекаться такие классы наблюдаемых, для которых каркасная редукция недостаточна, но и чистое фазовое или чистое внутреннее чтение также не объясняет устойчивый паттерн. В первую очередь это:
\begin{itemize}
    \item коэффициенты смешанного отклика типа \([C_{a\alpha}]_{pat}\);
    \item наблюдаемые, одновременно чувствительные к глобальной фазе и внутреннему каналу;
    \item корреляционные структуры, неразложимые на чисто геометрический и чисто внутренний вклад;
    \item классы, для которых \(\Xi\) велико, а \(\Upsilon\) близко к нулю при устойчиво ненулевом смешанном остатке.
\end{itemize}

Именно этот режим наиболее ценен для всей программы. Если он действительно существует как устойчивый аттрактор, то теория получает сильнейший аргумент в пользу того, что глубинная сущность мира по своей природе не обязана сводиться к сумме почти независимых чтений.

\subsection{Переходные классы наблюдаемых}

Нужно, однако, учитывать и такие классы наблюдаемых, которые не принадлежат устойчиво ни одному режиму, а локализуются в переходных областях между ними. Эти наблюдаемые особенно важны для ранней стадии теории, потому что именно они могут сигнализировать о бифуркациях, смене доминирования и переходах между бассейнами притяжения.

Переходные классы не следует считать теоретическим браком. Напротив, они могут оказаться главными индикаторами структурной динамики. Если устойчивые режимы задают скелет карты, то переходные классы показывают, как именно этот скелет связан в живую систему.

\subsection{Пробное распределение уже введённого материала}

На текущей стадии трактат может уже предложить первый пробный словарь сопоставлений:
\begin{align}
\rho_{as} &\rightsquigarrow P_{core}, \\
[\Delta_{\alpha}] &\rightsquigarrow P_{\Phi}, \\
[m_a] &\rightsquigarrow P_{int}, \\
[C_{a\alpha}]_{pat} &\rightsquigarrow P_{mix} \text{ или } \partial \mathcal{B}(P_{mix}).
\end{align}

Этот словарь пока ещё не обладает статусом окончательной классификации. Но он уже задаёт рабочую архитектуру будущей теории. С его помощью можно перечитывать накопленные результаты и проверять, не образуют ли они естественные кластеры вокруг указанных режимов.

\subsection{Первый критерий содержательной наполненности карты}

Теперь можно сформулировать важный новый критерий. Карта режимов считается содержательно наполненной тогда и только тогда, когда значимые классы наблюдаемых не просто распределяются по ней, а распределяются с заметной концентрацией вокруг устойчивых режимов и с осмысленными переходными слоями между ними.

Если же каждый новый класс наблюдаемых приходится вводить как исключение, не связанное ни с одним режимом и ни с одной границей бассейна, то сама карта ещё недостаточно адекватна. Следовательно, фазовый портрет теории проверяется не только своей внутренней красивостью, но и тем, насколько естественно он собирает реальные семейства наблюдаемых.

\subsection{От сопоставления к архитектуре теории}

Когда сопоставление становится устойчивым, карта перестаёт быть просто диаграммой. Она превращается в архитектуру теории. Каркасный режим начинает отвечать за универсальные ведущие структуры. Фазовый режим --- за глобальную память и ветвление. Внутренний режим --- за канализацию и мультиплетность. Смешанный режим --- за genuinely hybrid phenomena, в которых глубинная сущность говорит сразу несколькими голосами.

Именно в этой точке трактат впервые получает возможность перерасти из исследовательской программы в зарождающуюся систематику нового типа. Он уже не только предлагает понятия, но и раскладывает материал по собственной внутренней топологии.

\subsection{Промежуточный вывод}

Сопоставление классов наблюдаемых с устойчивыми режимами завершает ещё один принципиальный шаг. Фазовый портрет перестаёт быть абстрактной геометрией и становится рабочей картой интерпретации. Теперь теория умеет не только говорить о потоках и фиксированных областях, но и помещать в них реальные типы данных.

Следующий шаг должен состоять в том, чтобы проверить, можно ли на основе этой карты сформулировать первые строгие критерии отбора между конкурирующими моделями одной и той же области наблюдаемых. Именно это даст теории настоящий механизм внутренней селекции.


\section{Критерии отбора между конкурирующими моделями}

До этого момента трактат в основном строил пространство понятий, инвариантов, режимов и их динамики. Но зрелая теория должна делать ещё один шаг: она должна уметь предпочитать одни модели другим. Если две или более модели одинаково хорошо описывают один и тот же класс наблюдаемых, новая программа обязана предложить критерии, по которым между ними можно производить внутренний отбор. Иначе карта режимов останется лишь богатым языком интерпретации, но не станет инструментом теоретической селекции.

\subsection{Что значит конкуренция моделей в рамках трактата}

Под конкурирующими моделями здесь будем понимать не обязательно полностью разные физические теории. Достаточно того, что существует несколько структурно несовпадающих способов объяснить один и тот же класс наблюдаемых \(\mathcal{C}\): одна модель может делать акцент на каркасной геометрии, другая --- на фазовой ветви, третья --- на внутренней канализации, четвёртая --- на смешанном механизме.

Следовательно, задача отбора состоит не просто в сравнении численных подгонок. Она состоит в сравнении того, насколько естественно каждая модель встраивается в уже построенную архитектуру \(\mathfrak{S}\), \((\Xi,\Upsilon)\), устойчивых режимов и их динамики.

\subsection{Первый критерий: режимная совместимость}

Наиболее естественный критерий отбора --- это совместимость модели с фазовым портретом. Если класс наблюдаемых \(\mathcal{C}\) устойчиво локализуется вблизи режима \(P\), то предпочтение должна получать та модель, которая воспроизводит именно этот режим без искусственного переназначения вкладов.

Формально это можно мыслить так: пусть модель \(M\) индуцирует карту
\begin{equation}
\mathfrak{R}_M[\mathcal{C}] = (\Xi_M[\mathcal{C}],\Upsilon_M[\mathcal{C}]).
\end{equation}

Тогда модель считается режимно совместимой, если точка \(\mathfrak{R}_M[\mathcal{C}]\) лежит в той же устойчивой области, что и реконструированное по данным положение класса. Если же модель требует перевода класса в принципиально другой режим, то она должна считаться менее естественной, даже если локально воспроизводит часть численных характеристик.

\subsection{Второй критерий: структурная экономия}

Однако одной режимной совместимости недостаточно. Две модели могут помещать класс данных в один и тот же режим, но делать это с разной степенью искусственности. Поэтому второй критерий --- экономия структуры. Предпочтение получает та модель, которая объясняет наблюдаемое меньшим числом независимых допущений, меньшим числом специальных поправок и меньшим числом отдельно вводимых механизмов.

В языке трактата это означает: лучшая модель --- та, которая минимизирует дополнительные компоненты сверх необходимого инвариантного ядра. Если для воспроизведения одного и того же класса наблюдаемых модель вынуждена вводить длинную цепочку дополнительных исключений, то она проигрывает модели, в которой тот же результат возникает как естественный след уже имеющейся структуры.

\subsection{Третий критерий: устойчивость под деформациями}

Следующий критерий связан с динамикой. Хорошая модель должна не только попадать в нужный режим, но и сохранять это попадание при разумных допустимых деформациях. Если малая деформация параметров, окна чтения или условий редукции выбрасывает модель в другой режим без структурного объяснения, это признак хрупкости.

Следовательно, пусть \(M\) --- модель для класса \(\mathcal{C}\). Тогда предпочтение получает та модель, для которой траектории
\begin{equation}
\delta : \mathfrak{R}_M[\mathcal{C}] \longmapsto \mathfrak{R}_{\delta M}[\mathcal{C}]
\end{equation}
остаются внутри одного и того же бассейна притяжения или переходят между режимами только через осмысленные промежуточные слои. Такая устойчивость делает модель не только совместимой, но и structurally robust.

\subsection{Четвёртый критерий: межъязыковая переносимость}

Одно из главных обещаний трактата состоит в том, что одна и та же глубинная сущность должна читаться несколькими языками. Поэтому хорошая модель обязана быть переносимой между естественными окнами чтения. Если модель работает только в одном языке и распадается при переводе в другой, то её explanatory power ограничена.

Значит, конкурентное преимущество имеет та модель, которая позволяет согласованно переходить между геометрическим, спектральным и корреляционным описанием, сохраняя то же структурное положение на карте режимов. Это особенно сильный критерий, потому что он проверяет модель не только на локальную точность, но и на глубину её внутренней архитектуры.

\subsection{Пятый критерий: предсказательная селекция}

Наконец, самый сильный критерий связан с будущими следствиями. Если две модели одинаково хорошо объясняют уже известный класс наблюдаемых, предпочтение должна получать та, которая сильнее ограничивает ещё не классифицированные случаи. Иначе говоря, лучшая модель --- не та, что изящнее комментирует прошлое, а та, что уже сейчас сужает пространство допустимого будущего.

В рамках трактата это можно выразить так: если из модели \(M_1\) следуют более жёсткие ограничения на допустимые траектории по карте, на возможные переходы между режимами или на форму смешанных наблюдаемых, чем из модели \(M_2\), то \(M_1\) структурно предпочтительнее.

\subsection{Составной функционал отбора}

Чтобы собрать все критерии вместе, полезно ввести схематический функционал качества модели:
\begin{equation}
\mathcal{Q}(M\mid \mathcal{C}) =
\lambda_1 \mathcal{R}_{comp}
+ \lambda_2 \mathcal{E}_{str}
+ \lambda_3 \mathcal{S}_{stab}
+ \lambda_4 \mathcal{T}_{lang}
+ \lambda_5 \mathcal{P}_{pred},
\end{equation}
где \(\mathcal{R}_{comp}\) измеряет режимную совместимость, \(\mathcal{E}_{str}\) --- структурную экономию, \(\mathcal{S}_{stab}\) --- устойчивость, \(\mathcal{T}_{lang}\) --- межъязыковую переносимость, а \(\mathcal{P}_{pred}\) --- предсказательную селекцию.

Разумеется, на этой стадии ни сами функционалы, ни коэффициенты \(\lambda_i\) не заданы окончательно. Но сам факт их введения уже важен. Теория получает внутренний принцип предпочтения, а не остаётся в ситуации, где несколько несовместимых моделей сосуществуют без способа сравнения.

\subsection{Пробное применение к уже введённым типам моделей}

Теперь можно схематически представить, как этот отбор работает. Если для класса \([\Delta_{\alpha}]\) одна модель объясняет его через чисто локальную геометрию, а другая --- через глобальную фазовую ветвь, то карта режимов сразу будет благоприятствовать второй, поскольку первая модель искусственно переводит фазовый класс в каркасную область. Если для класса \([m_a]\) одна модель требует множества дополнительных внутренних меток, а другая выводит внутреннюю иерархию из минимального секторного ядра, предпочтение получает вторая. Если же для \([C_{a\alpha}]_{pat}\) одна модель зануляет смешанный вклад, а другая допускает устойчивый смешанный режим, карта и динамика снова будут склонять отбор ко второй.

Тем самым трактат впервые показывает, что его аппарат может не только интерпретировать, но и сравнивать.

\subsection{Критерий поражения модели}

Если есть критерии отбора, то должен быть и критерий поражения. Модель должна считаться структурно неадекватной, если она одновременно:
\begin{itemize}
    \item помещает класс наблюдаемых в неверный режим;
    \item делает это ценой большого числа искусственных допущений;
    \item оказывается неустойчивой при малых деформациях;
    \item не переносится между языками;
    \item не даёт новых ограничений на соседние классы наблюдаемых.
\end{itemize}

Такой набор условий важен, потому что он защищает программу от произвольного удержания плохо работающих моделей. Теория должна не только выбирать лучших кандидатов, но и уметь отпускать слабые конструкции.

\subsection{От отбора моделей к отбору теорий}

В перспективе описанный механизм может работать и на более высоком уровне. Если несколько больших теоретических корпусов претендуют на описание одной и той же глубинной структуры, то и между ними отбор должен проводиться по тем же принципам: режимная совместимость, экономия, устойчивость, переносимость, предсказательная сила. Тогда трактат получит возможность сравнивать не только локальные модели, но и целые исследовательские программы.

Именно здесь становится виден настоящий масштаб построенного аппарата. Он стремится быть не просто новой интерпретацией мира, а новой системой внутренней теоретической селекции.

\subsection{Промежуточный вывод}

Критерии отбора между конкурирующими моделями завершают ещё один важный виток. Теория получает механизм предпочтения, то есть способность не только организовывать и описывать, но и выбирать. А без такого механизма невозможен переход от богатой исследовательской программы к строгой развивающейся теории.

Следующий шаг должен состоять в том, чтобы проверить, как именно этот механизм отбора может работать на одном полном сравнении двух конкурирующих схем для одного и того же класса наблюдаемых. Именно это станет первым реальным испытанием зрелости трактата.


\section{Первое полное сравнение конкурирующих схем}

Теперь аппарат трактата должен пройти первое настоящее испытание. До сих пор он строил язык, инварианты, карту режимов, фазовый портрет и критерии внутреннего отбора. Но всё это останется предварительной архитектурой, если теория не покажет, что умеет сравнивать две реальные конкурирующие схемы объяснения одного и того же класса наблюдаемых. Именно поэтому теперь необходимо провести первый полный сравнительный цикл.

\subsection{Выбор класса наблюдаемых}

В качестве пробного класса возьмём смешанные коэффициенты отклика типа \([C_{a\alpha}]_{pat}\). Этот выбор не случаен. Именно такие величины наиболее трудны для теории, потому что они одновременно чувствительны к фазовой ветви и к внутреннему каналу. Следовательно, они особенно подходят для различения сильной и слабой моделей. Если новая программа не умеет работать именно с такими наблюдаемыми, то её притязание на описание глубинной согласованности остаётся неполным.

\subsection{Две конкурирующие схемы}

Пусть имеются две конкурирующие интерпретации.

\textbf{Схема A: редукция к каркасно-фазовой модификации.} В этой схеме предполагается, что наблюдаемое \([C_{a\alpha}]_{pat}\) в основе своей определяется каркасом и глобальной фазой, а внутренняя индексная структура \(a\) есть лишь вторичная маркировка, не несущая самостоятельного динамического содержания.

\textbf{Схема B: подлинно смешанный секторный механизм.} Здесь предполагается, что \([C_{a\alpha}]_{pat}\) есть несводимый результат совместного действия фазового сектора и внутренней нагрузки, так что ни один из них не может быть исключён без потери устойчивого паттерна.

Сразу видно, что обе схемы могут быть подогнаны под часть данных. Следовательно, именно здесь и нужен аппарат структурного отбора.

\subsection{Шаг первый: режимная локализация}

Сначала сравним, в какие области карты режимов каждая схема помещает рассматриваемый класс. Схема A естественно стремится интерпретировать \([C_{a\alpha}]_{pat}\) как почти фазовое наблюдаемое с малой внутренней коррекцией. Это означает локализацию в области с
\begin{equation}
\Xi_A \gg 1,
\qquad
\Upsilon_A > 0,
\qquad
|\Upsilon_A| \text{ заметно отлично от } 0.
\end{equation}

Схема B, напротив, помещает класс в окрестность смешанного режима:
\begin{equation}
\Xi_B \gg 1,
\qquad
\Upsilon_B \approx 0,
\qquad
w_{\times} \not\approx 0.
\end{equation}

Если уже исходные данные реконструкции указывают на устойчивый ненулевой смешанный остаток и на отсутствие чёткой фазовой ориентации, то первая проверка сразу работает против схемы A. Она локализует класс не в том режиме, куда его ведёт сама карта наблюдаемых.

\subsection{Шаг второй: структурная экономия}

Теперь сравним степень искусственности двух схем. Схема A вынуждена интерпретировать внутренний индекс \(a\) как почти декоративную метку, но при этом всё же как-то объяснять устойчивый паттерн зависимости по \(a\). Это значит, что внутренняя структура фактически сохраняется в данных, но официально не признаётся в ядре модели. В результате схема A вынуждена добавлять специальные правила для восстановления того, что сама же заранее объявляет вторичным.

Схема B, напротив, напрямую включает внутреннюю нагрузку в механизм объяснения. Она не требует дополнительных заплаток для объяснения зависимостей по \(a\), потому что эти зависимости входят в саму архитектуру модели. Следовательно, по критерию структурной экономии схема B оказывается предпочтительнее: она меньше скрывает в исключениях и больше выводит из основного инвариантного ядра.

\subsection{Шаг третий: устойчивость при деформациях}

Теперь важнейшая проверка: как ведут себя схемы при допустимых слабых деформациях? Если слегка менять внутреннюю канализацию, не разрушая глобальную фазу, то схема A должна была бы сохранять почти тот же характер объяснения, поскольку внутренний сектор в ней вторичен. Но именно это и становится проблемой: данные класса \([C_{a\alpha}]_{pat}\) по построению чувствительны к внутреннему индексу. Следовательно, малая внутренняя перестройка должна менять их паттерн сильнее, чем допускает схема A.

Схема B проходит этот тест лучше. Она уже заранее учитывает, что класс расположен в смешанной области, а значит чувствителен к обоим типам деформации. Для неё изменение паттерна при внутренней перестройке не есть аномалия, а естественное следствие положения класса возле \(P_{mix}\).

Поэтому по критерию устойчивости под деформациями схема B снова выигрывает: она не ломается именно там, где данные и должны быть тонко чувствительны.

\subsection{Шаг четвёртый: межъязыковая переносимость}

Рассмотрим теперь переносимость между окнами чтения. В спектральном языке схема A ещё может выглядеть приемлемо, если трактовать паттерн \([C_{a\alpha}]_{pat}\) как фазово-модифицированный спектральный отклик с дополнительной индексацией. Но при переходе к корреляционному или внутренне-операторному языку выясняется, что внутренняя структура перестаёт быть вторичной: она непосредственно входит в организацию корреляций и каналов.

Схема B, напротив, переносится между языками естественнее. В спектральном окне она описывает смешанное расщепление, в корреляционном --- смешанную организацию связей, в операторном --- совместную фазово-внутреннюю канализацию. Следовательно, одна и та же структурная логика сохраняется в разных редукциях. Именно это и делает схему B сильнее по критерию межъязыковой переносимости.

\subsection{Шаг пятый: предсказательная селекция}

Теперь посмотрим на наиболее сильный критерий. Схема A фактически допускает широкое семейство моделей, в которых внутренняя зависимость может быть почти произвольной малой коррекцией к фазовой структуре. Это делает её гибкой, но одновременно и слабо ограничивающей. Схема B, напротив, накладывает более жёсткое требование: если класс действительно принадлежит смешанному режиму, то изменения фазовой ветви и внутреннего канала не могут происходить независимо. Они должны отражаться в согласованной перестройке паттерна \([C_{a\alpha}]_{pat}\).

Следовательно, схема B исключает больше допустимых форм поведения и потому обладает большей селективной силой. В рамках трактата именно это и должно считаться признаком более зрелой модели.

\subsection{Сравнение через функционал качества}

Теперь можно собрать всё вместе. Пусть для схем A и B вычисляются или хотя бы качественно оцениваются компоненты функционала
\begin{equation}
\mathcal{Q}(M\mid [C_{a\alpha}]_{pat}).
\end{equation}

Тогда картина оказывается следующей:
\begin{itemize}
    \item по режимной совместимости выигрывает схема B;
    \item по структурной экономии выигрывает схема B;
    \item по устойчивости под деформациями выигрывает схема B;
    \item по межъязыковой переносимости выигрывает схема B;
    \item по предсказательной селекции выигрывает схема B.
\end{itemize}

Это означает, что даже без полной численной реализации трактат уже обладает достаточной структурой, чтобы предпочесть одну схему другой на основе внутренних теоретических критериев, а не только на основе интуитивного вкуса.

\subsection{Почему схема A всё же важна}

Нужно, однако, подчеркнуть, что проигрыш схемы A не делает её бесполезной. Она может сохранять эвристическую ценность как приближённая эффективная редукция в тех режимах, где внутренний сектор действительно мал или плохо разрешим. Иначе говоря, трактат не обязан рассматривать конкурирующие схемы только по принципу ``истинно/ложно''. Он может включать их в иерархию приближений.

Это важный шаг зрелости. Теория начинает различать не только победителя и проигравшего, но и относительный статус моделей: фундаментальная, эффективная, переходная, локально полезная. Такая иерархия особенно ценна для реальной научной работы.

\subsection{Первый вывод о зрелости аппарата}

Проведённое сравнение показывает нечто принципиальное. Новый трактат уже не просто создаёт красивый философский язык. Он умеет строить реальные процедуры различения между схемами объяснения, опираясь на собственные инварианты, режимы, динамику и фазовый портрет. Это и есть первый признак теоретической зрелости.

Особенно важно, что победа схемы B достигается не за счёт декларации, а за счёт согласованности сразу нескольких независимых критериев. Именно такая многокритериальная поддержка и должна стать нормой для следующих сравнений.

\subsection{Промежуточный вывод}

Первое полное сравнение двух конкурирующих схем становится для трактата рубежным моментом. Теория впервые демонстрирует способность работать как система внутреннего отбора, а не только как система описания. Тем самым она делает ещё один шаг от исследовательской программы к зарождающейся строгой теории структурной согласованности.

Следующий шаг должен состоять в том, чтобы расширить этот тип анализа: не сравнивать две схемы только для одного класса наблюдаемых, а посмотреть, сохраняется ли их относительный порядок на нескольких разных классах сразу. Именно это даст первый набросок иерархии моделей в полном смысле.


\section{Первая иерархия моделей на нескольких классах наблюдаемых}

Один успешный сравнительный тест ещё не создаёт устойчивой теоретической иерархии. Он показывает, что аппарат трактата способен работать, но не доказывает, что его выводы сохраняются при переходе к другим классам наблюдаемых. Поэтому следующим естественным шагом становится расширенный отбор: нужно проверить, остаётся ли относительный порядок конкурирующих схем тем же самым на нескольких структурно различных классах данных. Только тогда можно говорить не о разовой селекции, а о первой иерархии моделей.

\subsection{Зачем нужен многоклассовый тест}

Причина проста. Модель может выглядеть сильной на одном классе наблюдаемых и слабой на другом. Если теория умеет отбирать только локально, но не умеет строить устойчивый порядок предпочтения на семействе классов, то она ещё не получила общего механизма внутренней селекции. Следовательно, задача состоит в том, чтобы сравнить схемы не в одной точке материала, а на малом, но уже разнообразном корпусе классов.

На текущей стадии естественно взять три опорных класса:
\begin{equation}
\mathcal{C}_{core} = \{\rho_{as}\},
\qquad
\mathcal{C}_{\Phi} = \{[\Delta_{\alpha}]\},
\qquad
\mathcal{C}_{mix} = \{[C_{a\alpha}]_{pat}\}.
\end{equation}

Эти три класса образуют минимальный пробный корпус, потому что покрывают три существенно разные структурные ситуации: каркасное доминирование, фазовое доминирование и смешанный режим.

\subsection{Три пробные схемы}

Для большей содержательности сравнения теперь рассмотрим уже не две, а три схемы.

\textbf{Схема A: каркасно-фазовая редукция.} Она стремится объяснять максимум наблюдаемых через геометрический каркас с фазовыми модификациями.

\textbf{Схема B: смешанный секторный механизм.} Она признаёт самостоятельную роль внутреннего сектора и допускает устойчивые смешанные режимы.

\textbf{Схема C: гипертрофированная внутренняя редукция.} В этой схеме большинство наблюдаемых трактуются как следствия богатой внутренней канализации, тогда как геометрия и глобальная фаза играют подчинённую роль.

Такой набор полезен потому, что создаёт не только оппозицию A/B, но и полноценный треугольник конкуренции: слишком каркасная, слишком внутренняя и смешанно-согласованная схемы.

\subsection{Сравнение на каркасном классе}

Начнём с класса \(\mathcal{C}_{core}\). Здесь ведущая спектральная плотность \(\rho_{as}\) по самому построению должна лежать вблизи каркасного режима. Следовательно:
\begin{itemize}
    \item схема A здесь получает сильное преимущество, поскольку локализует класс в \(P_{core}\) естественно и экономно;
    \item схема B также может объяснить этот класс, но с меньшей специфичностью, поскольку её архитектура шире, чем требуется для чисто каркасного наблюдаемого;
    \item схема C выглядит перегруженной: ей приходится объяснять почти каркасную величину через чрезмерно богатую внутреннюю структуру.
\end{itemize}

На первом классе, следовательно, возникает порядок
\begin{equation}
A \succ B \succ C.
\end{equation}

Это важно: теория не обязана везде отдавать победу одной и той же схеме. Она должна уметь различать области локальной силы.

\subsection{Сравнение на фазовом классе}

Теперь рассмотрим класс \(\mathcal{C}_{\Phi}\). Для фазового расщепления \([\Delta_{\alpha}]\) ожидается устойчивое положение вблизи фазового режима. Тогда:
\begin{itemize}
    \item схема A снова выглядит сильной, поскольку фазовая редукция для неё естественна;
    \item схема B остаётся конкурентоспособной, так как признаёт фазовый сектор как одну из равноправных компонент и не разрушает его автономию;
    \item схема C снова проигрывает, потому что внутренний механизм здесь явно перенасыщает объяснение и требует искусственного перевода фазового класса во внутренний язык.
\end{itemize}

Поэтому на втором классе порядок остаётся тем же:
\begin{equation}
A \succ B \succ C.
\end{equation}

Но здесь уже появляется важный нюанс: разрыв между A и B меньше, чем в каркасном классе, поскольку схема B сохраняет хорошую режимную совместимость и межъязыковую переносимость.

\subsection{Сравнение на смешанном классе}

Наконец, рассмотрим класс \(\mathcal{C}_{mix}\). Здесь ситуация радикально меняется. Для \([C_{a\alpha}]_{pat}\) требуется устойчивая локализация в смешанном режиме, чувствительность к фазе и внутреннему индексу и ненулевой смешанный остаток. Тогда:
\begin{itemize}
    \item схема B получает сильное преимущество, поскольку именно она естественно размещает класс возле \(P_{mix}\) и сохраняет его устойчивость под деформациями;
    \item схема A проигрывает, потому что переводит смешанный класс в слишком фазовую область и вынуждена скрывать внутреннюю структуру в исключениях;
    \item схема C также проигрывает, но по противоположной причине: она пытается внутренне поглотить фазовый вклад и потому теряет баланс смешанного режима.
\end{itemize}

Следовательно, на третьем классе порядок уже иной:
\begin{equation}
B \succ A, C.
\end{equation}

Именно этот переворот особенно важен: он показывает, что многоклассовая иерархия не обязана быть линейной и глобально постоянной без оговорок.

\subsection{От локальных побед к глобальному ранжированию}

Теперь встаёт вопрос: как перейти от трёх локальных порядков к единой иерархии моделей? В рамках трактата естественно использовать агрегированный функционал качества:
\begin{equation}
\overline{\mathcal{Q}}(M) = \sum_{k} \omega_k \, \mathcal{Q}(M\mid \mathcal{C}_k),
\end{equation}
где \(\omega_k\) задают значимость различных классов наблюдаемых.

Если все три класса считать одинаково важными, то схема B получает мощное преимущество за счёт своей сбалансированности. Схема A выигрывает на более простых классах, но резко теряет на смешанном. Схема C проигрывает почти везде, потому что её внутренняя редукция оказывается чрезмерной для каркасных и фазовых классов и не даёт достаточно чистого преимущества даже в смешанном режиме.

Это приводит к первому глобальному порядку:
\begin{equation}
B \succ A \succ C,
\end{equation}
если теория придаёт решающее значение способности работать не только на чистых, но и на genuinely mixed observables.

\subsection{Смысл глобального порядка}

Такой результат имеет глубокий смысл. Схема A оказывается сильной как эффективное приближение для простых режимов. Схема B оказывается наиболее зрелой как общая архитектура, потому что охватывает и чистые, и смешанные классы. Схема C показывает, что простая гипертрофия внутреннего механизма не ведёт к лучшей теории, а лишь перенасыщает объяснение там, где требуется баланс.

Следовательно, первая иерархия моделей уже задаёт структурную типологию:
\begin{itemize}
    \item \textbf{B} --- универсально-перспективная схема;
    \item \textbf{A} --- сильная эффективная схема для ограниченного класса режимов;
    \item \textbf{C} --- локально мотивированная, но глобально перегруженная схема.
\end{itemize}

Это именно тот тип вывода, который невозможен без ранее построенного аппарата режимов, инвариантов и фазового портрета.

\subsection{Зависимость иерархии от веса классов}

Важно, однако, подчеркнуть, что иерархия не абсолютно жёстка. Она зависит от того, какие классы наблюдаемых считать центральными. Если исследователь интересуется почти исключительно каркасными и фазовыми величинами, схема A может выглядеть почти равной B или даже временно предпочтительной. Но как только теория требует охвата смешанных классов и устойчивой межъязыковой переносимости, преимущество смещается к B.

Это очень важный результат зрелости. Теория не навязывает одну-единственную иерархию независимо от задач. Она показывает, как и почему порядок моделей меняется при изменении исследовательского центра тяжести. Тем самым она становится не догмой, а регулируемой системой рационального выбора.

\subsection{Первый критерий универсальности модели}

Теперь можно сформулировать и новый критерий. Универсальной в смысле трактата следует считать такую модель, которая не обязательно выигрывает на каждом отдельном чистом классе, но занимает устойчиво высокое положение на разнообразном наборе структурно разных классов и не проваливается на наиболее трудных смешанных наблюдаемых.

Именно по этому критерию схема B получает первое серьёзное преимущество: она не всегда локально максимальна, но глобально наиболее устойчива по спектру задач. Это и есть признак модели следующего уровня зрелости.

\subsection{Промежуточный вывод}

Первая иерархия моделей на нескольких классах наблюдаемых показывает, что трактат уже умеет строить не только локальные сравнительные акты, но и более общую систему ранжирования. Теория начинает различать эффективные локальные победы и универсальную структурную силу. Именно это приближает её к форме настоящей метатеории отбора.

Следующий шаг должен состоять в том, чтобы перейти от пробной иерархии к вопросу о минимальном ядре теории: какие элементы аппарата \(\mathfrak{S}\), \((\Xi,\Upsilon)\), режимов и селекции являются действительно необходимыми, а какие могут оказаться производными. Это будет означать переход от расширения теории к её внутреннему сжатию.

\section{Минимальное ядро теории}

Теперь теория должна сделать движение, противоположное раннему этапу своего роста. До сих пор трактат сознательно расширял словарь: вводил новые сектора, карты режимов, наблюдаемые классы, функционалы селекции и иерархии моделей. Это расширение было необходимо, потому что без него невозможно было увидеть общую архитектуру. Но зрелость программы определяется не только способностью наращивать аппарат, а и способностью затем сжать его до действительно необходимого остатка.

Вопрос о минимальном ядре теории есть поэтому не вопрос эстетики, а вопрос строгости. Если теория не умеет отличить несущие элементы от вспомогательных, то она рискует спутать собственный язык с собственным предметом. Тогда каждый удачный промежуточный инструмент будет ошибочно объявляться фундаментальным, и аппарат начнёт расти быстрее, чем его объяснительная сила.

Под \textit{минимальным ядром теории} будем понимать такой набор структур, без которого одновременно разрушаются:
\begin{itemize}
    \item идентичность глубинной сущности \(\mathfrak{S}\);
    \item различимость режимов чтения;
    \item возможность построить карту допустимых редукций;
    \item возможность осуществить нетривиальный отбор между конкурирующими моделями.
\end{itemize}

Иначе говоря, минимальное ядро должно быть меньше полного аппарата, но достаточно богато, чтобы теория не распалась на набор несвязанных эвристик.

\subsection{Первый принцип сжатия}

Переход к ядру требует принять простой, но жёсткий принцип:
\begin{quote}
элемент аппарата считается фундаментальным только в том случае, если удаление этого элемента уничтожает хотя бы одну из несущих функций теории, а не merely делает описание менее удобным.
\end{quote}

Этот принцип сразу отделяет необходимость от риторического комфорта. Многие конструкции могут быть полезными для вычисления, иллюстрации или первичного ориентирования, но оставаться нефундаментальными. Теория не обязана сохранять в ядре все полезные объекты; она обязана сохранять только те, без которых теряется сама возможность строгого различения и реконструкции.

\subsection{Что, по-видимому, входит в ядро}

На текущей стадии наиболее естественно считать, что в минимальное ядро должны входить четыре типа структур.

Во-первых, сама сущность \(\mathfrak{S}\) как носитель общей идентичности. Без неё теория снова рассыпается на отдельные геометрические, спектральные и корреляционные языки, между которыми уже нельзя утверждать содержательное тождество.

Во-вторых, инвариантное ядро \(\mathcal{I}(\mathfrak{S})\). Именно оно удерживает различие между допустимой деформацией одной и той же сущности и переходом к иной сущности. Если убрать этот слой, исчезает критерий того, что вообще сохраняется при переходе между реализациями.

В-третьих, семейство режимов чтения. Не обязательно в той подробной форме, в какой они были введены на промежуточных этапах, но по крайней мере как различение каркасного, фазового и смешанного типов проявления. Без этого теория утрачивает способность объяснять, почему одни наблюдаемые локализуются возле \(P_{\mathcal{B}}\), другие --- возле \(P_{\Phi}\), а третьи требуют genuinely mixed структуры.

В-четвёртых, принцип селекции между моделями. Теория, которая не может сравнивать конкурирующие реализации и ранжировать их по качеству, остаётся описательной онтологией, но ещё не становится исследовательской программой. Следовательно, некоторый минимальный функционал отбора должен сохраняться и после сжатия аппарата.

Эту мысль можно записать схематически в виде условия
\begin{equation}
\mathfrak{K}_{min}=\bigl(\mathfrak{S},\mathcal{I}(\mathfrak{S}),\mathcal{R}_{type},\mathcal{Q}_{sel}\bigr),
\end{equation}
где \(\mathcal{R}_{type}\) обозначает минимальную структуру типов редукции, а \(\mathcal{Q}_{sel}\) --- минимальный принцип сравнительного отбора.

\subsection{Что может оказаться производным}

Напротив, целый ряд объектов, введённых на предыдущих этапах, имеет шанс оказаться не элементами ядра, а развёрнутыми представлениями ядра в более удобной форме. К ним, по-видимому, относятся:
\begin{itemize}
    \item подробные фазовые портреты, если они выводимы из более компактной карты режимных переходов;
    \item конкретные координатные схемы для размещения классов наблюдаемых;
    \item разветвлённые системы пробных коэффициентов, если их роль может быть заменена несколькими структурными индексами;
    \item локальные иерархии моделей, если они восстанавливаются из общего селекционного функционала при выборе класса \(\mathcal{C}_k\).
\end{itemize}

Это не означает, что такие объекты следует отбросить. Напротив, они могут остаться незаменимыми на вычислительном, иллюстративном или прикладном уровне. Но теория должна различать: одно дело --- объект, необходимый для существования схемы, и другое --- объект, полезный для её развертывания.

\subsection{Минимальность как проверяемое требование}

Важно подчеркнуть, что минимальность не должна провозглашаться декларативно. Она сама должна стать предметом проверки. Для этого естественно потребовать, чтобы каждый кандидат в ядро проходил следующий тест исключения:
\begin{enumerate}
    \item удалить данный элемент из аппарата;
    \item проверить, сохраняется ли идентичность \(\mathfrak{S}\) и карта допустимых деформаций;
    \item проверить, можно ли по-прежнему различать типы наблюдаемых режимов;
    \item проверить, остаётся ли возможным нетривиальный отбор между моделями.
\end{enumerate}

Если после удаления все четыре функции сохраняются, то соответствующий элемент не принадлежит ядру. Если хотя бы одна из них радикально разрушается, то элемент получает право на фундаментальный статус.

Так возникает более строгая формула экономии: минимальным следует считать не просто самый короткий аппарат, а самый короткий аппарат, который ещё удерживает все несущие операции теории.

\subsection{Сжатие языка и сохранение силы}

У этого шага есть и более глубокий смысл. На ранних стадиях новая теория почти неизбежно говорит длиннее, чем знает. Она вынуждена строить множество промежуточных конструкций, чтобы нащупать собственный предмет. Но по мере роста зрелости язык должен становиться не богаче, а плотнее. Хорошая теория отличается от плохо дисциплинированной не количеством введённых сущностей, а тем, какая часть её объяснительной силы переживает внутреннее сжатие.

Если после сжатия оказывается, что большая часть аппарата была лишь временными лесами, это не дефект, а признак нормального развития. Дефектом было бы обратное: если после удаления декоративных надстроек вместе с ними исчезает вся содержательность. Тогда это означает, что теория ещё не нашла собственное ядро и держалась на множестве вспомогательных подпорок.

\subsection{Новый ориентир для следующего этапа}

Тем самым следующая задача трактата может быть сформулирована уже более строго. Нужно не просто развивать аппарат \(\mathfrak{S}\), а построить процедуру последовательного редукционного сжатия, которая:
\begin{itemize}
    \item проверяет каждый крупный элемент на фундаментальность;
    \item переводит производные конструкции в статус выводимых представлений;
    \item выделяет окончательный короткий список несущих структур;
    \item показывает, что после такого сжатия теория сохраняет предсказательную и различающую силу.
\end{itemize}

Именно этот шаг должен открыть следующую стадию трактата. До сих пор теория училась расширять своё поле зрения. Теперь она должна научиться удерживать мир меньшим числом исходных сущностей. Если ей это удастся, то речь действительно пойдёт не о богатом комментарии к разным языкам физики, а о формировании собственного минимального основания.

\section{Процедура редукционного сжатия}

После постановки вопроса о минимальном ядре теория должна перейти от общей формулы экономии к воспроизводимой процедуре. Недостаточно просто утверждать, что одни элементы фундаментальны, а другие производны. Необходимо описать такой порядок действий, при котором это различие может быть установлено последовательно, повторяемо и с сохранением общей логики трактата.

Под \textit{редукционным сжатием} будем понимать управляемую процедуру, которая переводит расширенный аппарат теории в более короткую форму без утраты его несущих функций. В символическом виде этот переход можно записать как отображение
\begin{equation}
\mathfrak{Red}: \mathfrak{A}_{ext} \longrightarrow \mathfrak{K}_{min},
\end{equation}
где \(\mathfrak{A}_{ext}\) обозначает полный на текущем этапе аппарат трактата, а \(\mathfrak{K}_{min}\) --- его минимальное несущее ядро.

Смысл этого отображения не в механическом удалении терминов, а в проверке того, какие конструкции действительно необходимы для удержания идентичности, режимной различимости, карты редукций и селекционной силы. Поэтому процедура должна быть не лексической, а структурной.

\subsection{Исходный расширенный аппарат}

На текущей стадии разумно считать, что расширенный аппарат теории включает по меньшей мере следующие блоки:
\begin{equation}
\mathfrak{A}_{ext}=\bigl(\mathfrak{S},\mathcal{I}(\mathfrak{S}),\mathcal{R},(\Xi,\Upsilon),\mathfrak{O}_{cls},\mathcal{Q}_{sel},\mathfrak{H}_{mod}\bigr),
\end{equation}
где:
\begin{itemize}
    \item \(\mathcal{R}\) обозначает полный набор режимов и карт переходов;
    \item \((\Xi,\Upsilon)\) обозначают координатно-режимное и селекционное представление структуры;
    \item \(\mathfrak{O}_{cls}\) обозначает семейство классов наблюдаемых;
    \item \(\mathfrak{H}_{mod}\) обозначает текущую иерархию моделей и локальных ранжирований.
\end{itemize}

Эта запись важна тем, что делает явным объект сжатия. Теория должна уметь указать не только то, что она хочет сохранить, но и то, что именно она подвергает проверке на производность.

\subsection{Четыре шага процедуры}

В первом рабочем варианте редукционное сжатие можно организовать как последовательность четырёх шагов.

\textbf{Шаг первый: локализация функции.} Для каждого элемента аппарата необходимо указать, какую именно несущую функцию он выполняет. Элемент, чья функция не может быть названа ясно, уже тем самым подозрителен: возможно, он лишь повторяет работу другого слоя под более длинным именем.

\textbf{Шаг второй: пробное исключение.} После локализации функции элемент временно удаляется из схемы, а его вклад считается запрещённым. На этом этапе важно проверять не удобство изложения, а возможность сохранить основные операции теории.

\textbf{Шаг третий: реконструкция без элемента.} Если после удаления теория всё ещё способна восстановить ту же структурную роль через другие, более короткие средства, то рассматриваемый элемент должен быть переведён в класс производных представлений.

\textbf{Шаг четвёртый: фиксация статуса.} После проверки элемент получает один из трёх статусов: фундаментальный, выводимый или временно неопределённый. Последний статус особенно важен, потому что не всякая структура уже сейчас может быть надёжно классифицирована.

\subsection{Критерии сохранения после сжатия}

Чтобы процедура не свелась к произвольному минимализму, необходимо заранее фиксировать, что именно обязано сохраниться после каждого шага. Минимально следует требовать сохранения четырёх классов свойств:
\begin{enumerate}
    \item \textbf{идентичности} --- теория должна по-прежнему различать, где речь идёт об одной сущности \(\mathfrak{S}\), а где уже о переходе к иной сущности;
    \item \textbf{режимной различимости} --- теория должна сохранять различие между каркасным, фазовым и смешанным проявлениями;
    \item \textbf{редукционной управляемости} --- должна оставаться возможность переходить от глубинного описания к наблюдаемым классам без распада схемы;
    \item \textbf{селекционной нетривиальности} --- теория должна сохранять способность различать конкурирующие модели не риторически, а по структурному качеству.
\end{enumerate}

Если хотя бы одно из этих свойств исчезает, сжатие оказалось чрезмерным. Если же все четыре свойства сохраняются, то сокращение можно считать допустимым и внутренне легитимным.

\subsection{Три возможных исхода для каждого элемента}

После применения процедуры к конкретному элементу \(X\subset \mathfrak{A}_{ext}\) теория должна относить его к одному из трёх классов.

\textbf{Фундаментальный класс.} Элемент входит в ядро, если его исключение разрушает хотя бы одну несущую функцию и если эта функция не восстанавливается более коротким путём.

\textbf{Производный класс.} Элемент считается производным, если его структурная роль полностью реконструируется из уже сохранённых компонентов ядра. В этом случае он остаётся в теории, но уже как развертка, координатизация или удобная надстройка.

\textbf{Пограничный класс.} Элемент считается пограничным, если на текущей стадии нельзя решить, фундаментален он или производен, поскольку для этого ещё не хватает строгой реконструкции. Такой статус не является слабостью; напротив, он фиксирует точку реальной незавершённости программы.

Схематически это можно выразить так:
\begin{equation}
\mathrm{Stat}(X) \in \{\mathrm{fund},\mathrm{der},\mathrm{border}\}.
\end{equation}

\subsection{Первая пробная классификация текущего аппарата}

Если применить описанную процедуру к уже введённому языку трактата, то можно получить первую, ещё не окончательную классификацию.

Наиболее сильными кандидатами в фундаментальный класс остаются \(\mathfrak{S}\), \(\mathcal{I}(\mathfrak{S})\), минимальное различение типов редукции и сам принцип селекции. Без них теория почти сразу утрачивает либо идентичность предмета, либо способность сравнивать модели.

Напротив, объекты вида \((\Xi,\Upsilon)\), детальные карты расположения классов наблюдаемых и локальные иерархии семейства \(\mathfrak{H}_{mod}\) пока естественнее рассматривать как кандидаты в производный класс. Их роль кажется очень важной для анализа, но ещё не очевидно, что без них теория теряет саму возможность существовать.

Наконец, пограничный статус, по-видимому, получают некоторые промежуточные структуры, связывающие режимы с конкретными формами наблюдаемой локализации. Возможно, позже они будут выведены из более короткого механизма. Но возможно и обратное: именно они окажутся тем отсутствующим связующим слоем, без которого ядро не сможет работать на реальном материале.

\subsection{Сжатие как новый тип строгости}

Важнейший результат этой процедуры состоит в том, что строгость теории начинает измеряться не только способностью что-либо ввести, но и способностью затем что-либо удалить без потери содержательности. Это меняет сам стиль дальнейшей работы. Теперь теория должна доказывать свою зрелость не экспансией, а дисциплинированным сохранением необходимого.

Такой поворот особенно важен для больших синтетических программ. Они почти всегда соблазняются избыточностью: каждый новый успешный инструмент кажется слишком ценным, чтобы отказаться от него как от элемента ядра. Но именно здесь и начинается зрелая селекция. Теория должна научиться различать внутренне необходимое и исторически полезное.

\subsection{Следующая конкретная задача}

После введения самой процедуры следующий этап можно поставить уже почти в технической форме. Нужно построить таблицу или карту редукционного статуса для основных элементов аппарата и пройти её последовательно, начиная с наиболее нагруженных объектов. Иначе говоря, дальнейшая работа должна ответить на вопросы:
\begin{itemize}
    \item можно ли вывести \((\Xi,\Upsilon)\) из более короткой режимно-селекционной структуры;
    \item можно ли восстановить локальные иерархии моделей из одного общего функционала \(\mathcal{Q}_{sel}\);
    \item какие элементы семейства наблюдаемых \(\mathfrak{O}_{cls}\) являются самостоятельными, а какие лишь удобными агрегатами;
    \item существует ли ещё более короткая запись для самого набора \(\mathfrak{K}_{min}\).
\end{itemize}

Тем самым трактат получает уже не только философское направление, но и внутренний протокол самосжатия. Если предыдущие разделы строили язык расширения, то настоящий раздел вводит язык дисциплины. Именно их соединение и должно со временем сделать программу не просто богатой, но и действительно фундаментальной.

\section{Статусная карта основных элементов аппарата}

Теперь можно сделать следующий шаг и перевести общую процедуру в форму первой явной карты статусов. Такая карта ещё не является окончательным решением. Но она важна как переход от общих рассуждений о сжатии к дисциплинированной инвентаризации тех объектов, с которыми теория уже фактически работает.

Именно здесь становится ясно, насколько программа способна смотреть на саму себя как на объект анализа. Если теория умеет не только строить новые структуры, но и присваивать им строгий внутренний статус, то она начинает приобретать форму зрелой архитектуры, а не просто растущего корпуса мотивов.

\subsection{Формат статусной записи}

Для каждого базового элемента \(X\) введём запись вида
\begin{equation}
\mathbf{Card}(X)=\bigl(\mathrm{role}(X),\mathrm{risk}(X),\mathrm{stat}(X)\bigr),
\end{equation}
где:
\begin{itemize}
    \item \(\mathrm{role}(X)\) обозначает основную несущую функцию элемента;
    \item \(\mathrm{risk}(X)\) обозначает, что именно разрушается при его удалении;
    \item \(\mathrm{stat}(X)\) принимает одно из значений \(\mathrm{fund}\), \(\mathrm{der}\) или \(\mathrm{border}\).
\end{itemize}

Такая запись полезна тем, что она отделяет смысл элемента от его исторической роли в тексте. Теория должна судить не по тому, насколько часто объект использовался, а по тому, какую незаменимую работу он действительно выполняет.

\subsection{Карта для ядровых кандидатов}

Для первой рабочей версии трактата естественно предложить следующую карту.

\begin{center}
\begin{tabular}{|>{\raggedright\arraybackslash}p{2.4cm}|>{\raggedright\arraybackslash}p{3.8cm}|>{\raggedright\arraybackslash}p{5.4cm}|>{\raggedright\arraybackslash}p{2.0cm}|}
\hline
Элемент & Основная роль & Риск при удалении & Статус \\
\hline
\(\mathfrak{S}\) & носитель идентичности & распад единого предмета & fund \\
\hline
\(\mathcal{I}(\mathfrak{S})\) & сохранение тождества под деформациями & исчезновение критерия инвариантности & fund \\
\hline
\(\mathcal{R}_{type}\) & различение типов чтения & смешение режимов проявления & fund \\
\hline
\(\mathcal{Q}_{sel}\) & сравнительный отбор моделей & утрата нетривиальной селекции & fund \\
\hline
\((\Xi,\Upsilon)\) & координатизация и режимная визуализация & потеря аналитического удобства & der \\
\hline
\(\mathfrak{O}_{cls}\) & организация наблюдаемых по классам & ослабление операциональной читаемости & border \\
\hline
\(\mathfrak{H}_{mod}\) & локальные иерархии моделей & потеря детального ранжирования & der \\
\hline
\end{tabular}
\end{center}

Разумеется, эта таблица пока имеет рабочий характер. Её цель не в том, чтобы закрыть вопрос, а в том, чтобы задать первую дисциплинированную матрицу для дальнейшей проверки.

\subsection{Почему \texorpdfstring{$\mathfrak{S}$}{S} получает фундаментальный статус}

Статус \(\mathfrak{S}\) как фундаментального элемента почти неизбежен. Именно он удерживает саму возможность говорить, что геометрическое, спектральное и корреляционное чтения относятся к одному и тому же источнику. Если убрать \(\mathfrak{S}\), то теория не просто потеряет одно из обозначений; она утратит свой предмет как единый предмет.

В этом смысле \(\mathfrak{S}\) не является одним объектом среди других. Он занимает метаструктурное место: через него трактат утверждает, что различие языков не есть различие сущностей. Именно поэтому его нельзя заменить простой совокупностью редукций без риска распада всей программы.

\subsection{Почему инвариантное ядро незаменимо}

Если \(\mathfrak{S}\) удерживает единство предмета, то \(\mathcal{I}(\mathfrak{S})\) удерживает его идентичность при деформациях. Без этого объекта теория потеряла бы критерий различия между вариацией одной и той же сущности и переходом к иной сущности. Тогда все разговоры о режимах, редукциях и устойчивых классах становились бы двусмысленными.

Поэтому \(\mathcal{I}(\mathfrak{S})\) следует считать вторым безусловным кандидатом в ядро. Если \(\mathfrak{S}\) отвечает на вопрос ``что является единым источником?'', то \(\mathcal{I}(\mathfrak{S})\) отвечает на вопрос ``что в этом источнике не имеет права исчезнуть?''

\subsection{Промежуточный статус семейства наблюдаемых}

Наиболее интересен сейчас не фундаментальный и не производный класс, а именно пограничный. В первую очередь это относится к семейству \(\mathfrak{O}_{cls}\). С одной стороны, без организации наблюдаемых по классам теория быстро теряет операциональную ясность: ей становится труднее связывать глубинные различия с проверяемыми режимами проявления. С другой стороны, пока ещё не доказано, что это семейство не может быть полностью восстановлено из более короткой комбинации типов редукции и селекционного функционала.

Поэтому для \(\mathfrak{O}_{cls}\) наиболее честным остаётся статус \(\mathrm{border}\). Он фиксирует реальную неопределённость: возможно, перед нами самостоятельный несущий слой; возможно, лишь удобная агрегирующая оболочка для более простых структур.

\subsection{Производные элементы как необходимые инструменты}

Важно также не спутать производность с незначительностью. Объекты \((\Xi,\Upsilon)\) и \(\mathfrak{H}_{mod}\) разумно пока считать производными не потому, что они слабы, а потому, что их функция, по-видимому, заключается прежде всего в развёртывании, визуализации и уточнении уже существующих различий. Они делают аппарат более управляемым, но ещё не доказано, что делают его возможным в принципе.

Именно эта разница и должна стать центральной для следующего этапа. Теория обязана уважать силу производных конструкций как исследовательских инструментов, не превращая их автоматически в онтологическое ядро.

\subsection{Первый почти-аксиоматический вывод}

На этом этапе можно зафиксировать первую почти-аксиоматическую формулу статусов:
\begin{align}
&\mathrm{Stat}(\mathfrak{S})=\mathrm{fund}, \\
&\mathrm{Stat}(\mathcal{I}(\mathfrak{S}))=\mathrm{fund}, \\
&\mathrm{Stat}(\mathcal{R}_{type})=\mathrm{fund}, \\
&\mathrm{Stat}(\mathcal{Q}_{sel})=\mathrm{fund}, \\
&\mathrm{Stat}(\Xi,\Upsilon)=\mathrm{der}, \\
&\mathrm{Stat}(\mathfrak{O}_{cls})=\mathrm{border}, \\
&\mathrm{Stat}(\mathfrak{H}_{mod})=\mathrm{der}.
\end{align}

Эта формула не претендует на окончательность. Её статус --- методологический. Она нужна как временная ось, относительно которой уже можно проверять дальнейшие реконструкции и уточнения.

\subsection{Что должно последовать дальше}

Теперь у трактата появляется уже очень конкретная следующая задача: для каждого элемента из статусной карты нужно либо построить явную реконструкцию из ядра, либо доказать, что такая реконструкция невозможна. Иначе говоря, следующая стадия должна состоять не просто в повторении статусов, а в их доказательной развязке.

Это означает по меньшей мере три направления работы:
\begin{itemize}
    \item построить вывод \((\Xi,\Upsilon)\) из более короткой структуры \((\mathcal{R}_{type},\mathcal{Q}_{sel})\);
    \item выяснить, можно ли вывести \(\mathfrak{O}_{cls}\) как образ редукционного действия ядра на семейство наблюдаемых;
    \item проверить, является ли \(\mathfrak{H}_{mod}\) всего лишь развёрткой глобального качества \(\mathcal{Q}_{sel}\) по различным классам данных.
\end{itemize}

Если эти шаги удастся выполнить, то трактат перейдёт от предварительной карты статусов к первому собственному доказательству внутренней экономии. Именно это и станет настоящим экзаменом для всей программы редукционного сжатия.

\section{Первая реконструкция: вывод \texorpdfstring{$(\Xi,\Upsilon)$}{(Xi,Upsilon)} из \texorpdfstring{$(\mathcal{R}_{type},\mathcal{Q}_{sel})$}{(Rtype,Qsel)}}

Теперь можно сделать первый по-настоящему содержательный шаг в программе редукционного сжатия. До сих пор статус \((\Xi,\Upsilon)\) как производного элемента был лишь рабочей гипотезой. Чтобы этот статус не оставался декларативным, необходимо показать хотя бы в первом приближении, каким образом данная пара может быть восстановлена из более короткого набора структур.

Именно поэтому в качестве первого теста следует взять наиболее удобный случай: попытаться вывести \((\Xi,\Upsilon)\) из сочетания типологии редукций \(\mathcal{R}_{type}\) и селекционного функционала \(\mathcal{Q}_{sel}\). Если такая реконструкция удаётся хотя бы на структурном уровне, то производный статус пары \((\Xi,\Upsilon)\) получает реальное основание.

\subsection{Что именно должны означать \texorpdfstring{$\Xi$}{Xi} и \texorpdfstring{$\Upsilon$}{Upsilon}}

Для начала нужно уточнить, что именно мы хотим реконструировать. В языке трактата \(\Xi\) естественно понимать как координатно-режимное представление пространства проявлений, то есть как способ размещения классов наблюдаемых и моделей относительно основных типов чтения. Соответственно, \(\Upsilon\) следует понимать как слой селекционной оценки, который упорядочивает эти размещения по качеству, устойчивости и межъязыковой переносимости.

Иначе говоря, \(\Xi\) отвечает прежде всего на вопрос ``где в структурном пространстве находится данный объект?'', а \(\Upsilon\) --- на вопрос ``какой вес и какое качество имеет это положение с точки зрения теории?'' Если удастся показать, что оба этих вопроса уже содержатся в \((\mathcal{R}_{type},\mathcal{Q}_{sel})\), то отдельный фундаментальный статус пары \((\Xi,\Upsilon)\) станет избыточным.

\subsection{Реконструкция координатного слоя}

Начнём с \(\Xi\). Пусть типология редукций задаёт минимальное множество режимов
\begin{equation}
\mathcal{R}_{type}=\{R_{\mathcal{B}},R_{\Phi},R_{mix}\},
\end{equation}
где каждый элемент обозначает соответственно каркасное, фазовое или смешанное проявление. Тогда любой объект теории, будь то наблюдаемый класс, модель или переход, может быть охарактеризован не собственной абсолютной координатой, а своим \textit{профилем принадлежности} к этим режимам.

Этот профиль удобно записать как отображение
\begin{equation}
\Xi_X: X \longmapsto \bigl(\xi_{\mathcal{B}}(X),\xi_{\Phi}(X),\xi_{mix}(X)\bigr),
\end{equation}
где компоненты \(\xi_{\alpha}(X)\) выражают не геометрические координаты в обычном смысле, а относительную структурную близость объекта \(X\) к каждому из трёх режимов.

Но именно такая конструкция уже фактически порождается самой типологией \(\mathcal{R}_{type}\). Как только теория умеет различать три типа чтения и относить объект к тому или иному типу доминирования, она тем самым уже задаёт прообраз координатного пространства. Следовательно, \(\Xi\) может быть понята не как дополнительная первичная структура, а как развертка режима принадлежности, содержащегося в \(\mathcal{R}_{type}\).

\subsection{Реконструкция селекционного слоя}

Теперь перейдём к \(\Upsilon\). Если \(\mathcal{Q}_{sel}\) есть глобальный функционал качества, то для каждого объекта или класса объектов можно определить индуцированный вес
\begin{equation}
\Upsilon(X)=\mathcal{Q}_{sel}(X\mid \mathcal{R}_{type}),
\end{equation}
то есть значение селекционного качества, вычисляемое не в пустоте, а относительно режимного профиля данного объекта.

Это означает, что \(\Upsilon\) не обязана вводиться как отдельный фундаментальный слой. Она может быть понята как действие \(\mathcal{Q}_{sel}\) на уже реконструированное пространство режимных профилей. Тогда \(\Upsilon\) есть не новая сущность, а оценочная тень селекционного функционала на пространстве \(\Xi\).

Именно здесь становится виден ключевой принцип экономии: сначала типология редукций создаёт пространство различимости, затем селекционный функционал создаёт на этом пространстве иерархию качеств. В таком порядке пара \((\Xi,\Upsilon)\) рождается естественно и не требует самостоятельного онтологического постулирования.

\subsection{Схематическая формула реконструкции}

Всё сказанное можно записать в компактном виде. Если обозначить через
\begin{equation}
\Xi = \mathrm{Prof}(\mathcal{R}_{type}),
\end{equation}
пространство профилей режимной принадлежности, а через
\begin{equation}
\Upsilon = \mathcal{Q}_{sel}\big|_{\Xi}
\end{equation}
индуцированный слой качества на этом пространстве, то получаем реконструкцию
\begin{equation}
(\Xi,\Upsilon)=\mathfrak{Rec}\bigl(\mathcal{R}_{type},\mathcal{Q}_{sel}\bigr).
\end{equation}

Смысл этой формулы в том, что пара \((\Xi,\Upsilon)\) не исчезает из теории, а меняет свой статус. Раньше она выглядела как самостоятельный несущий блок. Теперь она читается как выводимая форма организации уже более короткого ядра.

\subsection{Почему это ещё не полное доказательство}

Разумеется, полученная реконструкция пока не является строгой теоремой. Она остаётся структурной схемой. Здесь ещё не задано точное пространство значений для \(\xi_{\mathcal{B}},\xi_{\Phi},\xi_{mix}\), не определён полный класс допустимых объектов \(X\), и не выписаны аксиомы, которым должен удовлетворять функционал \(\mathcal{Q}_{sel}\).

Однако этого уже достаточно для методологического результата. Производный статус \((\Xi,\Upsilon)\) теперь опирается не на интуицию, а на ясную схему порождения: сначала различимость по типам редукции, затем селекционная оценка на пространстве этой различимости.

\subsection{Критерий успеха реконструкции}

Чтобы считать реконструкцию успешной не только риторически, нужно указать, что именно она обязана сохранить. Минимально следует требовать трёх свойств:
\begin{enumerate}
    \item реконструированное \(\Xi\) должно уметь различать каркасное, фазовое и смешанное доминирование;
    \item реконструированное \(\Upsilon\) должно воспроизводить хотя бы тот же качественный порядок моделей и классов, что и прежняя версия аппарата;
    \item переход от \((\mathcal{R}_{type},\mathcal{Q}_{sel})\) к \((\Xi,\Upsilon)\) не должен требовать добавления новых фундаментальных сущностей.
\end{enumerate}

Если эти три условия выполнены, то пара \((\Xi,\Upsilon)\) действительно получает право считаться производной. Если же хотя бы одно условие рушится, статус \(\mathrm{der}\) придётся пересмотреть.

\subsection{Промежуточный вывод}

Итак, первый тест редукционного сжатия даёт положительный, хотя и пока предварительный результат. Координатно-режимный слой \(\Xi\) можно понимать как пространство профилей, индуцированное самой типологией редукций. Селекционный слой \(\Upsilon\) можно понимать как действие общего функционала качества на это пространство. Тем самым пара \((\Xi,\Upsilon)\) оказывается не независимым основанием теории, а организованной развёрткой уже имеющихся фундаментальных структур.

Это и есть первый настоящий успех программы внутренней экономии. Теория впервые не просто объявляет некоторый блок производным, а показывает направление его реконструкции из более короткого ядра.

\subsection{Следующий шаг после первой реконструкции}

После этого естественно перейти к следующей, более трудной цели. Если пара \((\Xi,\Upsilon)\) действительно реконструируема, то следующий вопрос состоит в том, можно ли аналогичным образом реконструировать и семейство \(\mathfrak{O}_{cls}\). Именно там находится самый интересный пограничный случай: объект, который ещё не ясно, считать ли самостоятельным слоем теории или лишь образом действия уже выделенного ядра.

Следовательно, дальнейший ход трактата должен двигаться от доказуемой производности к анализу пограничности. Это будет означать следующий уровень зрелости: теория научится различать не только фундаментальное и выводимое, но и по-настоящему проблемные случаи своего собственного аппарата.

\section{Пограничный случай: можно ли вывести \texorpdfstring{$\mathfrak{O}_{cls}$}{Ocls} из действия ядра}

После первой успешной реконструкции естественно перейти к более трудной задаче. Если пара \((\Xi,\Upsilon)\) оказалась хотя бы структурно выводимой, то теперь нужно рассмотреть объект иного типа --- не явно производный, а по-настоящему пограничный. Именно таким объектом в текущем аппарате выступает семейство классов наблюдаемых \(\mathfrak{O}_{cls}\).

Его положение принципиально сложнее. В отличие от \((\Xi,\Upsilon)\), семейство \(\mathfrak{O}_{cls}\) не выглядит просто координатной или оценочной разверткой уже введённых структур. Оно связано с самой организацией того, что теория вообще считает наблюдаемым, сравнимым и типологически устойчивым. Поэтому здесь задача состоит не в быстром объявлении производности, а в аккуратной проверке того, насколько далеко может продвинуться реконструкция.

\subsection{Что понимается под семейством \texorpdfstring{$\mathfrak{O}_{cls}$}{Ocls}}

Под \(\mathfrak{O}_{cls}\) будем понимать не просто список наблюдаемых величин, а организованное семейство классов, в котором отдельные наблюдаемые группируются по типу структурной чувствительности. Иначе говоря, речь идёт не о множестве изолированных чисел, а о классификации того, к какому режиму, какому сектору и какому типу смешения преимущественно чувствителен тот или иной наблюдаемый объект.

Именно поэтому \(\mathfrak{O}_{cls}\) играет в трактате особую роль. Оно образует мост между глубинной архитектурой сущности и тем уровнем, на котором теория начинает делать различимые утверждения о классах явлений. Если этот мост целиком выводим, то семейство \(\mathfrak{O}_{cls}\) оказывается производным. Если же нет, то оно, возможно, выражает самостоятельный слой теории.

\subsection{Пробная схема реконструкции}

Пусть ядро теории задано набором
\begin{equation}
\mathfrak{K}_{min}=\bigl(\mathfrak{S},\mathcal{I}(\mathfrak{S}),\mathcal{R}_{type},\mathcal{Q}_{sel}\bigr).
\end{equation}
Тогда естественно попробовать определить семейство наблюдаемых классов как образ действия ядра на допустимое множество наблюдаемых величин \(\mathfrak{O}\):
\begin{equation}
\mathfrak{O}_{cls}=\mathrm{Cls}_{\mathfrak{K}}(\mathfrak{O}),
\end{equation}
где оператор \(\mathrm{Cls}_{\mathfrak{K}}\) должен для каждого наблюдаемого объекта определять:
\begin{enumerate}
    \item его режимный профиль относительно \(\mathcal{R}_{type}\);
    \item его селекционный вес относительно \(\mathcal{Q}_{sel}\);
    \item его устойчивость под допустимыми деформациями, контролируемыми \(\mathcal{I}(\mathfrak{S})\).
\end{enumerate}

Если такой оператор существует, то классы наблюдаемых действительно могут быть поняты как результат действия ядра, а не как отдельный первичный словарь.

\subsection{Как ядро может порождать классы}

Идея здесь следующая. Типология редукций \(\mathcal{R}_{type}\) даёт первый разрез: она разделяет наблюдаемые по доминирующему режиму проявления. Функционал \(\mathcal{Q}_{sel}\) даёт второй разрез: он различает более и менее устойчивые, более и менее универсальные, более и менее межъязыково переносимые наблюдаемые структуры. Наконец, инвариантное ядро \(\mathcal{I}(\mathfrak{S})\) даёт третий разрез: оно отделяет те различия, которые сохраняются под допустимыми деформациями, от тех, которые являются лишь поверхностными вариациями.

Вместе эти три слоя уже способны породить предварительную классификацию. Два наблюдаемых попадают в один и тот же класс тогда и только тогда, когда:
\begin{itemize}
    \item их режимные профили эквивалентны в грубом смысле;
    \item их селекционный ранг лежит в одной и той же устойчивой полосе;
    \item их поведение под допустимыми деформациями не разрушает общей структурной идентичности класса.
\end{itemize}

Следовательно, сама идея наблюдаемого класса действительно может быть частично реконструирована как композиция режимного, селекционного и инвариантного действия ядра.

\subsection{Почему реконструкция остаётся неполной}

Однако именно здесь обнаруживается граница, отличающая пограничный случай от чисто производного. Для полной реконструкции \(\mathfrak{O}_{cls}\) недостаточно лишь знать, как ядро различает режимы и ранги. Нужно ещё знать, \textit{какие именно величины допускаются в качестве наблюдаемых носителей этих различий}. Иначе говоря, требуется не только внутренний аппарат классификации, но и правило доступа к самому пространству \(\mathfrak{O}\).

Вот здесь и возникает трудность. Возможно, множество допустимых наблюдаемых полностью задаётся редукциями сущности \(\mathfrak{S}\). Но возможно и другое: чтобы выделить физически осмысленные классы, теория должна дополнительно постулировать некоторую операциональную структуру наблюдаемости. Если верно второе, то \(\mathfrak{O}_{cls}\) уже не сводится без остатка к действию ядра.

Именно поэтому текущая реконструкция доводит нас лишь до частичного вывода:
\begin{equation}
\mathfrak{O}_{cls}\subseteq \mathrm{Im}\,\mathrm{Cls}_{\mathfrak{K}},
\end{equation}
но ещё не позволяет отождествить эти объекты полностью.

\subsection{Что именно удерживает пограничный статус}

Теперь можно точнее понять, почему статус \(\mathrm{border}\) для \(\mathfrak{O}_{cls}\) остаётся оправданным. С одной стороны, семейство классов очевидным образом зависит от ядра: без режимов, селекции и инвариантов оно почти сразу теряет свою структуру. С другой стороны, ещё не доказано, что ядро само по себе уже достаточно, чтобы породить весь язык наблюдаемых классов без дополнительных предпосылок о допустимой операциональной читаемости.

Иначе говоря, \(\mathfrak{O}_{cls}\) не выглядит полностью независимым от ядра, но и не сводится к нему с той лёгкостью, с какой сводится пара \((\Xi,\Upsilon)\). Именно этот избыток и есть настоящий смысл пограничности.

\subsection{Проверочный критерий для \texorpdfstring{$\mathfrak{O}_{cls}$}{Ocls}}

Чтобы в будущем перевести статус \(\mathfrak{O}_{cls}\) либо в \(\mathrm{fund}\), либо в \(\mathrm{der}\), необходимо будет проверить по меньшей мере следующее:
\begin{enumerate}
    \item можно ли определить пространство допустимых наблюдаемых \(\mathfrak{O}\) исключительно через редукции \(\mathfrak{S}\);
    \item можно ли задать оператор \(\mathrm{Cls}_{\mathfrak{K}}\) без привлечения новых фундаментальных сущностей;
    \item совпадает ли полученная классификация с теми различиями, которые трактат уже использует при анализе классов наблюдаемых.
\end{enumerate}

Если ответы на все три вопроса окажутся положительными, статус \(\mathrm{border}\) должен будет понизиться до \(\mathrm{der}\). Если же хотя бы один из них потребует нового несущего слоя, тогда семейство \(\mathfrak{O}_{cls}\) придётся признать более фундаментальным, чем предполагалось.

\subsection{Промежуточный итог второго теста}

Таким образом, второй тест программы редукционного сжатия даёт более тонкий результат, чем первый. Для \((\Xi,\Upsilon)\) удалось показать прямую схему реконструкции. Для \(\mathfrak{O}_{cls}\) удалось показать лишь частичную реконструируемость: классы наблюдаемых явно несут на себе отпечаток действия ядра, но ещё не доказано, что этим действием они исчерпываются.

Именно это различие и делает текущую стадию трактата зрелее. Теория учится не только подтверждать собственные гипотезы о производности, но и честно фиксировать случаи, где реконструкция пока остаётся неполной. Это уже не риторика цельности, а дисциплина различения.

\subsection{Следующая задача после пограничного случая}

После анализа \(\mathfrak{O}_{cls}\) следующая естественная цель состоит в том, чтобы проверить последний крупный кандидат на производность --- семейство локальных иерархий \(\mathfrak{H}_{mod}\). Там вопрос будет звучать ещё конкретнее: является ли локальное ранжирование моделей всего лишь частной разверткой глобального функционала \(\mathcal{Q}_{sel}\), или же в нём скрыт отдельный несущий механизм сравнения.

Если этот шаг удастся выполнить, то трактат получит уже почти полную первую карту доказанных, пограничных и ещё неразрешённых статусов. Именно тогда программа редукционного сжатия сможет претендовать на форму внутренней теории, а не только методологического обещания.

\section{Локальные иерархии моделей как развёртка глобальной селекции}

После реконструкции \((\Xi,\Upsilon)\) и анализа пограничного статуса \(\mathfrak{O}_{cls}\) остаётся проверить последний крупный элемент первой статусной карты --- семейство локальных иерархий моделей \(\mathfrak{H}_{mod}\). Именно здесь вопрос о производности принимает максимально конкретную форму. Нужно понять, является ли локальное ранжирование моделей самостоятельным слоем сравнения или же оно возникает как частный срез уже введённого глобального функционала \(\mathcal{Q}_{sel}\).

Если удастся показать второе, то первый контур программы редукционного сжатия получит почти завершённый вид: один блок будет реконструирован как производный, один --- подтверждён как пограничный, и ещё один --- сведен к глобальному ядру сравнения.

\subsection{Что следует понимать под \texorpdfstring{$\mathfrak{H}_{mod}$}{Hmod}}

Под \(\mathfrak{H}_{mod}\) будем понимать не просто множество моделей, а совокупность их локальных порядков, возникающих при сравнении на отдельных классах наблюдаемых, отдельных режимах или отдельных исследовательских задачах. Иначе говоря, \(\mathfrak{H}_{mod}\) фиксирует не абсолютную истину о моделях, а их относительное расположение при конкретном выборе контекста сравнения.

Именно это и делает данный объект удобным кандидатом в производный класс. Если локальный порядок всякий раз определяется выбором класса \(\mathcal{C}_k\), режима чтения и меры качества, то можно ожидать, что сама структура локальных иерархий не требует нового фундаментального механизма, а строится из уже имеющегося аппарата отбора.

\subsection{Схема вывода локального ранжирования}

Пусть для каждой модели \(M\) и каждого класса наблюдаемых \(\mathcal{C}_k\) глобальный селекционный функционал индуцирует локальное значение качества
\begin{equation}
\mathcal{Q}_k(M)=\mathcal{Q}_{sel}(M\mid \mathcal{C}_k).
\end{equation}
Тогда локальная иерархия на классе \(\mathcal{C}_k\) естественно определяется как порядок
\begin{equation}
M_i \succ_k M_j \quad \Longleftrightarrow \quad \mathcal{Q}_k(M_i)>\mathcal{Q}_k(M_j).
\end{equation}

Если теперь собрать все такие порядки по всем допустимым классам и режимам, то получим семейство
\begin{equation}
\mathfrak{H}_{mod}=\bigl\{\succ_k\bigr\}_{k\in K},
\end{equation}
где множество индексов \(K\) обозначает все контексты локального сравнения, признанные трактатом существенными.

В этой записи уже видно, что \(\mathfrak{H}_{mod}\) возникает как развертка одного и того же селекционного принципа на множество частных условий. Следовательно, в первом приближении семейство локальных иерархий действительно может быть понято как производный объект.

\subsection{Почему новый механизм сравнения не обязателен}

Возникает естественное возражение: не скрывается ли в локальных иерархиях нечто большее, чем просто частные значения \(\mathcal{Q}_{sel}\)? На первый взгляд может показаться, что локальный порядок обладает собственной логикой, поскольку разные классы наблюдаемых могут менять относительные позиции моделей. Но это ещё не означает существования нового фундаментального механизма.

Напротив, сама зависимость порядка от класса \(\mathcal{C}_k\) как раз и говорит в пользу производности. Если глобальный функционал уже допускает контекстную оценку качества, то изменчивость локальных порядков является не признаком нового слоя, а естественным следствием применения одного и того же принципа к разным структурным ситуациям.

Иначе говоря, \(\mathfrak{H}_{mod}\) не добавляет к теории новую способность сравнивать; оно лишь распределяет уже имеющуюся способность по множеству локальных контекстов.

\subsection{Формула производности для \texorpdfstring{$\mathfrak{H}_{mod}$}{Hmod}}

Это позволяет записать компактную формулу реконструкции:
\begin{equation}
\mathfrak{H}_{mod}=\mathrm{LocOrd}(\mathcal{Q}_{sel},\mathfrak{O}_{cls}),
\end{equation}
где оператор \(\mathrm{LocOrd}\) обозначает семейство локальных порядков, порождённых ограничением глобального функционала качества на различные классы наблюдаемых.

Если же в качестве промежуточной структуры использовать уже введённую координатно-селекционную развёртку, то можно записать более детальную форму:
\begin{equation}
\mathfrak{H}_{mod}=\mathrm{LocOrd}(\Upsilon,\mathfrak{O}_{cls}),
\end{equation}
где \(\Upsilon\) есть не самостоятельная сущность, а селекционный слой, индуцированный \(\mathcal{Q}_{sel}\). В обоих вариантах семейство \(\mathfrak{H}_{mod}\) остаётся производным относительно уже выделенного ядра.

\subsection{Где остаётся тонкость}

Тем не менее и здесь есть одна важная тонкость. Чтобы локальные иерархии действительно были полностью выводимы, необходимо, чтобы трактат заранее имел достаточно определённый набор классов сравнения и правило допустимого ограничения функционала \(\mathcal{Q}_{sel}\). Иначе говоря, производность \(\mathfrak{H}_{mod}\) частично зависит от того, насколько устойчиво задан объект \(\mathfrak{O}_{cls}\).

Но в отличие от самого \(\mathfrak{O}_{cls}\), эта зависимость не выглядит онтологической. Она скорее техническая и организационная. Локальная иерархия не создаёт новый принцип различения; она лишь наследует недоопределённость того пространства, на котором производится локальное сравнение. Поэтому даже при неполной ясности в отношении \(\mathfrak{O}_{cls}\) статус \(\mathfrak{H}_{mod}\) всё ещё разумно сохранять как производный.

\subsection{Критерий успешной редукции}

Чтобы окончательно удержать этот вывод, достаточно потребовать двух условий:
\begin{enumerate}
    \item любое локальное ранжирование должно восстанавливаться из ограничения одного и того же глобального критерия качества;
    \item при переходе между разными классами наблюдаемых не должен появляться новый фундаментальный параметр сравнения, отсутствующий в \(\mathcal{Q}_{sel}\).
\end{enumerate}

Если оба условия соблюдены, то \(\mathfrak{H}_{mod}\) не имеет права претендовать на фундаментальный статус. Оно оказывается лишь множественным представлением уже имеющегося селекционного закона.

\subsection{Итог по статусу локальных иерархий}

Следовательно, третий тест программы редукционного сжатия даёт результат, близкий к положительному завершению первой карты. Семейство \(\mathfrak{H}_{mod}\) можно понимать как организованную развертку глобального функционала \(\mathcal{Q}_{sel}\) по различным классам наблюдаемых. Оно не вводит новый тип сущности, а лишь эксплицирует то, как один и тот же принцип качества ведёт себя в разных локальных режимах сравнения.

Тем самым производный статус \(\mathfrak{H}_{mod}\) получает достаточное методологическое основание. В отличие от \(\mathfrak{O}_{cls}\), здесь не возникает сильного подозрения на наличие самостоятельного несущего слоя.

\subsection{Промежуточное закрытие первой карты}

Теперь можно подвести промежуточный итог всей первой фазы редукционного сжатия. На текущем этапе трактат располагает следующей более дисциплинированной картиной:
\begin{itemize}
    \item \(\mathfrak{S}\), \(\mathcal{I}(\mathfrak{S})\), \(\mathcal{R}_{type}\) и \(\mathcal{Q}_{sel}\) остаются ядровыми кандидатами;
    \item \((\Xi,\Upsilon)\) и \(\mathfrak{H}_{mod}\) получают осмысленную схему производной реконструкции;
    \item \(\mathfrak{O}_{cls}\) сохраняет пограничный статус как главный ещё неразрешённый случай.
\end{itemize}

Это уже существенно более сильная позиция, чем в начале трактата. Теория больше не просто перечисляет свои элементы, а умеет различать, какие из них действительно несут архитектуру, какие являются развертками, а какие образуют область настоящей нерешённой проблемы.

\subsection{Следующий рубеж}

После такого промежуточного закрытия первый цикл самосжатия можно считать состоявшимся. Следующий рубеж должен состоять уже не просто в дальнейших классификациях, а в переходе к более жёсткой форме: либо к почти-аксиоматической записи минимального ядра, либо к списку лемм о сохранении структуры при редукциях и деформациях.

Именно там трактат сможет сделать следующий качественный шаг --- от дисциплинированной карты статусов к собственному короткому формальному костяку.

\section{Почти-аксиоматический каркас минимального ядра}

После завершения первого цикла редукционного сжатия у трактата появляется возможность сделать следующий шаг к компактной форме. Если раньше текст в основном строил карту различий и проверял статус отдельных элементов, то теперь он может попробовать собрать минимальный набор утверждений, без которых вся найденная архитектура перестаёт работать. Такой набор ещё не является окончательной аксиоматикой, но уже может служить коротким формальным каркасом.

Под \textit{почти-аксиоматическим каркасом} будем понимать систему утверждений, которые не претендуют на полную математическую завершённость, но уже достаточно сжаты и жёстки, чтобы играть роль внутреннего основания для дальнейших выводов.

\subsection{Аксома единой сущности}

\textbf{Аксома A1.} Существует единая сущность \(\mathfrak{S}\), допускающая несколько режимов чтения, которые не являются чтениями разных сущностей, а относятся к одному и тому же источнику.

Смысл этой аксиомы не в простом введении нового имени. Она фиксирует главное онтологическое требование трактата: различие языков теории не должно автоматически означать различие предметов. Без этого пункта вся дальнейшая программа редукций и сопоставлений распадается.

\subsection{Аксома инвариантной идентичности}

\textbf{Аксома A2.} Сущности \(\mathfrak{S}\) сопоставлено инвариантное ядро \(\mathcal{I}(\mathfrak{S})\), сохраняющее идентичность при всех допустимых слабых деформациях.

Эта аксиома отделяет сущность как носитель единства от сущности как носитель устойчивости. Если \(\mathfrak{S}\) утверждает существование единого источника, то \(\mathcal{I}(\mathfrak{S})\) утверждает, что в этом источнике есть нечто неустранимое и структурно сохраняемое.

\subsection{Аксома режимной различимости}

\textbf{Аксома A3.} Существует минимальная типология редукций \(\mathcal{R}_{type}=\{R_{\mathcal{B}},R_{\Phi},R_{mix}\}\), достаточная для различения каркасного, фазового и смешанного проявлений одной и той же сущности.

Тем самым трактат утверждает, что по крайней мере три типа проявления должны входить в минимальный аппарат. Не потому, что это окончательно исчерпывает мир, а потому, что именно без такого различения уже теряется способность организованно читать накопленный материал.

\subsection{Аксома селекционной нетривиальности}

\textbf{Аксома A4.} На множестве допустимых реализаций и моделей задан глобальный функционал качества \(\mathcal{Q}_{sel}\), позволяющий различать конкурирующие реализации не риторически, а структурно.

Эта аксиома принципиальна. Без неё трактат мог бы оставаться интересной онтологией, но ещё не становился бы программой рационального отбора. Именно наличие \(\mathcal{Q}_{sel}\) переводит схему из состояния описания в состояние сравнения.

\subsection{Сжатая запись ядра}

Тогда минимальное ядро теории можно записать в форме
\begin{equation}
\mathfrak{K}_{min}=\bigl(\mathfrak{S},\mathcal{I}(\mathfrak{S}),\mathcal{R}_{type},\mathcal{Q}_{sel}\bigr).
\end{equation}

Эта запись теперь уже не просто рабочая эвристика, а итог проведённого сжатия. Она собирает в одной строке именно те структуры, которые на текущей стадии трактата удерживают единство предмета, инвариантную идентичность, режимную различимость и селекционную нетривиальность.

\subsection{Первая лемма о производных развёртках}

\textbf{Лемма L1.} Если заданы \(\mathcal{R}_{type}\) и \(\mathcal{Q}_{sel}\), то координатно-селекционная пара \((\Xi,\Upsilon)\) может быть реконструирована как производная развёртка ядра.

Содержательно это выражает уже установленный ранее результат:
\begin{equation}
(\Xi,\Upsilon)=\mathfrak{Rec}\bigl(\mathcal{R}_{type},\mathcal{Q}_{sel}\bigr).
\end{equation}

Лемма важна тем, что впервые показывает: теория уже умеет не только перечислять производные элементы, но и явно восстанавливать их из более короткого аппарата.

\subsection{Вторая лемма о локальных порядках}

\textbf{Лемма L2.} Если задан глобальный функционал качества \(\mathcal{Q}_{sel}\) и пространство допустимых контекстов сравнения, то локальные иерархии моделей \(\mathfrak{H}_{mod}\) являются производной локализацией одного и того же селекционного закона.

Схематически это выражается формулой
\begin{equation}
\mathfrak{H}_{mod}=\mathrm{LocOrd}(\mathcal{Q}_{sel},\mathfrak{O}_{cls}).
\end{equation}

Тем самым локальное ранжирование моделей лишается статуса самостоятельного несущего механизма и переводится в класс организованных производных представлений.

\subsection{Третья лемма о пограничном остатке}

\textbf{Лемма L3.} Семейство классов наблюдаемых \(\mathfrak{O}_{cls}\) на текущем этапе зависит от действия ядра, но ещё не доказано, что оно без остатка выводимо из \(\mathfrak{K}_{min}\).

В краткой форме это можно зафиксировать так:
\begin{equation}
\mathfrak{O}_{cls}\subseteq \mathrm{Im}\,\mathrm{Cls}_{\mathfrak{K}},
\end{equation}
но равенство здесь пока не установлено. Эта лемма фиксирует важный тип знания: не доказанную производность и не доказанную фундаментальность, а дисциплинированно удержанную пограничность.

\subsection{Первая теорема-схема текущего этапа}

На основе аксиом и лемм можно сформулировать и первую общую теорему-схему трактата.

\textbf{Теорема-схема T1.} На текущем этапе минимальное ядро теории достаточно для реконструкции части расширенного аппарата и для объяснения того, почему другая часть остаётся пограничной.

Иначе говоря, существует разложение вида
\begin{equation}
\mathfrak{A}_{ext}=\mathfrak{K}_{min} \oplus \mathfrak{D}_{rec} \oplus \mathfrak{B}_{open},
\end{equation}
где \(\mathfrak{D}_{rec}\) обозначает уже реконструируемые производные структуры, а \(\mathfrak{B}_{open}\) --- ещё не закрытые пограничные элементы.

На текущей стадии трактата естественно считать, что
\begin{equation}
\mathfrak{D}_{rec}\supset \{(\Xi,\Upsilon),\mathfrak{H}_{mod}\},
\qquad
\mathfrak{B}_{open}\supset \{\mathfrak{O}_{cls}\}.
\end{equation}

Эта теорема-схема важна не как окончательное доказательство, а как компактное выражение достигнутого состояния программы.

\subsection{Что этот каркас уже позволяет делать}

У предложенного каркаса есть непосредственная польза. Он позволяет:
\begin{itemize}
    \item оценивать новые элементы аппарата по отношению к уже сжатому ядру;
    \item различать доказанную производность и честно удержанную пограничность;
    \item строить дальнейшие разделы трактата не от интуитивного расширения, а от короткого списка структурных обязательств.
\end{itemize}

Именно это и означает переход к настоящему формальному костяку. Теория получает не просто красивую карту мотивов, а внутреннюю норму того, какие утверждения должны считаться исходными, какие --- выводимыми, а какие --- открытыми.

\subsection{Следующий шаг после каркаса}

Теперь следующий естественный рубеж уже почти предопределён. После фиксации почти-аксиоматического каркаса трактат может двигаться в одном из двух направлений. Либо он начнёт доказывать более строгие леммы о сохранении структуры при деформациях и редукциях, либо перейдёт к построению почти-формальной динамики на ядре \(\mathfrak{K}_{min}\).

В обоих случаях ситуация уже отличается от исходной философской стадии. Теперь текст располагает коротким формальным минимумом, от которого можно двигаться дальше не только интуитивно, но и структурно-дедуктивно.

\section{Леммы о сохранении структуры при деформациях и редукциях}

После фиксации почти-аксиоматического каркаса следующий шаг естественно состоит в том, чтобы уточнить, как именно ядро теории ведёт себя при допустимых изменениях. Без этого даже короткий каркас оставался бы статическим описанием. Между тем сама идея трактата требует большего: нужно показать, что некоторые структуры переживают деформации и переходы между режимами чтения, а другие, напротив, могут меняться, не разрушая общей идентичности.

Именно поэтому здесь уместно сформулировать первые леммы о сохранении. Они ещё не являются строгими математическими теоремами в полном смысле, но уже задают тот тип внутренней дисциплины, без которого минимальное ядро не сможет работать как настоящее основание теории.

\subsection{Лемма о сохранении инвариантного ядра при слабых деформациях}

\textbf{Лемма S1.} Если две реализации \(\mathfrak{S}_1\) и \(\mathfrak{S}_2\) связаны слабой допустимой деформацией, то их инвариантные ядра совпадают в структурно существенной части:
\begin{equation}
\mathfrak{S}_1 \sim_w \mathfrak{S}_2 \,\Longrightarrow\, \mathcal{I}(\mathfrak{S}_1)\cong \mathcal{I}(\mathfrak{S}_2).
\end{equation}

Смысл этой леммы состоит в том, что слабая деформация не может считаться внутренней вариацией одной и той же сущности, если она разрушает именно то, что теория объявила неустранимым. Следовательно, инвариантное ядро должно быть устойчиво не по риторическому соглашению, а по самой логике различения сущностей.

\subsection{Лемма о сохранении типологии редукций}

\textbf{Лемма S2.} При слабых и промежуточных деформациях минимальная типология редукций \(\mathcal{R}_{type}\) сохраняется как типология, даже если меняются относительные веса отдельных режимов.

Это можно выразить схематически так:
\begin{equation}
\mathfrak{S}_1 \sim_{w/int} \mathfrak{S}_2
\,\Longrightarrow\,
\mathcal{R}_{type}(\mathfrak{S}_1)\simeq \mathcal{R}_{type}(\mathfrak{S}_2).
\end{equation}

Данная лемма особенно важна для всей программы чтений. Она утверждает, что допустимые деформации могут смещать объект ближе к каркасному, фазовому или смешанному режиму, но не обязаны каждый раз порождать новый набор фундаментальных типов проявления.

\subsection{Лемма о ковариантности производных развёрток}

\textbf{Лемма S3.} Если ядро \(\mathfrak{K}_{min}\) сохраняется под допустимой деформацией, то его производные развёртки изменяются ковариантно, а не произвольно.

В частности, для уже реконструированных объектов это означает:
\begin{equation}
\mathfrak{K}_{min}^{(1)} \rightsquigarrow \mathfrak{K}_{min}^{(2)}
\,\Longrightarrow\,
(\Xi,\Upsilon)^{(1)} \rightsquigarrow (\Xi,\Upsilon)^{(2)},
\quad
\mathfrak{H}_{mod}^{(1)} \rightsquigarrow \mathfrak{H}_{mod}^{(2)}.
\end{equation}

Иначе говоря, производные объекты не должны иметь собственной независимой динамики, оторванной от динамики ядра. Их допустимое изменение обязано быть наведённым, а не фундаментально автономным.

\subsection{Лемма о селекционной устойчивости}

\textbf{Лемма S4.} Если деформация не разрушает \(\mathcal{I}(\mathfrak{S})\) и не выводит реализацию за пределы допустимой типологии \(\mathcal{R}_{type}\), то глобальный функционал качества \(\mathcal{Q}_{sel}\) сохраняет свою различающую способность.

Содержательно это означает, что допустимые деформации могут менять значения качества, но не должны превращать функционал отбора в тривиальный. Теория обязана сохранять возможность различать более и менее структурно удачные реализации даже после контролируемого изменения условий.

Эту мысль удобно записать как требование нетривиальности:
\begin{equation}
\exists M_i,M_j:
\quad
\mathcal{Q}_{sel}^{(deformed)}(M_i)\neq \mathcal{Q}_{sel}^{(deformed)}(M_j).
\end{equation}

Если это условие систематически нарушается, значит, либо деформация вышла за пределы допустимой области, либо сама схема селекции задана слишком слабо.

\subsection{Лемма о границе сильной деформации}

\textbf{Лемма S5.} Если деформация разрушает инвариантное ядро, стирает типологию редукций или устраняет селекционную нетривиальность, то такую деформацию следует считать сильной и выводящей объект за пределы прежней сущности.

В сокращённой форме:
\begin{equation}
\neg \mathcal{I}\text{-stable}
\;\vee\;
\neg \mathcal{R}_{type}\text{-stable}
\;\vee\;
\neg \mathcal{Q}_{sel}\text{-nontrivial}
\;\Longrightarrow\;
\mathfrak{S}_1 \not\sim \mathfrak{S}_2.
\end{equation}

Эта лемма важна, потому что она фиксирует не только то, что теория сохраняет, но и то, где именно проходит граница допустимого. Без такой границы язык деформаций быстро стал бы слишком эластичным и перестал бы различать изменения внутри сущности и переходы к другой сущности.

\subsection{Редукции как контролируемые проекции}

Из перечисленных лемм вытекает и общий принцип работы редукций. Геометрическая, спектральная и корреляционная редукции не должны пониматься как произвольные потери информации. Они должны работать как контролируемые проекции, которые:
\begin{itemize}
    \item сохраняют хотя бы часть инвариантного ядра;
    \item не разрушают типологию режимов;
    \item допускают обратное сопоставление с исходной сущностью хотя бы на уровне структурной идентификации;
\end{itemize}

В этом смысле корректная редукция есть не простое ``схлопывание'', а дисциплинированный переход, при котором сущность становится читаемой в более узком окне, но не исчезает как тот же самый источник.

\subsection{Первая схема сохранения}

Теперь можно записать и первую обобщённую схему сохранения:
\begin{equation}
\mathfrak{S}
\xrightarrow{\text{допустимая деформация}}
\mathfrak{S}'
\xrightarrow{\mathcal{R}_*}
\mathcal{R}_*(\mathfrak{S}')
\end{equation}
должна быть устроена так, чтобы сохранялись:
\begin{equation}
\mathcal{I}(\mathfrak{S}) \leadsto \mathcal{I}(\mathfrak{S}'),
\qquad
\mathcal{R}_{type}(\mathfrak{S}) \leadsto \mathcal{R}_{type}(\mathfrak{S}'),
\qquad
\mathcal{Q}_{sel} \leadsto \mathcal{Q}_{sel}'.
\end{equation}

Знак \(\leadsto\) здесь обозначает не тождество значений, а контролируемую преемственность структуры. Именно такая форма записи удобна для дальнейших, уже более строгих версий лемм.

\subsection{Что эти леммы дают трактату}

У этих лемм есть важная функция. Они переводят разговор о минимальном ядре из состояния ``что есть'' в состояние ``что обязано сохраняться''. Тем самым трактат впервые получает не только статическую архитектуру, но и зародыш собственной структурной кинематики.

Это особенно существенно для дальнейшего роста текста. Пока ядро задано лишь как список исходных структур, теория остаётся описательной. Как только появляются правила сохранения и правила нарушения, она начинает приобретать форму дисциплинированной внутренней механики.

\subsection{Следующий рубеж после лемм сохранения}

После этого следующий шаг становится ещё яснее. Теперь можно либо уточнять каждую из этих лемм до почти-формального доказательства, либо строить поверх них динамический принцип: как именно эволюционируют допустимые реализации \(\mathfrak{S}\) внутри пространства, где инвариантное ядро сохраняется, а производные слои перестраиваются.

Именно этот переход --- от аксиоматики и лемм сохранения к почти-формальной динамике --- и может открыть следующую крупную фазу трактата.

\section{Почти-формальная динамика на ядре \texorpdfstring{$\mathfrak{K}_{min}$}{Kmin}}

После аксиоматики и лемм сохранения трактат может сделать следующий качественный шаг: перейти от вопроса ``что сохраняется?'' к вопросу ``как происходит допустимое изменение?'' Если ядро \(\mathfrak{K}_{min}\) действительно играет роль минимального основания, то теория должна уметь описывать не только статические отношения между его компонентами, но и формы их контролируемой эволюции.

Речь пока не идёт о полной динамической теории в привычном физическом смысле. Здесь достаточно построить \textit{почти-формальную динамику}: схему, в которой допустимые реализации \(\mathfrak{S}\) движутся по пространству состояний так, что инвариантное ядро сохраняется, типология редукций не разрушается, а производные слои перестраиваются согласованно.

\subsection{Пространство допустимых реализаций}

Обозначим через \(\mathfrak{M}_{adm}\) пространство всех допустимых реализаций сущности \(\mathfrak{S}\), совместимых с аксиомами \(A1\)--\(A4\). Тогда допустимая динамика должна быть задана не на всём мыслимом пространстве конфигураций, а только на той его части, где сохраняется принадлежность к одной и той же сущности.

Схематически это можно записать как семейство траекторий
\begin{equation}
\gamma: \tau \longmapsto \mathfrak{S}(\tau),
\qquad
\mathfrak{S}(\tau)\in \mathfrak{M}_{adm},
\end{equation}
где параметр \(\tau\) пока не обязан интерпретироваться как физическое время. Он обозначает скорее внутренний параметр деформации или эволюции внутри пространства допустимых реализаций.

\subsection{Условие допустимой динамики}

Главное требование к такой траектории состоит в том, что она не должна разрушать минимальное ядро. Поэтому будем называть траекторию \(\gamma\) допустимой, если для всех значений параметра выполняются три условия:
\begin{enumerate}
    \item сохраняется структурная идентичность \(\mathcal{I}(\mathfrak{S}(\tau))\);
    \item не исчезает типология редукций \(\mathcal{R}_{type}(\mathfrak{S}(\tau))\);
    \item функционал \(\mathcal{Q}_{sel}^{(\tau)}\) остаётся нетривиальным.
\end{enumerate}

Иначе говоря, допустимая динамика есть такая эволюция, при которой движение происходит \textit{внутри} сущности, а не через её распад.

\subsection{Динамика ядра и динамика производных слоёв}

Теперь важно различить два уровня изменения. Первый уровень относится к самому ядру:
\begin{equation}
\mathfrak{K}_{min}(\tau)=\bigl(\mathfrak{S}(\tau),\mathcal{I}(\mathfrak{S}(\tau)),\mathcal{R}_{type}(\tau),\mathcal{Q}_{sel}^{(\tau)}\bigr).
\end{equation}

Второй уровень относится к производным структурам, которые должны следовать за ядром:
\begin{equation}
(\Xi,\Upsilon)(\tau)=\mathfrak{Rec}\bigl(\mathcal{R}_{type}(\tau),\mathcal{Q}_{sel}^{(\tau)}\bigr),
\qquad
\mathfrak{H}_{mod}(\tau)=\mathrm{LocOrd}\bigl(\mathcal{Q}_{sel}^{(\tau)},\mathfrak{O}_{cls}(\tau)\bigr).
\end{equation}

Тем самым динамика устроена послойно: ядро эволюционирует в допустимой области, а производные объекты перестраиваются как наведённые тени этого движения.

\subsection{Квазиградиентный принцип}

Чтобы сделать эту схему содержательнее, полезно ввести первый пробный принцип движения. Наиболее естественно предположить, что допустимая динамика тяготеет к росту структурной согласованности, то есть к таким направлениям в пространстве \(\mathfrak{M}_{adm}\), где глобальный функционал качества не убывает слишком резко. Это можно выразить как квазиградиентное условие:
\begin{equation}
\frac{d}{d\tau}\mathcal{Q}_{sel}\bigl(\mathfrak{S}(\tau)\bigr) \gtrsim 0,
\end{equation}
при дополнительном ограничении, что инвариантное ядро сохраняется.

Смысл этого условия не в том, что всякая динамика обязана монотонно максимизировать качество. Оно выражает более мягкий тезис: допустимые направления эволюции не должны систематически вести к разрушению структурной организованности, если только теория не описывает переход к сильной деформации или распаду прежней сущности.

\subsection{Переключение между режимами как динамический эффект}

Внутри такой схемы естественно понимать смену доминирующего режима чтения как динамический эффект. Если по мере изменения параметра \(\tau\) меняются относительные веса компонентов режимного профиля, то один и тот же объект может переходить, например, из преимущественно каркасной области в смешанную, а затем в преимущественно фазовую.

Это означает, что типология \(\mathcal{R}_{type}\) сохраняется, но режимная локализация объекта в пространстве \(\Xi\) может эволюционировать. Тем самым теория получает первый механизм объяснения того, как одна и та же сущность может предъявлять разные феноменологические профили без потери идентичности.

\subsection{Динамика пограничного слоя}

Особый интерес представляет поведение пограничного объекта \(\mathfrak{O}_{cls}(\tau)\). В отличие от \((\Xi,\Upsilon)\) и \(\mathfrak{H}_{mod}\), его эволюция не может считаться полностью наведённой без дополнительных оговорок. Здесь возможны два сценария.

В первом сценарии семейство классов наблюдаемых действительно индуцируется ядром и постепенно стабилизируется вместе с ним. Во втором сценарии для его определения требуется дополнительная операциональная структура, и тогда динамика \(\mathfrak{O}_{cls}(\tau)\) уже не сводится без остатка к динамике ядра. Именно это различие и делает пограничный слой столь важным для дальнейшей работы.

\subsection{Слабая динамика и сильный разрыв}

Теперь можно естественно различить два типа эволюции. \textbf{Слабой динамикой} назовём такую траекторию \(\gamma\), вдоль которой сохраняются условия допустимости, инвариантное ядро и селекционная нетривиальность. \textbf{Сильным разрывом} назовём такую точку \(\tau_*\), в которой по крайней мере одно из этих условий перестаёт выполняться.

Тем самым теория получает первый динамический аналог различия между слабой и сильной деформацией:
\begin{equation}
\gamma_{weak}: \tau\mapsto \mathfrak{S}(\tau)
\quad \text{vs.} \quad
\tau_*: \mathfrak{S}(\tau_*) \notin \mathfrak{M}_{adm}^{same}.
\end{equation}

Именно вблизи таких сильных разрывов, вероятно, и должны возникать переходы к другой сущности, другой режимной архитектуре или другой области теории.

\subsection{Первая динамическая схема трактата}

Всё это можно собрать в одну первую динамическую схему:
\begin{equation}
\mathfrak{K}_{min}(\tau)
\longrightarrow
\bigl((\Xi,\Upsilon)(\tau),\mathfrak{H}_{mod}(\tau),\mathfrak{O}_{cls}(\tau)\bigr),
\end{equation}
где:
\begin{itemize}
    \item ядро эволюционирует внутри допустимой области;
    \item производные слои следуют за ядром ковариантно;
    \item пограничный слой может либо стабилизироваться как производный, либо выявить собственную несводимость.
\end{itemize}

Эта схема важна тем, что впервые соединяет аксиомы, леммы сохранения и статусную карту в одном динамическом образе.

\subsection{Что меняется после введения динамики}

С введением почти-формальной динамики трактат начинает по-настоящему выходить за пределы расширенного введения. Теперь он описывает не только то, из чего теория состоит, но и то, как она должна мыслить допустимые изменения своих объектов. Это и есть первый признак внутренней зрелости: теория становится способной не просто перечислять структуры, а задавать правила их движения и границы их устойчивости.

\subsection{Следующий шаг после динамики}

После такого раздела следующий шаг становится ещё более конкретным. Теперь можно либо попробовать построить квазидействие или функционал, генерирующий допустимую динамику, либо перейти к конкретным примерам траекторий в пространстве реализаций \(\mathfrak{M}_{adm}\). В обоих случаях трактат уже располагает достаточным минимумом, чтобы двигаться от общей философии к почти-рабочему формализму.

\section{Квазидействие как генератор допустимой динамики}

Если предыдущий раздел описал, как должна выглядеть допустимая динамика, то теперь естественно поставить следующий вопрос: что именно её генерирует? Чтобы теория не оставалась на уровне словесной кинематики, ей нужен хотя бы первый пробный функционал, относительно которого допустимые траектории можно понимать как предпочтительные или стационарные.

Здесь ещё рано вводить полное действие в строгом физико-математическом смысле. Достаточно более слабого объекта, который будет играть роль \textit{квазидействия}: функционала, чувствительного к сохранению инвариантного ядра, режимной различимости и селекционной согласованности, но пока не требующего окончательной локальной лагранжевой формы.

\subsection{Общий вид квазидействия}

Обозначим такой функционал через \(\mathfrak{A}_{q}\). В первом приближении естественно требовать, чтобы он зависел от четырёх типов данных:
\begin{equation}
\mathfrak{A}_{q}=\mathfrak{A}_{q}\bigl[\mathfrak{S}(\tau),\mathcal{I}(\mathfrak{S}(\tau)),\mathcal{R}_{type}(\tau),\mathcal{Q}_{sel}^{(\tau)}\bigr].
\end{equation}

Эта запись выражает простую мысль: генератор допустимой динамики не может зависеть только от одной внешней координаты или только от одного качества. Он обязан одновременно учитывать идентичность сущности, устойчивое ядро, режимную архитектуру и структурную меру качества.

\subsection{Три вклада в функционал}

На текущей стадии разумно разложить квазидействие на три основных вклада:
\begin{equation}
\mathfrak{A}_{q}=\mathfrak{A}_{inv}+\mathfrak{A}_{mode}+\mathfrak{A}_{sel}.
\end{equation}

Здесь:
\begin{itemize}
    \item \(\mathfrak{A}_{inv}\) штрафует траектории, разрушающие инвариантное ядро;
    \item \(\mathfrak{A}_{mode}\) контролирует устойчивость режимной архитектуры и допустимость переходов между режимами;
    \item \(\mathfrak{A}_{sel}\) поощряет такие направления эволюции, в которых структурное качество не размывается и не теряет различающую силу.
\end{itemize}

Такое разложение важно тем, что оно отражает уже найденную архитектуру трактата. Теория не выдумывает новый словарь для динамики, а перерабатывает собственные уже выделенные несущие структуры в виде динамического функционала.

\subsection{Инвариантный вклад}

Первый вклад должен обеспечивать, чтобы допустимая траектория не уходила из класса одной и той же сущности. Поэтому в первом приближении его можно понимать как меру отклонения инвариантного ядра от устойчивой области:
\begin{equation}
\mathfrak{A}_{inv}\sim \int d\tau\; \mathcal{N}_{inv}\bigl(\mathcal{I}(\mathfrak{S}(\tau)),\mathcal{I}_0\bigr),
\end{equation}
где \(\mathcal{I}_0\) обозначает эталонный класс инвариантной идентичности, а \(\mathcal{N}_{inv}\) --- некоторую меру несогласованности.

Чем сильнее траектория стремится вывести реализацию за пределы прежнего инвариантного ядра, тем больше должен становиться этот вклад. В пределе сильной деформации он должен переставать быть контролируемым, сигнализируя о выходе из допустимой области.

\subsection{Режимный вклад}

Второй вклад должен учитывать, что допустимая динамика может менять относительные веса режимов, но не обязана разрушать саму типологию \(\mathcal{R}_{type}\). Поэтому естественно ввести функционал, чувствительный к плавности и допустимости смены режимного профиля:
\begin{equation}
\mathfrak{A}_{mode}\sim \int d\tau\; \mathcal{N}_{mode}\bigl(\Xi(\tau),\dot{\Xi}(\tau)\bigr),
\end{equation}
где \(\Xi(\tau)\) выступает как реконструированное пространство режимной локализации.

Этот вклад не должен запрещать переключения между каркасным, смешанным и фазовым доминированием. Но он должен различать контролируемый переход и хаотическую потерю режимной структуры. Иначе говоря, он регулирует не факт изменения, а его структурную дисциплину.

\subsection{Селекционный вклад}

Третий вклад должен связывать динамику с уже введённым квазиградиентным принципом. Его простейшая функция --- не дать эволюции регулярно скатываться в область, где \(\mathcal{Q}_{sel}\) теряет различающую способность или где структурная согласованность систематически убывает. Поэтому можно записать схематически:
\begin{equation}
\mathfrak{A}_{sel}\sim -\int d\tau\; W\bigl(\mathcal{Q}_{sel}^{(\tau)}\bigr),
\end{equation}
где функция \(W\) возрастает с ростом структурного качества.

Знак минус здесь выражает простую идею: более качественные траектории должны быть динамически предпочтительнее, то есть соответствовать меньшему значению квазидействия.

\subsection{Принцип стационарности}

Теперь можно сформулировать и первый динамический принцип трактата:
\begin{quote}
допустимыми траекториями сущности следует считать такие траектории в \(\mathfrak{M}_{adm}\), которые являются стационарными или квазистационарными для функционала \(\mathfrak{A}_{q}\) при сохранении условий допустимости.
\end{quote}

Символически это можно записать как условие
\begin{equation}
\delta \mathfrak{A}_{q}=0
\qquad
\text{при}
\qquad
\delta \mathcal{I}=0 \text{ в допустимом смысле.}
\end{equation}

Это ещё не полноценное уравнение движения, но уже вполне ясная схема того, как теория хочет мыслить предпочтительные эволюции.

\subsection{Стационарные, метастабильные и разрывные траектории}

После введения квазидействия полезно различить три типа траекторий.

\textbf{Стационарные траектории} соответствуют устойчивым минимумам или плато функционала \(\mathfrak{A}_{q}\). Они выражают особенно устойчивые формы организации сущности.

\textbf{Метастабильные траектории} соответствуют локально устойчивым режимам, которые сохраняются в ограничённой области параметров, но могут перейти в другой режим при накоплении деформации.

\textbf{Разрывные траектории} возникают тогда, когда вариационный принцип перестаёт удерживать инвариантное ядро или режимную типологию. Такие траектории естественно интерпретировать как подход к сильной деформации, фазовому перелому или смене области теории.

\subsection{Квазидействие и статусная карта}

Важно подчеркнуть, что введение квазидействия не отменяет статусную карту, а наоборот, делает её динамически осмысленной. Ядровые элементы входят в \(\mathfrak{A}_{q}\) как несущие аргументы. Производные структуры возникают как наведённые конфигурации вдоль допустимых траекторий. Пограничный слой \(\mathfrak{O}_{cls}\) становится местом, где динамика может либо закрыть нерешённость, либо подтвердить её принципиальный характер.

Тем самым динамический функционал превращает статическую карту в рабочий механизм отбора: уже не просто ``что есть фундаментальное?'', а ``какие эволюции удерживают фундаментальное и какие режимы они порождают?''

\subsection{Первая компактная динамическая запись}

Теперь можно собрать всё в одну короткую схему:
\begin{equation}
\mathfrak{A}_{q}[\mathfrak{K}_{min}] 
\quad \Longrightarrow \quad
\delta \mathfrak{A}_{q}=0
\quad \Longrightarrow \quad
\mathfrak{K}_{min}(\tau)
\quad \Longrightarrow \quad
\bigl((\Xi,\Upsilon)(\tau),\mathfrak{H}_{mod}(\tau),\mathfrak{O}_{cls}(\tau)\bigr).
\end{equation}

Эта запись важна не своей математической окончательностью, а тем, что она впервые собирает в одной цепочке основание, вариационный принцип и производные структуры.

\subsection{Следующий шаг после квазидействия}

После этого трактат получает уже совсем конкретное продолжение. Следующий шаг может состоять либо в построении пробных форм для \(\mathcal{N}_{inv}\), \(\mathcal{N}_{mode}\) и \(W\), либо в разборе нескольких модельных траекторий --- устойчивой, метастабильной и разрывной. В обоих случаях текст начинает двигаться не просто к философской ясности, а к зарождающемуся рабочему формализму.

\section{Пробные формы квазидействия}

Чтобы квазидействие \(\mathfrak{A}_{q}\) перестало быть только структурной схемой, полезно задать для его основных вкладов первые пробные формы. Речь пока не идёт о конечной и единственно верной записи. На этой стадии важнее другое: показать, что даже очень простые функциональные формы уже позволяют различать устойчивые, метастабильные и разрывные траектории и тем самым переводят динамику трактата в полуколичественный режим.

Именно поэтому ниже будут введены не окончательные законы, а \textit{рабочие модели} для трёх составляющих квазидействия. Их задача --- не закрыть формализм, а создать первую площадку, на которой разные типы эволюции можно сравнивать по величине вклада, штрафа и выигрыша в структурной согласованности.

\subsection{Пробная форма инвариантного штрафа}

Для \(\mathfrak{A}_{inv}\) естественно взять форму, которая обращается в минимум, когда инвариантное ядро остаётся внутри эталонного класса, и быстро растёт при удалении от него. В простейшей схеме можно записать:
\begin{equation}
\mathcal{N}_{inv}=\alpha_{inv}\, d_{inv}\bigl(\mathcal{I}(\mathfrak{S}(\tau)),\mathcal{I}_0\bigr)^2,
\end{equation}
где \(d_{inv}\) обозначает меру расстояния до эталонного класса идентичности, а \(\alpha_{inv}>0\) --- коэффициент жёсткости сохранения ядра.

Тогда соответствующий вклад принимает вид
\begin{equation}
\mathfrak{A}_{inv}=\alpha_{inv}\int d\tau\; d_{inv}\bigl(\mathcal{I}(\mathfrak{S}(\tau)),\mathcal{I}_0\bigr)^2.
\end{equation}

Эта форма удобна тем, что она сразу различает три режима: малые флуктуации вокруг инвариантного класса, контролируемый дрейф и резкий уход в область разрушения идентичности.

\subsection{Пробная форма режимного вклада}

Для режимного слоя \(\Xi(\tau)\) естественно использовать комбинацию потенциала и кинетического штрафа. Тогда можно написать:
\begin{equation}
\mathcal{N}_{mode}=\alpha_{m}\,\|\dot{\Xi}(\tau)\|^2 + V_{mode}(\Xi(\tau)),
\end{equation}
где \(V_{mode}\) --- потенциальный профиль режимной локализации.

В простейшей картине потенциал можно мыслить как имеющий несколько долин, отвечающих каркасному, смешанному и фазовому доминированию:
\begin{equation}
V_{mode}(\Xi)=\beta_{\mathcal{B}}U_{\mathcal{B}}(\Xi)+\beta_{mix}U_{mix}(\Xi)+\beta_{\Phi}U_{\Phi}(\Xi).
\end{equation}

Тогда:
\begin{itemize}
    \item устойчивая траектория сидит в одной долине;
    \item метастабильная траектория медленно движется по пологому барьеру между долинами;
    \item разрывная траектория фактически выходит из области, где такой потенциал ещё описывает допустимый режимный профиль.
\end{itemize}

Соответствующий вклад в квазидействие принимает вид
\begin{equation}
\mathfrak{A}_{mode}=\int d\tau\;\Bigl[\alpha_{m}\,\|\dot{\Xi}(\tau)\|^2 + V_{mode}(\Xi(\tau))\Bigr].
\end{equation}

\subsection{Пробная форма селекционного выигрыша}

Для третьего вклада нужен функционал, который награждает траектории, сохраняющие высокую структурную согласованность, но не становится сингулярным при умеренных колебаниях качества. В качестве простейшей формы удобно взять линейно-квадратичную комбинацию:
\begin{equation}
W(\mathcal{Q}_{sel})=\lambda_1\mathcal{Q}_{sel}-\lambda_2\mathcal{Q}_{sel}^2,
\qquad
\lambda_1,\lambda_2>0.
\end{equation}

Тогда селекционный вклад записывается как
\begin{equation}
\mathfrak{A}_{sel}=-\int d\tau\;\Bigl(\lambda_1\mathcal{Q}_{sel}^{(\tau)}-\lambda_2(\mathcal{Q}_{sel}^{(\tau)})^2\Bigr).
\end{equation}

Такая запись полезна тем, что она не просто предпочитает большие значения \(\mathcal{Q}_{sel}\), но и допускает существование устойчивых диапазонов качества, а не только неограниченного роста. Тем самым она лучше согласуется с идеей стационарных и метастабильных режимов.

\subsection{Итоговая пробная форма квазидействия}

Собирая три вклада вместе, получаем первую полуколичественную запись:
\begin{align}
\mathfrak{A}_{q}
= {} & \alpha_{inv}\int d\tau\; d_{inv}\bigl(\mathcal{I}(\mathfrak{S}(\tau)),\mathcal{I}_0\bigr)^2 \\
& + \int d\tau\;\Bigl[\alpha_{m}\,\|\dot{\Xi}(\tau)\|^2 + V_{mode}(\Xi(\tau))\Bigr] \\
& - \int d\tau\;\Bigl(\lambda_1\mathcal{Q}_{sel}^{(\tau)}-\lambda_2(\mathcal{Q}_{sel}^{(\tau)})^2\Bigr).
\end{align}

Эта формула, конечно, не окончательна. Но она уже делает нечто важное: соединяет в одном объекте штраф за разрушение идентичности, цену режимной перестройки и выигрыш от структурной согласованности.

\subsection{Как различаются три траектории в этой форме}

Теперь можно увидеть, как новые формы работают на ранее введённых модельных траекториях.

Для устойчивой траектории \(\gamma_{st}\) расстояние \(d_{inv}\) мало, скорость режимного дрейфа \(\|\dot{\Xi}\|\) мала, а \(\mathcal{Q}_{sel}\) находится в благоприятной области. Следовательно, \(\mathfrak{A}_{q}\) получает малое или локально минимальное значение.

Для метастабильной траектории \(\gamma_{meta}\) инвариантный штраф всё ещё остаётся умеренным, но растёт режимный вклад за счёт дрейфа между долинами, а селекционный выигрыш уже не максимален. Поэтому \(\mathfrak{A}_{q}\) остаётся локально приемлемым, но больше, чем у устойчивой траектории.

Для разрывной траектории \(\gamma_{break}\) резко возрастает либо инвариантный штраф, либо потенциал \(V_{mode}\), либо разрушается сама применимость селекционного выигрыша. В результате \(\mathfrak{A}_{q}\) теряет стационарный характер, и траектория выходит из области допустимого описания.

\subsection{Что эти формы пока не делают}

Принципиально важно отметить и ограничения введённой схемы. Здесь ещё не задана конкретная метрика на пространстве инвариантных ядер, не построена явная геометрия пространства \(\Xi\), не объяснено происхождение коэффициентов \(\alpha_{inv}\), \(\alpha_{m}\), \(\lambda_1\), \(\lambda_2\), и не показано, как именно \(\mathcal{Q}_{sel}\) вычисляется в нетривиальных примерах.

Но на этой стадии это скорее достоинство, чем недостаток. Трактат получает достаточно формальную модель, чтобы начать различать сценарии динамики, и при этом сохраняет свободу для последующей специализации и уточнения.

\subsection{Следующий шаг после пробных форм}

После введения этих полуколичественных форм дальнейший шаг становится уже почти экспериментом внутри самого трактата. Теперь можно либо построить первую фазовую диаграмму по параметрам \(\alpha_{inv}\), \(\alpha_{m}\), \(\lambda_1\), \(\lambda_2\), либо разобрать несколько конкретных эволюционных сценариев с указанием того, какой вклад начинает доминировать. Именно это и станет переходом от полуформальной динамики к первому внутреннему анализу режимов движения.

\section{Иллюстративные схемы текущего тома}

Чтобы завершающая часть трактата не оставалась чисто вербальной, соберём в одном месте несколько рабочих иллюстраций. Эти рисунки не взяты извне и не вставлены как произвольные декоративные схемы: они заново сгенерированы локальным скриптом \texttt{generate\_shadow\_figures.py} на основе вычислительных результатов текущего корпуса файлов. Поэтому они не подменяют доказательства и не фиксируют окончательную формализацию, но помогают увидеть, как разные линии исследования уже складываются в общую конфигурацию. В этом смысле рисунки ниже играют роль не украшения, а компактных карт: каждая из них подчёркивает один из тех переходов, которые весь том пытался сделать явными --- от внутренней секторности к наблюдаемым, от фазовой структуры к допустимым ветвям, от окон устойчивости к режимной селекции и от локальных данных к общей тепловой карте согласованности.

Дополнительно в завершающую часть тома включим и прямую фазовую карту, построенную по вычислительной схеме перевода сырых чувствительностей в инварианты \(\log_{10}(\Xi)\) и \(\Upsilon\). В отличие от предыдущих рисунков, которые суммировали уже проведённые аудиты, эта диаграмма даёт более синтетический портрет: она размещает различные физические модели на одной плоскости режимов и делает наглядным само разделение между каркасной, фазовой, внутренней и смешанной областями.

\begin{figure}[htbp]
    \centering
    \includegraphics[width=0.97\textwidth]{fig_unified_phase_portrait.png}
    \caption{Фазовый портрет структурных режимов: размещение модельных физических теорий на плоскости инвариантов \(\log_{10}(\Xi)\) и \(\Upsilon\).}
    \label{fig:unified-phase-portrait}
\end{figure}

Эта карта особенно полезна тем, что сразу показывает два независимых направления различения. Горизонтальная координата измеряет степень ухода от чисто геометрического каркаса, а вертикальная --- баланс между фазовым и внутренним вкладами. Поэтому рисунок не просто иллюстрирует уже введённые понятия, а работает как компактный классификатор: гравитационные, топологические, спектрально-внутренние и смешанные режимы оказываются размещёнными в одном общем пространстве чтения.

\begin{figure}[htbp]
    \centering
    \includegraphics[width=0.9\textwidth]{fig_sector_attribution.png}
    \caption{Схема секторной атрибуции: каркасный, фазовый и внутренний слои читаются как различные проекции одной структурной основы.}
    \label{fig:sector-attribution}
\end{figure}

Первый рисунок особенно важен для онтологической части трактата. Он делает наглядным тезис о том, что различие между несущим каркасом, фазовой организацией и внутренней нагрузкой не обязано означать различие независимых сущностей. Наоборот, эти слои можно читать как согласованные режимы разложения одного и того же минимального ядра. Тем самым рисунок поддерживает центральную линию текста: глубинная сущность не обязана быть одномерной по способу своего проявления.

\begin{figure}[htbp]
    \centering
    \includegraphics[width=0.92\textwidth]{fig_gauge_holonomy_phases.png}
    \caption{Иллюстрация связи калибровочных фаз, голономий и допустимых ветвей глобальной согласованности.}
    \label{fig:gauge-holonomy}
\end{figure}

Второй рисунок связан уже не столько с внутренней секторностью, сколько с тем, как локальные калибровочные и фазовые данные вписываются в глобальную структуру. Именно здесь наиболее отчётливо видно, почему трактат постоянно возвращается к голономии, фазовому обходу и нестягиваемым циклам: эти элементы не являются побочным математическим декором, а помогают различать допустимые и недопустимые ветви эволюции.

\begin{figure}[htbp]
    \centering
    \includegraphics[width=0.92\textwidth]{fig_scalar_window_stability.png}
    \caption{Окна устойчивости и переходные области: полуколичественная схема, в которой различаются устойчивые, метастабильные и пограничные режимы.}
    \label{fig:scalar-window}
\end{figure}

Третий рисунок естественно читать вместе с разделами о статусной карте и квазидействии. Если предыдущие главы вводили различие между устойчивой, метастабильной и разрывной траекториями на языке понятий и функционалов, то здесь это различие получает компактную визуальную форму. Благодаря этому становится яснее, что квазидействие --- не абстрактный риторический жест, а средство отбора между реально различимыми областями поведения.

\begin{figure}[htbp]
    \centering
    \includegraphics[width=0.94\textwidth]{fig_enriched_heatmap.png}
    \caption{Обогащённая тепловая карта согласованности: сводное представление локальных вкладов, режимных различий и общей селекционной структуры.}
    \label{fig:enriched-heatmap}
\end{figure}

Наконец, четвёртый рисунок полезен как сжатый образ всей исследовательской программы. Он показывает, что даже при неполноте формализма уже возможно собирать разнородные признаки --- инвариантные, спектральные, фазовые и корреляционные --- в одну карту, где становятся различимыми области высокой, промежуточной и слабой согласованности. Именно такая карта и делает переход от философской гипотезы к рабочему протоколу сравнения моделей практически осуществимым.

\section{Заключение: статус первого тома}

Теперь можно сформулировать, в чём именно состоит результат этого тома. Он не дал окончательной теории мира и не вывел замкнутую систему уравнений, из которых автоматически следовали бы все наблюдаемые структуры. Но он сделал нечто менее окончательное и в то же время принципиально необходимое: он собрал разрозненные геометрические, спектральные, топологические, корреляционные и динамические мотивы в единое исследовательское пространство, где их совпадения перестают выглядеть случайными.

Главное достижение текущего текста состоит в переводе интуиции в аппарат различений. На первых страницах трактата речь шла о предчувствии единой сущности, лежащей глубже привычных делений между геометрией, полем, спектром и наблюдаемым. К концу тома это предчувствие уже получило несколько более строгих форм: язык минимального ядра, различение инвариантного, режимного и селекционного вкладов, статусную карту элементов аппарата, а также первое квазидействие, позволяющее говорить о допустимой динамике не только метафорически, но и структурно.

Тем самым том выполняет три функции сразу. Во-первых, он фиксирует онтологическую ставку: первичен не набор готовых объектов, а закон согласованного становления допустимых структур. Во-вторых, он предлагает рабочий словарь, в котором можно сравнивать конкурирующие описания одного и того же глубинного содержания. В-третьих, он задаёт исследовательский мост от философской формулировки к полуформальному механизму отбора, где уже можно различать устойчивые, метастабильные и разрушительные траектории.

Из этого следует и более трезвый вывод. Настоящий том не завершает программу, а только делает её внятной. Он не устраняет нерешённость, но локализует её. Теперь уже ясно, где именно остаётся пробел: в явном определении пространства инвариантных ядер, в геометрии режимного слоя, в вычислении селекционного функционала, в согласовании полуколичественных коэффициентов с содержательными классами наблюдаемых и в доказательстве того, что единый источник действительно способен порождать всю заявленную иерархию редукций.

Однако именно такая незавершённость и является честным признаком продвижения. До этого текста разрыв между общей интуицией и рабочим формализмом был слишком велик. После него разрыв не исчез, но стал измеримым. Мы уже можем назвать основные объекты, выделить их статусы, проследить линии редукции, задать первые функционалы отбора и наметить сценарии, в которых гипотеза о глубинной ткани согласованности либо усилится, либо будет опровергнута.

Поэтому завершение первого тома должно быть сформулировано не как окончательный триумф, а как точка дисциплинированного перехода. Исследование больше не находится в состоянии чистого рассеяния. Оно получило каркас. Если исходная интуиция верна, следующие тома должны будут превратить этот каркас в более строгую теорию реконструкции: вывести явные модельные пространства, уточнить динамику минимального ядра, связать селекционные показатели с наблюдаемыми классами величин и построить проверяемые сценарии совместимости с известной физикой.

Итак, первый том можно считать завершённым в собственном смысле: не как замкнутую теорию, а как основание программы. Он формулирует предмет, очерчивает язык, вводит первичные различения, предлагает начальный вариационный механизм и показывает, что разнородные результаты действительно допускают чтение как следы одной глубинной структуры. Если последующая работа окажется успешной, именно этот том будет читаться как место, где впервые стало видно не только то, что следует объяснить, но и в каком именно пространстве понятий такое объяснение вообще может быть найдено.

\section{Послесловие после проверки вторым томом}

Настоящее послесловие написано после завершения основной серии вычислительных и операторных аудитов второго тома. Оно не отменяет исходную постановку первого тома, но уточняет её научный статус. Первый том формулировал гипотезу о едином структурном источнике и предлагал язык, на котором такую гипотезу можно сделать проверяемой. Второй том впервые применил к этой гипотезе достаточно жёсткие критерии: единство меры, отсутствие скрытых коэффициентов, согласованность нормировок, воспроизводимость спектральных расчётов и обязательное прохождение независимых секторных тестов.

Результат этой проверки оказался смешанным. Программа не исчезла, но разделилась на три существенно различных слоя.

\subsection{Что сохранилось как математическое ядро}

Наиболее устойчивой частью программы остаётся не совокупность численных формул, а геометрический и операторный каркас. Конкретное содержание имеют:
\begin{itemize}
    \item носитель \(\mathbb{RP}^3\times S^1\), его накрытие, объёмная нормировка и нетривиальная \(\mathbb Z_2\)-голономия;
    \item различение локальных мод, глобальных фаз, spin-структур и quotient-секторов;
    \item спектральные башни и внешне воспроизводимые свойства scalar- и coexact-операторов на \(\mathbb{RP}^3\);
    \item defect-конструкции, в которых топология действительно ограничивает ранг ядра и допустимые ветви;
    \item сформулированные no-go критерии, запрещающие выдавать подгонку нормировки за геометрический вывод.
\end{itemize}

Эти результаты не образуют завершённой физической теории, однако они уже являются содержательной математической программой. Особенно важно, что отрицательные результаты здесь не уничтожают каркас, а показывают границы конкретных механизмов.

\subsection{Что сохраняется только как условная связь}

Несколько численно сильных соотношений выдержали проверки на воспроизводимость и простоту, но не получили единого операторного происхождения. К ним относятся электромагнитная вакуумная комбинация, отношение масс \(\tau\)-лептона и мюона, а также абсолютный хиггсовский масштаб, построенный из этих величин.

Их корректный статус теперь таков:
\begin{equation}
\text{численно сильная структурная связь}
\quad\neq\quad
\text{вывод из единого действия}.
\end{equation}
Для \(\tau\)-формулы остаются невыведенными как исходный объёмный seed, так и нормировка компактной петли. Хиггсовский абсолютный масштаб наследует эту условность. В электромагнитном секторе точное малое спектральное дополнение не было восстановлено из полной Maxwell--ghost схемы в одной фиксированной нормировке.

Поэтому эти формулы допустимо сохранять как мотивы и кандидаты на будущую реконструкцию, но недопустимо использовать как доказательство уже найденной единой сущности.

\subsection{Какие ветви закрыты отрицательно}

Второй том закрыл несколько естественных минимальных маршрутов. Обычный промежуточный RG-бег дополнительного split-сектора имеет неверное направление поправок. Минимальный логарифмический threshold-cone не достигает требуемого электрослабого и QCD-сдвига в физическом диапазоне масс. Каноническая charged-lepton нормировка не порождает требуемые объёмные коэффициенты. Наконец, единое минимальное trace--Hodge действие сохраняет нейтринную collective-норму \(23+\pi^{-1}\), но не воспроизводит независимо второй normalization-sensitive сектор.

Последний результат является центральным. Он означает не то, что идея общего источника логически опровергнута, а то, что текущий набор секторных формул ещё нельзя считать редукциями одного найденного действия. Между общей онтологической гипотезой и физической теорией остаётся именно тот разрыв, который первый том стремился сделать измеримым.

\subsection{Новая статусная формула программы}

После проведённых проверок программу следует записывать не как цепочку уже закрытых физических выводов, а как трёхуровневую структуру:
\begin{equation}
\boxed{
\begin{gathered}
\text{геометрическое ядро}
\;\longrightarrow\;
\text{условные секторные модели}
\\[2mm]
\dashrightarrow\;
\text{непостроенное единое действие}
\end{gathered}
}
\end{equation}
Сплошная стрелка обозначает реально построенные специализации и операторы. Пунктирная стрелка обозначает главный открытый переход: необходимость вывести общую меру и относительные секторные веса до обращения к наблюдаемым числам.

Такое уточнение не является поражением исходного замысла. Напротив, оно выполняет методологическое обещание первого тома: гипотеза о структуре согласованности должна быть способна не только порождать красивые аналогии, но и запрещать собственные неверные реализации.

Финальный falsification-аудит второго тома делает статус ещё точнее. Из двенадцати остаточных групп утверждений пять относятся к наблюдаемым величинам, но ни одна не сохраняет одновременно заранее фиксированные оператор и меру, prospective blind-статус и происхождение из единого двухсекторного действия. Поэтому версия II.A отвергается как замкнутая единая предсказательная физическая теория. Это не обнуляет математическое ядро первого тома, но прекращает использование его численных мотивов как физического доказательства.

\subsection{Новый запрет и условие переоткрытия}

После второго тома запрещено продолжать программу путём добавления новых численных совпадений, новых секторных коэффициентов или специальных нормировок, выбранных после сравнения с экспериментом. Следующий механизм считается допустимым только при одновременном выполнении трёх условий:
\begin{enumerate}
    \item мера, stiffness-оператор или boundary functional выводятся из предварительно объявленной симметрии, топологии либо квантовой меры;
    \item все относительные веса фиксируются до сравнения с массами и константами связи;
    \item один и тот же механизм воспроизводит не менее двух независимых normalization-sensitive секторов.
\end{enumerate}

Если такой механизм не будет найден, программа сохранит ценность как геометрическая и философско-математическая гипотеза, но не должна называться единой предсказательной физической теорией. Если же он будет найден, именно этот двухсекторный тест станет первым реальным мостом от языка первого тома к физическому замыканию.

\section{Перспектива следующего тома}

Следующий том должен начаться не с новой феноменологической формулы, а с ревизии самого перехода от геометрического ядра к физическому действию. Его первейшей задачей станет классификация тех мер и stiffness-операторов на пространстве \(\mathfrak{K}_{min}\), которые могут возникать из симметрии, граничных данных, quotient-конструкции или квантового якобиана без введения секторных коэффициентов вручную.

Затем потребуется построить строгую геометрию пространства режимов и деформаций, но уже с учётом нормировочного no-go. Метрика на пространстве допустимых конфигураций должна не просто измерять расстояния между режимами, а единственным образом фиксировать относительные веса нуль-форм, связностей, defect-мод и фермионных операторов. В противном случае квазидействие останется полезным классификатором, но не станет физическим действием.

Решающим результатом следующего тома должен стать не очередной близкий численный ответ, а один из двух честных исходов. Положительный исход состоит в выводе единого функционала, который без подгонки проходит как минимум два независимых секторных теста. Отрицательный исход состоит в доказательстве того, что выбранный класс геометрических мер принципиально не способен связать эти сектора. Оба результата будут продвижением: первый откроет путь к физической теории, второй точно очертит область, в которой исходная онтологическая гипотеза остаётся только исследовательской метафизикой.

Таким образом, перспектива программы становится уже и строже. Теперь требуется не расширение числа мотивов, а сокращение пространства допустимых механизмов. Только после прохождения этого этапа гипотеза о структурно-спектральной ткани согласованности сможет вновь претендовать на переход от философской глубины к теоретической ответственности.

\section{Финальная заморозка статуса II.A/II.B}

После окончательного аудита нейтринного selector-а и поколенчески-ковариантной
ветви количественные нейтринные формулы сняты со статуса предсказаний.
Поколенчески-слепой defect сохраняется только как partial topology result с
rank-three kernel и нечётным mod-two индексом. Ни \(23+\pi^{-1}\), ни
\(21+3\pi^{-1}\) не являются выведенным физическим знаменателем.

До полного внешнего computational reproduction и mathematical replication
фиксируется
\begin{equation}
 R_{\rm sci}=4/10,
 \qquad
 N_{\rm closed\ physical}=0.
\end{equation}
Прежняя внутренняя оценка \(5/10\) сохраняется только как исторический
milestone повторного вычисления. Окончательная формула программы имеет вид
\begin{equation}
 \boxed{
 \begin{gathered}
 \text{математическое ядро предъявлено к независимой проверке},\\
 \text{единая физическая теория не построена},\\
 \text{новый parent action относится уже к отдельной версии III}.
 \end{gathered}}
\end{equation}

\end{document}